% Les outils numériques pédagogiques et l'intelligence artificielle — Allemand Tle A — Cours signé OnBuch+
% Programme officiel MINESEC. Compiler avec : tectonic AL19-les-outils-numeriques-pedagogiques-et-l-intelligence-artific.tex
\newcommand{\DOCMATIERE}{Allemand}
\newcommand{\DOCNIVEAU}{Tle A}
\newcommand{\DOCMODULE}{Chapitre 4 : Média et communication}
\newcommand{\DOCLECON}{AL19}
\newcommand{\DOCTITRE}{Les outils numériques pédagogiques et l'intelligence artificielle}
% OnBuch+ — préambule « cours signés » ENRICHI (compilé avec tectonic / XeLaTeX)
% Système visuel « Pop » d'OnBuch+ : fond crème (jamais blanc), bordures encre
% épaisses, ombres « dures » décalées, titres Archivo Black en capitales.
% Version MATHÉMATIQUES : graphes (pgfplots/tikz), boîtes pédagogiques
% (définition, propriété, méthode, attention, le savais-tu, exercice résolu,
% exercice, TP, bilan), page de garde et signature OnBuch+ ancrée en bas.
%
% Macros attendues dans main.tex AVANT \input{preamble} :
%   \DOCMATIERE  \DOCNIVEAU  \DOCMODULE  \DOCLECON (numéro)  \DOCTITRE
\documentclass[11pt]{article}

\usepackage{fontspec}
\usepackage[ngerman,french]{babel} % français = langue principale ; l'allemand via \all{..} / textebox
\usepackage[a4paper,margin=2cm,top=3.4cm,bottom=2.8cm,headheight=1.4cm,footskip=1.3cm]{geometry}
\usepackage{amsmath,amssymb,amsfonts}
\usepackage{mathtools} % \xmapsto, \coloneqq... souvent utilisés par les modèles
\usepackage{cancel} % simplifications barrées (bilans de forces)
% \abs{z} (module, valeur absolue) et \norm{u} (norme), utilisés par les modèles
\providecommand{\abs}{}\let\abs\relax\DeclarePairedDelimiter{\abs}{\lvert}{\rvert}
\providecommand{\norm}{}\let\norm\relax\DeclarePairedDelimiter{\norm}{\lVert}{\rVert}
\usepackage[table]{xcolor} % [table] : fournit aussi \rowcolors (alternance de lignes)
\usepackage{pagecolor}
\usepackage{tikz}
\usetikzlibrary{arrows.meta,positioning,calc,shapes.geometric,shapes.misc,shapes.arrows,decorations.pathmorphing,decorations.pathreplacing,decorations.markings,patterns,shadows,mindmap,trees,fit,backgrounds,matrix,intersections,angles,quotes}
\usepackage{pgfplots}
\usepackage[version=4]{mhchem} % équations nucléaires \ce{^{235}_{92}U}
\usepackage[european,siunitx]{circuitikz} % schémas électriques
\pgfplotsset{compat=1.18}
\usepgfplotslibrary{fillbetween,groupplots} % aires entre courbes (calcul intégral)
\usepackage{enumitem}
\usepackage{fancyhdr}
% [most] charge toutes les bibliothèques tcolorbox (listings, minted, raster,
% documentation...) et gonfle inutilement la mémoire TeX sur un document long
% (~300 boîtes) : on ne charge que ce qui est réellement utilisé.
\usepackage[skins,breakable]{tcolorbox}
\usepackage{graphicx}
\usepackage{colortbl}
\usepackage{booktabs}
\usepackage{tabularx}
\usepackage{diagbox} % case d'en-tête diagonale des tableaux à double entrée
% Colonnes extensibles centrée / gauche / droite (C, L, R), souvent utilisées par les modèles
\newcolumntype{C}{>{\centering\arraybackslash}X}
\newcolumntype{L}{>{\raggedright\arraybackslash}X}
\newcolumntype{R}{>{\raggedleft\arraybackslash}X}
\usepackage{array}
\usepackage{multicol}
\usepackage{multirow}
\usepackage{siunitx}
% parse-numbers=false : accepte aussi \SI{2,5 \times 10^{-3}}{...}
\sisetup{locale=FR,output-decimal-marker={,},group-minimum-digits=4,parse-numbers=false}
\DeclareSIUnit\ohm{\ensuremath{\symup{\Omega}}} % U+2126 absent de la police
\DeclareSIUnit\division{div} % réglages d'oscilloscope (ms/div, V/div)
\DeclareSIUnit\tr{tr} % tours (vitesse de rotation, ex. \SI{1500}{\tr\per\minute})
\DeclareSIUnit\molar{M} % concentration molaire (ex. \SI{3}{\molar}, \SI{1}{\micro\molar})
\DeclareSIUnit\an{an} % année (le modèle confond parfois avec \year, un compteur TeX) — ex. \SI{5}{\kilo\meter\per\an}
\DeclareSIUnit\FCFA{FCFA} % franc CFA (ex. \SI{500}{\FCFA\per\kilo\gram})
\DeclareSIUnit\degreeBrix{\textdegree Brix} % teneur en sucre d'un sirop/jus (réfractométrie)
\DeclareSIUnit\month{mois} % le modèle utilise parfois \month, un compteur TeX, comme unité de durée
\usepackage{qrcode}
% \degree : commande fréquemment utilisée par les modèles pour un angle ou une
% température en dehors de \SI{}{} (non fournie par les paquets chargés).
\providecommand{\degree}{^{\circ}}
% \qed (fin de démonstration), utilisé par les modèles sans amsthm
\providecommand{\qed}{\hfill$\square$}
% \barre : double barre des valeurs interdites dans un tableau de variations (tkz-tab)
\providecommand{\barre}{\ensuremath{\|}}
% environnement proof (amsthm non chargé) utilisé par les modèles
\ifdefined\proof\else\newenvironment{proof}[1][Démonstration]{\par\noindent\textit{#1.}\ }{\qed\par}\fi
\usepackage{lastpage}
\usepackage{hyperref}
\hypersetup{hidelinks}

% ============================================================
% Palette « Pop » OnBuch+ — mêmes hex que la classe P (plus_ui.dart)
% ============================================================
\definecolor{poporange}{HTML}{F59321}
\definecolor{popdark}{HTML}{E07A0C}
\definecolor{popcream}{HTML}{FFF6E8}
\definecolor{popink}{HTML}{1C1714}
\definecolor{poppurple}{HTML}{7A5AE0}
\definecolor{popgreen}{HTML}{2E9E6A}
\definecolor{poppink}{HTML}{E0457B}
\definecolor{popblue}{HTML}{2F7DE1}
\definecolor{popgold}{HTML}{C98A12}
\definecolor{popmuted}{HTML}{8A7F76}
\definecolor{popline}{HTML}{EADFD2}
\definecolor{popgray}{HTML}{8A7F76}
\colorlet{popgrey}{popgray}
% Alias tolérants (le modèle nomme parfois les couleurs différemment)
\colorlet{popwhite}{white}
\colorlet{popblack}{popink}
\colorlet{popred}{poppink}
\colorlet{popyellow}{popgold}
% Teintes claires (fonds de tableaux/encadrés) — le modèle nomme parfois
% les couleurs avec un suffixe L (ex. popblueL) sans les avoir définies.
\colorlet{poporangeL}{poporange!20}
\colorlet{popdarkL}{popdark!20}
\colorlet{poppurpleL}{poppurple!20}
\colorlet{popgreenL}{popgreen!20}
\colorlet{poppinkL}{poppink!20}
\colorlet{popblueL}{popblue!20}
\colorlet{popgoldL}{popgold!20}
\colorlet{popredL}{popred!20}
\colorlet{popgrayL}{popgray!20}
% Teintes claires pour fonds de tableaux / zones
\colorlet{popblueL}{popblue!10!white}
\colorlet{poporangeL}{poporange!14!white}
\colorlet{popgreenL}{popgreen!12!white}
\colorlet{poppurpleL}{poppurple!12!white}
\colorlet{poppinkL}{poppink!12!white}
\colorlet{popgoldL}{popgold!14!white}

\pagecolor{popcream}

% ============================================================
% Typographie
% ============================================================
\defaultfontfeatures{Ligatures=TeX}
\setmainfont{PlusJakartaSans-Medium.ttf}[Path=../../../fonts/,BoldFont=PlusJakartaSans-Bold.ttf]
% Symboles, API et emojis absents de la police principale : repli sur DejaVu Sans (caractères actifs)
\newfontface\symfont{DejaVuSans.ttf}[Path=../../../fonts/]
\makeatletter
\newcommand\symactive[1]{\catcode#1=13 \begingroup\lccode`\~=#1 \lowercase{\endgroup\def~}{{\symfont\char#1}}}
\newcommand\symblank[1]{\catcode#1=13 \begingroup\lccode`\~=#1 \lowercase{\endgroup\def~}{}}
\newcommand\symhyph[1]{\catcode#1=13 \begingroup\lccode`\~=#1 \lowercase{\endgroup\def~}{\mbox{-}}}
\newcount\symc
\newcommand\symrange[2]{\symc=#1 \@whilenum\symc<#2 \do{\expandafter\symactive\expandafter{\the\symc}\advance\symc by 1 }}
\newcommand\symblankrange[2]{\symc=#1 \@whilenum\symc<#2 \do{\expandafter\symblank\expandafter{\the\symc}\advance\symc by 1 }}
\makeatother
\symrange{"0250}{"02FF}\symactive{"2713}\symactive{"2714}\symactive{"2717}\symactive{"2718}\symactive{"2610}\symactive{"2611}\symactive{"2612}
\symhyph{"2011}\symhyph{"2E1A}\symblankrange{"1F300}{"1FAFF}\symblank{"FE0F}
% Maths en sans-serif (Fira Math) avec chiffres et lettres droites de Plus
% Jakarta, pour que les formules se fondent dans le texte.
\usepackage[math-style=ISO,bold-style=ISO]{unicode-math}
\setmathfont{FiraMath-Regular.otf}[Path=../../../fonts/]
\setmathfont{PlusJakartaSans-Medium.ttf}[Path=../../../fonts/,range={up/num,up/{Latin,latin},\mathup}]
\usepackage{icomma} % virgule décimale sans espace en mode math
% Caractères mathématiques Unicode absents des polices de texte (Plus Jakarta,
% Archivo Black) : rendus par leur équivalent mathématique, en texte comme en
% formule — sinon ils s'affichent en carré vide (ex. « Arithmétique dans ℤ »).
\usepackage{newunicodechar}
\newunicodechar{ℝ}{\ensuremath{\mathbb{R}}}
\newunicodechar{ℤ}{\ensuremath{\mathbb{Z}}}
\newunicodechar{ℂ}{\ensuremath{\mathbb{C}}}
\newunicodechar{ℕ}{\ensuremath{\mathbb{N}}}
\newunicodechar{ℚ}{\ensuremath{\mathbb{Q}}}
\newunicodechar{𝔻}{\ensuremath{\mathbb{D}}}
\newunicodechar{α}{\ensuremath{\alpha}}
\newunicodechar{β}{\ensuremath{\beta}}
\newunicodechar{γ}{\ensuremath{\gamma}}
\newunicodechar{δ}{\ensuremath{\delta}}
\newunicodechar{ε}{\ensuremath{\varepsilon}}
\newunicodechar{θ}{\ensuremath{\theta}}
\newunicodechar{λ}{\ensuremath{\lambda}}
\newunicodechar{μ}{\ensuremath{\mu}}
\newunicodechar{σ}{\ensuremath{\sigma}}
\newunicodechar{τ}{\ensuremath{\tau}}
\newunicodechar{φ}{\ensuremath{\varphi}}
\newunicodechar{ψ}{\ensuremath{\psi}}
\newunicodechar{ω}{\ensuremath{\omega}}
\newunicodechar{Σ}{\ensuremath{\Sigma}}
% majuscules produites par \MakeUppercase dans les titres (α-aminés -> Α)
\newunicodechar{Α}{\ensuremath{\alpha}}
\newunicodechar{Β}{\ensuremath{\beta}}
\newunicodechar{Γ}{\ensuremath{\Gamma}}
\newunicodechar{Δ}{\ensuremath{\Delta}}
\newunicodechar{Ε}{\ensuremath{\varepsilon}}
\newunicodechar{Θ}{\ensuremath{\Theta}}
\newunicodechar{Λ}{\ensuremath{\Lambda}}
\newunicodechar{Μ}{\ensuremath{\mu}}
\newunicodechar{Τ}{\ensuremath{\tau}}
\newunicodechar{Φ}{\ensuremath{\Phi}}
\newunicodechar{Ψ}{\ensuremath{\Psi}}
\newunicodechar{Ω}{\ensuremath{\Omega}}
\newunicodechar{↦}{\ensuremath{\mapsto}}
\newunicodechar{∈}{\ensuremath{\in}}
\newunicodechar{∉}{\ensuremath{\notin}}
\newunicodechar{∧}{\ensuremath{\wedge}}
\newunicodechar{⇔}{\ensuremath{\Leftrightarrow}}
\newunicodechar{⟺}{\ensuremath{\Longleftrightarrow}}
\newunicodechar{⇒}{\ensuremath{\Rightarrow}}
\newunicodechar{⟹}{\ensuremath{\Longrightarrow}}
\newunicodechar{∘}{\ensuremath{\circ}}
\newunicodechar{∪}{\ensuremath{\cup}}
\newunicodechar{∩}{\ensuremath{\cap}}
\newunicodechar{≡}{\ensuremath{\equiv}}
\newunicodechar{⊂}{\ensuremath{\subset}}
\newunicodechar{⊥}{\ensuremath{\perp}}
\newunicodechar{∃}{\ensuremath{\exists}}
\newunicodechar{∀}{\ensuremath{\forall}}
\newunicodechar{∛}{\ensuremath{\sqrt[3]{\,}}}
\newunicodechar{∜}{\ensuremath{\sqrt[4]{\,}}}
\newunicodechar{⃗}{\ensuremath{{}^{\to}}}
\newunicodechar{ⁿ}{\ifmmode^{n}\else\textsuperscript{\ensuremath{n}}\fi}
\newunicodechar{ˣ}{\ifmmode^{x}\else\textsuperscript{\ensuremath{x}}\fi}
\newunicodechar{ᵇ}{\ifmmode^{b}\else\textsuperscript{\ensuremath{b}}\fi}
\newunicodechar{ʳ}{\ifmmode^{r}\else\textsuperscript{\ensuremath{r}}\fi}
\newunicodechar{ᵘ}{\ifmmode^{u}\else\textsuperscript{\ensuremath{u}}\fi}
\newunicodechar{ʸ}{\ifmmode^{y}\else\textsuperscript{\ensuremath{y}}\fi}
\newunicodechar{ᵃ}{\ifmmode^{a}\else\textsuperscript{\ensuremath{a}}\fi}
\newunicodechar{ᵉ}{\ifmmode^{e}\else\textsuperscript{\ensuremath{e}}\fi}
\newunicodechar{ᵏ}{\ifmmode^{k}\else\textsuperscript{\ensuremath{k}}\fi}
\newunicodechar{ᵖ}{\ifmmode^{p}\else\textsuperscript{\ensuremath{p}}\fi}
\newunicodechar{ᴺ}{\ifmmode^{N}\else\textsuperscript{\ensuremath{N}}\fi}
\newunicodechar{ᶜ}{\ifmmode^{c}\else\textsuperscript{\ensuremath{c}}\fi}
\newunicodechar{ʰ}{\ifmmode^{h}\else\textsuperscript{\ensuremath{h}}\fi}
\newunicodechar{ᵠ}{\ifmmode^{q}\else\textsuperscript{\ensuremath{q}}\fi}
\newunicodechar{ₙ}{\ifmmode_{n}\else\textsubscript{\ensuremath{n}}\fi}
\newunicodechar{ₐ}{\ifmmode_{a}\else\textsubscript{\ensuremath{a}}\fi}
\newunicodechar{ᵢ}{\ifmmode_{i}\else\textsubscript{\ensuremath{i}}\fi}
\newunicodechar{ᵣ}{\ifmmode_{r}\else\textsubscript{\ensuremath{r}}\fi}
\newunicodechar{ₖ}{\ifmmode_{k}\else\textsubscript{\ensuremath{k}}\fi}
\newunicodechar{ᵦ}{\ifmmode_{\beta}\else\textsubscript{\ensuremath{\beta}}\fi}
\newunicodechar{ₓ}{\ifmmode_{x}\else\textsubscript{\ensuremath{x}}\fi}
\newfontfamily\popdisplay{ArchivoBlack-Regular.ttf}[Path=../../../fonts/]
\newfontfamily\poplogo{SpaceGrotesk-Bold.ttf}[Path=../../../fonts/]
\newfontfamily\popsemi{PlusJakartaSans-SemiBold.ttf}[Path=../../../fonts/]
\newfontfamily\popxbold{PlusJakartaSans-ExtraBold.ttf}[Path=../../../fonts/]
\newfontfamily\popxboldls{PlusJakartaSans-ExtraBold.ttf}[Path=../../../fonts/,LetterSpace=3.0]

\setlist[enumerate]{leftmargin=1.7em,itemsep=5pt,topsep=4pt}
\setlist[itemize]{leftmargin=1.7em,itemsep=4pt,topsep=4pt,label={\color{poporange}\textbullet}}
\setlength{\parindent}{0pt}
\setlength{\parskip}{5pt}
\renewcommand{\arraystretch}{1.25}

% pgfplots : style Pop par défaut
\pgfplotsset{
  every axis/.append style={axis line style={popink,line width=1pt},tick style={popink},
    grid style={popline,line width=0.5pt},label style={font=\small},tick label style={font=\footnotesize}},
  popcurve/.style={poporange,line width=1.8pt,no markers,smooth},
}
% Même style disponible dans un \draw TikZ simple (hors pgfplots)
\tikzset{popcurve/.style={poporange,line width=1.8pt}}
\tikzset{popgrid/.style={popline,very thin}}
\colorlet{popcurve}{poporange} % le nom du style est parfois utilisé comme couleur

% ============================================================
% Pill / squiggle
% ============================================================
\newcommand{\poppill}[3]{%
  \tikz[baseline=-0.55ex]\node[draw=popink,line width=1.2pt,rounded corners=2.6pt,%
    fill=#2,inner xsep=6.5pt,inner ysep=3pt]{%
    \popxboldls\fontsize{7}{7}\selectfont\color{#3}\MakeUppercase{#1}};%
}
\newcommand{\popsquiggle}[1]{%
  \tikz[baseline=-0.4ex]{
    \draw[#1,line width=2.6pt,line cap=round]
      plot [smooth] coordinates {(0,0.09) (0.34,0.01) (0.68,0.21) (1.02,0.01) (1.36,0.10)};
  }%
}
% Mot-clé surligné (marqueur orange sous le texte)
\newcommand{\cle}[1]{{\popxbold\color{popdark}#1}}

% ============================================================
% En-tête / pied
% ============================================================
\pagestyle{fancy}
\fancyhf{}
\renewcommand{\headrulewidth}{0pt}
\renewcommand{\footrulewidth}{0pt}
\lhead{\raisebox{-0.15ex}{\poplogo\fontsize{13}{13}\selectfont\color{popink}OnBuch\color{poporange}+}}
\rhead{\poppill{\DOCMATIERE\ \textbullet\ \DOCNIVEAU\ \textbullet\ Leçon \DOCLECON}{white}{popink}}
\lfoot{\popxbold\fontsize{7.6}{7.6}\selectfont\color{popmuted}\MakeUppercase{Cours signé OnBuch+}\ \textbullet\ {\popsemi\DOCTITRE}}
\rfoot{\tikz[baseline=-0.6ex]\node[rounded corners=8.5pt,fill=popink,minimum height=17pt,minimum width=17pt,inner xsep=5pt,inner sep=0]{\popxbold\fontsize{8}{8}\selectfont\color{popcream}\thepage\,/\,\pageref*{LastPage}};}

% ============================================================
% Titres
% ============================================================
\newcommand{\chaptitre}[2]{%
  \begin{tcolorbox}[enhanced,colback=poporange,colframe=popink,boxrule=2.6pt,arc=11pt,
      drop shadow=popink,left=15pt,right=15pt,top=12pt,bottom=11pt,before skip=3pt,after skip=18pt]
    \poppill{#2}{popink}{popcream}\par\vspace{9pt}
    {\raggedright\hyphenpenalty=10000\exhyphenpenalty=10000\popdisplay\fontsize{22}{24}\selectfont\color{popcream}\MakeUppercase{#1}\par}
  \end{tcolorbox}
}
% Section numérotée : pastille orange avec le numéro + titre Archivo
\newcounter{popsec}
\newcommand{\coursec}[1]{%
  \stepcounter{popsec}\setcounter{popsub}{0}%
  \par\vspace{16pt}\noindent
  \begin{minipage}{\linewidth}
  \tikz[baseline=-0.7ex]\node[draw=popink,line width=1.6pt,fill=poporange,rounded corners=4pt,minimum size=22pt,inner sep=0]{\popdisplay\fontsize{12}{12}\selectfont\color{popcream}\thepopsec};%
  \hspace{8pt}{\popdisplay\fontsize{15}{17}\selectfont\color{popink}\MakeUppercase{#1}}\par
  \vspace{4pt}\noindent{\color{poporange}\rule{36pt}{3.6pt}}
  \end{minipage}\par\nopagebreak\vspace{8pt}
}
\newcounter{popsub}
\newcommand{\courssub}[1]{%
  \stepcounter{popsub}%
  \par\vspace{9pt}\noindent{\popxbold\fontsize{11.5}{13}\selectfont\color{popink}{\color{poporange}\thepopsec.\thepopsub}\hspace{6pt}#1}\par\nopagebreak\vspace{4pt}
}

% ============================================================
% Boîtes pédagogiques — même recette : fond blanc, bordure épaisse colorée,
% ombre dure encre, badge plein. Titre optionnel [#1].
% ============================================================
\tcbset{popbase/.style={breakable,enhanced,colback=white,boxrule=2.3pt,arc=9pt,
  shadow={3pt}{-3pt}{0pt}{popink},left=14pt,right=14pt,top=11pt,bottom=11pt,before skip=11pt,after skip=15pt,
  fontupper=\popsemi\fontsize{10.6}{14.5}\selectfont\color{popink},
  fontlower=\popsemi\fontsize{10.6}{14.5}\selectfont\color{popink}}}
% Coche dessinée (le glyphe ✓ n'existe pas dans les polices du document)
\newcommand{\popcheck}[1]{\tikz[baseline=-0.5ex]\draw[#1,line width=1.6pt,line cap=round,line join=round](-0.13,0.0)--(-0.04,-0.09)--(0.13,0.1);}
\providecommand{\coche}{\popcheck{popgreen}}
\providecommand{\resultat}[1]{\par\smallskip\noindent\fcolorbox{popgreen}{popgreenL}{\parbox{\dimexpr\linewidth-2\fboxsep-2\fboxrule\relax}{\textbf{Résultat.} #1}}\par\smallskip}
\newcommand{\popboxhead}[3]{\poppill{#1}{#2}{white}\ifx\relax#3\relax\else\hspace{6pt}{\popxbold\fontsize{10.6}{12}\selectfont\color{popink}#3}\fi\par\vspace{7pt}}

\newtcolorbox{aretenir}[1][]{popbase,colframe=popgreen,before upper={\popboxhead{À retenir}{popgreen}{#1}}}
% Texte en allemand (document de lecture, dialogue, extrait) : cadre
% d'étude. Un titre optionnel (source, auteur) s'affiche dans l'en-tête.
\newtcolorbox{textebox}[1][]{popbase,colframe=popblue,before upper={\popboxhead{Text}{popblue}{#1}\begin{otherlanguage}{ngerman}},after upper={\end{otherlanguage}}}
% Marques ✓ / ✗ (forme correcte / fautive) souvent utilisées par les modèles
\providecommand{\cross}{\textcolor{poppink}{\ensuremath{\times}}}
\providecommand{\xmark}{\cross}\providecommand{\crossmark}{\cross}\providecommand{\cmark}{\ensuremath{\checkmark}}
% Ligne de réponse / justification à compléter (inventée par les modèles)
\providecommand{\justification}[1]{\par\noindent\textit{Justification~:} \dotfill\par}
% \textipa (package tipa non chargé) : la police ne couvre pas l'API -> simple passage
\providecommand{\textipa}[1]{#1}
% Soulignement (ulem non chargé) : \uline, \ul
\providecommand{\uline}[1]{\underline{#1}}\providecommand{\ul}[1]{\underline{#1}}
% \glossaire{\item ...} inventée par les modèles : liste de glossaire
\providecommand{\glossaire}[1]{\begin{itemize}#1\end{itemize}}
% \alert (beamer) inventée par les modèles : texte mis en évidence
\providecommand{\alert}[1]{\textbf{\textcolor{poppink}{#1}}}
% variantes ASCII d'ordinaux écrites en macro
\providecommand{\iere}{\textsuperscript{re}}\providecommand{\ieme}{\textsuperscript{e}}\providecommand{\ere}{\textsuperscript{re}}\providecommand{\eme}{\textsuperscript{e}}\providecommand{\er}{\textsuperscript{er}}\providecommand{\ie}{\textsuperscript{e}}
% Ordinaux français écrits en macro par les modèles : 1\ière, 3\ième
\newcommand{\ière}{\textsuperscript{re}}\newcommand{\ième}{\textsuperscript{e}}
% \exemple (inventée par les modèles) : étiquette « Exemple : »
\providecommand{\exemple}{\textbf{Exemple~:}\ }
% \square utilisable aussi hors mode math (cases à cocher)
\AtBeginDocument{\let\squareMath\square\renewcommand{\square}{\ensuremath{\squareMath}}}
% Mot ou phrase en allemand à l'intérieur d'un paragraphe français
\newcommand{\all}[1]{\ifmmode\text{\textit{#1}}\else\begingroup\selectlanguage{ngerman}\itshape #1\endgroup\fi}
\let\esp\all % alias toléré
\newtcolorbox{exemplebox}[1][]{popbase,colframe=poporange,before upper={\popboxhead{Exemple}{poporange}{#1}}}
\newtcolorbox{definition}[1][]{popbase,colframe=popblue,before upper={\popboxhead{Définition}{popblue}{#1}}}
\newtcolorbox{propriete}[1][]{popbase,colframe=poppurple,before upper={\popboxhead{Propriété}{poppurple}{#1}}}
\newtcolorbox{methode}[1][]{popbase,colframe=popgold,before upper={\popboxhead{Méthode}{popgold}{#1}}}
\newtcolorbox{attention}[1][]{popbase,colframe=poppink,before upper={\popboxhead{Attention}{poppink}{#1}}}
\newtcolorbox{experience}[1][]{popbase,colframe=popblue,colback=popblueL,before upper={\popboxhead{Expérience}{popblue}{#1}}}
% « Le savais-tu ? » : carte violette pleine (contexte, culture, Cameroun)
\newtcolorbox{savaistu}[1][]{popbase,colback=poppurple,colframe=popink,
  fontupper=\popsemi\fontsize{10.4}{14.2}\selectfont\color{white},
  before upper={\poppill{Le savais-tu\,?}{white}{poppurple}\ifx\relax#1\relax\else\hspace{6pt}{\popxbold\color{white}#1}\fi\par\vspace{7pt}}}
% Exercice résolu : énoncé en haut, \tcblower puis la solution sur fond crème
\newcounter{popexo}\newcounter{popexr}
\newtcolorbox{exoresolu}[1][]{popbase,colframe=poporange,colbacklower=popcream,
  segmentation style={popink,line width=1.2pt,dashed},
  before upper={\stepcounter{popexr}\popboxhead{Exercice résolu \thepopexr}{poporange}{#1}},
  before lower={\poppill{Solution}{popink}{popcream}\par\vspace{6pt}}}
% Exercice (non résolu en place) — la correction est dans la partie Corrigés
\newtcolorbox{exercice}[1][]{popbase,colframe=popink,
  before upper={\stepcounter{popexo}\popboxhead{Exercice \thepopexo}{popink}{#1}}}
% Corrigé
\newtcolorbox{corrige}[1][]{popbase,colframe=popgreen,colback=popgreenL,
  before upper={\popboxhead{Corrigé}{popgreen}{#1}}}
% Objectifs / prérequis / bilan
\newtcolorbox{objectifs}{popbase,colframe=popink,colback=poporangeL,
  before upper={\popboxhead{Objectifs de la leçon}{popink}{}}}
\newtcolorbox{prerequis}{popbase,colframe=popink,colback=popblueL,
  before upper={\popboxhead{Prérequis}{popblue}{}}}
\newtcolorbox{bilan}{popbase,colframe=popink,colback=white,boxrule=2.8pt,
  before upper={\popboxhead{Fiche bilan}{poporange}{L'essentiel à connaître pour l'examen}}}
% Figure encadrée avec légende
\newtcolorbox{popfigure}[1][]{enhanced,breakable=false,colback=white,colframe=popline,boxrule=1.4pt,arc=8pt,
  left=10pt,right=10pt,top=10pt,bottom=8pt,before skip=10pt,after skip=12pt,halign=center,
  lower separated=false,halign lower=center,
  fontlower=\popsemi\fontsize{9}{11.5}\selectfont\color{popmuted},
  #1}
\newcommand{\legende}[1]{\par\vspace{6pt}{\popsemi\fontsize{9}{11.5}\selectfont\color{popmuted}#1}}

% ============================================================
% Signature OnBuch+
% ============================================================
\newcommand{\poplogoinline}[1]{{\poplogo\fontsize{#1}{#1}\selectfont\color{popink}OnBuch\color{poporange}+}}
\newcommand{\signatureonbuch}{%
  \begin{tcolorbox}[enhanced,colback=popink,colframe=popink,boxrule=2.2pt,arc=10pt,
      drop shadow=poporange,left=14pt,right=14pt,top=10pt,bottom=10pt,before skip=0pt,after skip=0pt]
    \tikz[baseline=-0.7ex]\node[circle,fill=popgreen,draw=popcream,line width=1.3pt,minimum size=22pt,inner sep=0]{\popcheck{white}};%
    \hspace{9pt}\begin{minipage}[c]{0.62\linewidth}
      {\popxbold\fontsize{12}{13}\selectfont\color{popcream}Signé\ \textbullet\ {\poplogo OnBuch\color{poporange}+}}\par\vspace{2pt}
      {\raggedright\popsemi\fontsize{8}{10.5}\selectfont\color{popline}\MakeUppercase{Équipe pédagogique OnBuch+}\\\MakeUppercase{Conforme au programme officiel MINESEC}\par}
    \end{minipage}\hfill
    {\poplogo\fontsize{22}{22}\selectfont\color{popcream}OnBuch\color{poporange}+}
  \end{tcolorbox}%
}

% ============================================================
% Page de garde
% #1 titre · #2 matière · #3 niveau · #4 module · #5 n° leçon · #6 durée ·
% #7 description / objectifs (texte ou itemize)
% ============================================================
\newcommand{\pagegarde}[7]{%
  \thispagestyle{empty}
  \newgeometry{a4paper,margin=1.7cm,top=1.6cm,bottom=1.4cm}%
  \begin{tikzpicture}[remember picture,overlay]
    \fill[poporange!18] ([xshift=-3.2cm,yshift=-3.6cm]current page.north east) circle (4.6cm);
    \fill[poppurple!14] ([xshift=2.4cm,yshift=7.6cm]current page.south west) circle (3.8cm);
  \end{tikzpicture}%
  {\poplogo\fontsize{30}{30}\selectfont\color{popink}OnBuch\color{poporange}+}\hfill
  \raisebox{0.6ex}{\poppill{Cours signé \textbullet\ Premium}{popink}{popcream}}\par
  \vspace{1pt}\popsquiggle{poppurple}\par
  \vspace{8pt}
  \begin{tcolorbox}[enhanced,colback=poporange,colframe=popink,boxrule=2.8pt,arc=13pt,
      drop shadow=popink,left=18pt,right=18pt,top=12pt,bottom=13pt,after skip=10pt]
    \poppill{#2 \textbullet\ #3}{popink}{popcream}\hspace{5pt}\poppill{#4}{white}{popink}\par\vspace{8pt}
    {\popdisplay\fontsize{14}{14}\selectfont\color{popink}LEÇON #5}\par\vspace{4pt}
    {\raggedright\hyphenpenalty=10000\exhyphenpenalty=10000\popdisplay\fontsize{25}{27}\selectfont\color{popcream}\MakeUppercase{#1}\par}
    \vspace{8pt}
    {\popxbold\fontsize{9.5}{9.5}\selectfont\color{popink}\MakeUppercase{Durée conseillée : #6}}
  \end{tcolorbox}
  \begin{tcolorbox}[enhanced,colback=white,colframe=popink,boxrule=2.2pt,arc=10pt,
      drop shadow=popink,left=16pt,right=16pt,top=10pt,bottom=10pt,after skip=8pt]
    \poppill{Dans cette leçon}{poppurple}{white}\par\vspace{5pt}
    \popsemi\fontsize{9.3}{12.6}\selectfont\color{popink}#7
  \end{tcolorbox}
  \noindent\begin{minipage}[t]{0.487\linewidth}
    \begin{tcolorbox}[enhanced,colback=popgreen,colframe=popink,boxrule=2.2pt,arc=10pt,
        drop shadow=popink,left=11pt,right=11pt,top=10pt,bottom=10pt,height=3.1cm,valign=center]
      \tcbox[colback=white,colframe=popink,boxrule=1.3pt,arc=5pt,boxsep=0pt,nobeforeafter,
        left=3pt,right=3pt,top=3pt,bottom=3pt,valign=center]{\qrcode[height=1.9cm]{https://play.google.com/store/apps/details?id=com.luvvix.onbuch&hl=fr}}%
      \hspace{8pt}\begin{minipage}[c]{0.5\linewidth}
        {\popdisplay\fontsize{11}{12.5}\selectfont\color{white}\MakeUppercase{Télécharge OnBuch}}\par\vspace{4pt}
        {\raggedright\popsemi\fontsize{8.4}{11}\selectfont\color{white}Gratuit \textbullet\ Google Play\\Tuteur IA, annales, concours, résultats\par}
      \end{minipage}
    \end{tcolorbox}
  \end{minipage}\hfill
  \begin{minipage}[t]{0.487\linewidth}
    \begin{tcolorbox}[enhanced,colback=poppurple,colframe=popink,boxrule=2.2pt,arc=10pt,
        drop shadow=popink,left=11pt,right=11pt,top=10pt,bottom=10pt,height=3.1cm,valign=center]
      \tikz[baseline=-0.7ex]\node[circle,fill=white,draw=popink,line width=1.3pt,minimum size=1.6cm,inner sep=0]{\popdisplay\fontsize{24}{24}\selectfont\color{poppurple}+};%
      \hspace{8pt}\begin{minipage}[c]{0.6\linewidth}
        {\popdisplay\fontsize{11}{12.5}\selectfont\color{white}\MakeUppercase{Passe à OnBuch+}}\par\vspace{4pt}
        {\raggedright\popsemi\fontsize{8.4}{11}\selectfont\color{white}Tous les cours signés, Tuteur IA illimité, fiches PDF — onglet OnBuch+ de l'app\par}
      \end{minipage}
    \end{tcolorbox}
  \end{minipage}
  \vspace{8pt}
  \popcontenu
  \vfill
  \signatureonbuch
  \restoregeometry
  \newpage
}

% Rangée « Ce cours contient » (page de garde)
\newcommand{\poptile}[3]{% #1 couleur #2 gros texte #3 légende
  \begin{tcolorbox}[enhanced,colback=#1,colframe=popink,boxrule=2pt,arc=9pt,shadow={2.5pt}{-2.5pt}{0pt}{popink},
      left=6pt,right=6pt,top=9pt,bottom=9pt,halign=center,valign=center,height=1.75cm,width=\linewidth,nobeforeafter]
    {\popdisplay\fontsize{12.5}{13}\selectfont\color{white}\MakeUppercase{#2}}\par\vspace{3pt}
    {\centering\popsemi\fontsize{7.8}{9.5}\selectfont\color{white}#3\par}
  \end{tcolorbox}}
\newcommand{\popcontenu}{%
  \par\noindent{\popxboldls\fontsize{7.5}{7.5}\selectfont\color{popmuted}CE COURS SIGNÉ CONTIENT}\par\vspace{7pt}
  \noindent\begin{minipage}[t]{0.235\linewidth}\poptile{popblue}{Cours}{complet, illustré, pas à pas}\end{minipage}\hfill
  \begin{minipage}[t]{0.235\linewidth}\poptile{poporange}{Méthodes}{et exercices résolus}\end{minipage}\hfill
  \begin{minipage}[t]{0.235\linewidth}\poptile{poppink}{Exercices}{type examen + corrigés}\end{minipage}\hfill
  \begin{minipage}[t]{0.235\linewidth}\poptile{popgreen}{Fiche bilan}{l'essentiel à retenir}\end{minipage}\par}

% Sommaire
\newtcolorbox{sommaire}{enhanced,colback=white,colframe=popink,boxrule=2.2pt,arc=10pt,
  drop shadow=popink,left=18pt,right=18pt,top=13pt,bottom=13pt,before skip=6pt,after skip=20pt,
  before upper={\poppill{Sommaire}{poppurple}{white}\par\vspace{9pt}\popsemi\fontsize{10.6}{14.5}\selectfont\color{popink}}}

% Signature de fin de cours — poussée tout en bas de la dernière page
\newcommand{\findecours}{%
  \par\vfill\nopagebreak
  \begin{center}\popsquiggle{poporange}\end{center}\vspace{4pt}
  \signatureonbuch
}
\AtBeginDocument{\def\llbracket{\mathopen{[\mkern-3.5mu[}}\def\rrbracket{\mathclose{]\mkern-3.5mu]}}} % intervalles d'entiers (glyphe ⟦ absent de la police)
% Glyphes absents de Fira Math : redéfinis à partir de glyphes disponibles
\AtBeginDocument{%
  \def\setminus{\mathbin{\backslash}}\let\smallsetminus\setminus
  \def\vdots{\mathord{\vbox{\baselineskip3.2pt\lineskiplimit0pt\kern4pt\hbox{.}\hbox{.}\hbox{.}}}}%
  \def\ddots{\mathinner{\mkern1mu\raise6.5pt\hbox{.}\mkern2mu\raise3.6pt\hbox{.}\mkern2mu\raise0.7pt\hbox{.}\mkern1mu}}%
  \def\triangle{\mathord{\scalebox{0.9}{$\symup{\Delta}$}}}%
  \def\nabla{\mathord{\raisebox{\depth}{\scalebox{1}[-1]{$\symup{\Delta}$}}}}%
  \def\checkmark{\popcheck{popdark}}%
  \def\star{\mathbin{\ast}}\def\bigstar{\mathord{\ast}}%
  \def\diamond{\mathbin{\scalebox{0.6}{$\Diamond$}}}%
  \def\bigcup{\mathop{\mathchoice{\raisebox{-0.3em}{\scalebox{1.7}{$\cup$}}}{\scalebox{1.25}{$\cup$}}{\cup}{\cup}}\displaylimits}%
  \def\bigcap{\mathop{\mathchoice{\raisebox{-0.3em}{\scalebox{1.7}{$\cap$}}}{\scalebox{1.25}{$\cap$}}{\cap}{\cap}}\displaylimits}%
  \def\lesseqqgtr{\mathrel{\lessgtr}}\def\bigoplus{\mathop{\oplus}\displaylimits}%
}
\providecommand{\ssi}{si et seulement si} % abréviation fréquente
\AtBeginDocument{\providecommand{\wideparen}{\overparen}} % arc de cercle

%% FIGFIX-BEGIN v3 — correctifs globaux des figures (docs/onbuch-generation/tools/figfix)
\usepackage{adjustbox}\usepackage{etoolbox}
% toute figure TikZ/pgfplots plus large que la ligne est réduite à la largeur disponible
\BeforeBeginEnvironment{tikzpicture}{\begin{adjustbox}{max width=\linewidth}}
\AfterEndEnvironment{tikzpicture}{\end{adjustbox}}
% légendes pgfplots lisibles par-dessus les courbes
\pgfplotsset{every axis/.append style={legend style={fill=white,fill opacity=0.92,text opacity=1,draw=black!15,inner sep=3pt}}}
% étiquettes d'arêtes (syntaxe edge["..."]) avec fond blanc
\tikzset{every edge quotes/.append style={fill=white,inner sep=1.2pt}}
% bulles de cartes mentales : texte à la ligne dans la bulle, bulle un peu plus grande
\tikzset{every concept/.append style={text width=2.0cm,align=center,minimum size=2.3cm,inner sep=2pt}}
%% FIGFIX-END
\begin{document}

\pagegarde{Les outils numériques pédagogiques et l'intelligence artificielle}{Allemand}{Tle A}{Chapitre 4 : Média et communication}{AL19}{2 h}{Cette leçon propose la lecture et le commentaire d'un texte allemand sur les outils numériques pédagogiques et l'intelligence artificielle à l'école. Elle permet d'acquérir le vocabulaire thématique, de maîtriser le Futur I et les groupes infinitifs, et de produire un court texte argumenté.
\vspace{4pt}
\begin{multicols}{2}\small
\begin{itemize}
\item À la fin de cette leçon, je sais comprendre globalement puis en détail un texte allemand sur les outils numériques et l'intelligence artificielle.
\item À la fin de cette leçon, je sais employer correctement le vocabulaire du thème (die künstliche Intelligenz, das Lernprogramm, die App, das Tablet...) avec articles et pluriels.
\item À la fin de cette leçon, je sais former et employer le Futur I pour exprimer des prévisions et des projets.
\item À la fin de cette leçon, je sais construire et traduire des groupes infinitifs avec zu, um...zu, ohne...zu, anstatt...zu.
\item À la fin de cette leçon, je sais répondre par écrit et à l'oral à des questions de compréhension sur un texte.
\item À la fin de cette leçon, je sais résumer un texte allemand en quelques phrases avec mes propres mots.
\end{itemize}\end{multicols}}

\begin{sommaire}
\begin{multicols}{2}
\begin{enumerate}
\item Texte support : Lernen mit der KI
\item Compréhension du texte
\item Lexique thématique : digitale Medien und KI
\item Grammaire 1 : das Futur I
\item Grammaire 2 : die Infinitivgruppen
\item Méthode du commentaire et production guidée
\item Activité d'intégration
\item Méthodes et exercices résolus
\item Exercices
\item Corrigés des exercices
\item Fiche bilan
\end{enumerate}
\end{multicols}
\end{sommaire}

\begin{objectifs}
\begin{itemize}
\item À la fin de cette leçon, je sais comprendre globalement puis en détail un texte allemand sur les outils numériques et l'intelligence artificielle.
\item À la fin de cette leçon, je sais employer correctement le vocabulaire du thème (die künstliche Intelligenz, das Lernprogramm, die App, das Tablet...) avec articles et pluriels.
\item À la fin de cette leçon, je sais former et employer le Futur I pour exprimer des prévisions et des projets.
\item À la fin de cette leçon, je sais construire et traduire des groupes infinitifs avec zu, um...zu, ohne...zu, anstatt...zu.
\item À la fin de cette leçon, je sais répondre par écrit et à l'oral à des questions de compréhension sur un texte.
\item À la fin de cette leçon, je sais résumer un texte allemand en quelques phrases avec mes propres mots.
\item À la fin de cette leçon, je sais rédiger un court paragraphe argumenté (avantages/inconvénients) sur le numérique à l'école.
\item À la fin de cette leçon, je sais éviter les erreurs fréquentes des francophones (faux amis, place de l'infinitif, confusion Futur I / présent + adverbe).
\end{itemize}
\end{objectifs}

\begin{prerequis}
\begin{itemize}
\item La conjugaison des verbes faibles et forts au présent (Première)
\item Les auxiliaires werden, haben, sein au présent (Première)
\item La place du verbe en position 2 et l'infinitif en fin de phrase (Première)
\item Le vocabulaire scolaire de base : die Schule, der Unterricht, der Lehrer, lernen (Seconde/Première)
\item Les subordonnées introduites par dass, weil, wenn (début de Tle)
\end{itemize}
\end{prerequis}

\newpage

\coursec{Texte support : Lernen mit der KI}

\courssub{Situation d'accroche et problématique}

Au lycée de Yaoundé, Amina révise son bac d'allemand avec une application sur son smartphone : chaque soir, le logiciel lui propose des exercices de conjugaison adaptés à ses erreurs de la veille. Son correspondant bavarois Lukas, lui, utilise un programme d'apprentissage adaptatif qui corrige ses rédactions grâce à l'intelligence artificielle et lui explique ses fautes en quelques secondes. En Suisse, des élèves suivent leurs cours sur tablette ; en Autriche, des enseignants débattent dans les salles des professeurs : \all{Wird die KI den Lehrer ersetzen?} — « L'IA va-t-elle remplacer le professeur ? »

Ces outils numériques transforment l'école partout dans le monde, de Douala à Berlin. Mais quels sont leurs véritables avantages ? Quels dangers cachent-ils ? C'est la \cle{problématique} de cette leçon : comprendre un texte allemand sur l'école numérique, enrichir notre vocabulaire des médias, et préparer deux outils grammaticaux essentiels pour en parler : le \cle{Futur I} (pour exprimer l'avenir) et les \cle{groupes infinitifs} (pour construire des phrases complexes).

\courssub{Lecture globale : première approche du texte}

\begin{methode}[Lire un texte allemand au Bac : les 5 étapes]
\begin{enumerate}
\item Lire le \cle{titre} et la \cle{source} : ils annoncent le thème et le type de texte (article, dialogue, lettre).
\item Faire une \cle{première lecture rapide}, sans dictionnaire, pour saisir l'idée générale. Ne pas bloquer sur un mot inconnu.
\item Souligner les \cle{mots-clés} : noms (majuscules !), verbes, chiffres, connecteurs (\all{aber}, \all{also}, \all{um ... zu}).
\item Faire une \cle{deuxième lecture détaillée} avec le glossaire : vérifier le sens des mots soulignés.
\item \cle{Reformuler} l'idée principale en une phrase, en français ou en allemand : \all{Der Text handelt von ...} (« Le texte traite de ... »).
\end{enumerate}
\end{methode}

Appliquons l'étape 1 : le titre \all{Lernen mit der künstlichen Intelligenz} annonce un texte d'\emph{information et de réflexion} sur l'école ; la source indiquée est un magazine scolaire allemand. Le thème est donc prévisible : les nouvelles technologies dans l'enseignement.

\courssub{Le texte : Lernen mit der künstlichen Intelligenz}

\begin{textebox}[Aus einem Schülermagazin : Lernen mit der künstlichen Intelligenz]
In vielen Schulen in Deutschland, Österreich und der Schweiz lernen die Schüler heute mit Tablets, Computern und Apps. Ein Lernprogramm mit künstlicher Intelligenz passt die Aufgaben an das Niveau jedes Schülers an und korrigiert die Fehler sofort. So kann jeder Jugendliche in seinem eigenen Tempo arbeiten: Wer schnell lernt, bekommt schwierigere Übungen; wer mehr Zeit braucht, wiederholt die Lektion ohne Stress.

Viele Experten glauben: Die KI wird den Unterricht in den nächsten Jahren stark verändern, aber sie wird den Lehrer nicht ersetzen. Denn ein Computer kann Grammatik erklären, aber er kann keinen Schüler trösten, der traurig ist, und keine Klasse motivieren, die müde ist. Um gut zu lernen, braucht der Mensch auch Kontakt zu anderen Menschen.

Natürlich gibt es auch Probleme. Wer nur am Bildschirm lernt, ohne mit anderen zu sprechen, verliert oft die Motivation und fühlt sich isoliert. Außerdem sammeln viele Apps persönliche Daten der Nutzer, und nicht jede Familie kann ein teures Tablet kaufen. Die Schule der Zukunft wird also beides verbinden: moderne Technik und menschliche Beziehungen. Lehrer und Schüler werden gemeinsam entscheiden, wann der Computer hilft — und wann ein Gespräch wichtiger ist.
\end{textebox}

\textbf{Lecture à voix haute et prononciation.} L'enseignant lit le texte une première fois ; les élèves écoutent, livre fermé, puis répètent les mots difficiles :

\begin{center}
\begin{tabularx}{\linewidth}{|X|X|X|}\hline
\rowcolor{popblueL} \textbf{Mot allemand} & \textbf{Prononciation (lettres françaises)} & \textbf{Remarque} \\ \hline
\all{künstlich} & « künst-lich » (ü comme dans « tu ») & le \all{ch} final se prononce « h » doux \\ \hline
\all{Intelligenz} & « in-té-li-guèn-ts » & \all{-enz} = « ents », jamais « ense » \\ \hline
\all{Lernprogramm} & « lèr-n-pro-gram » & mot composé : \all{lernen} + \all{Programm} \\ \hline
\all{Bildschirm} & « bild-chirm » & \all{sch} = « ch » français \\ \hline
\all{verändern} & « fèr-èn-dèrn » & \all{v} = « f » ; \all{ä} = « è » \\ \hline
\all{Gespräch} & « guè-chprèh » & le \all{ch} après \all{ä} est doux \\ \hline
\end{tabularx}
\end{center}

\courssub{Glossaire et vocabulaire du texte}

\begin{definition}[Die künstliche Intelligenz (KI)]
La \cle{künstliche Intelligenz} (\all{die}, abrégé \all{die KI}) désigne la capacité d'un programme informatique à réaliser des tâches qui demandent normalement l'intelligence humaine : comprendre une langue, corriger une erreur, adapter un exercice au niveau de l'élève. En français : l'intelligence artificielle (IA).
\end{definition}

\begin{center}
\begin{tabularx}{\linewidth}{|X|X|}\hline
\rowcolor{popblueL} \textbf{Deutsch (article + pluriel)} & \textbf{Français} \\ \hline
\all{die künstliche Intelligenz, -en} & l'intelligence artificielle \\ \hline
\all{das Lernprogramm, -e} & le logiciel / programme d'apprentissage \\ \hline
\all{der Bildschirm, -e} & l'écran \\ \hline
\all{die Aufgabe, -n} & la tâche, l'exercice \\ \hline
\all{das Tablet, -s} & la tablette \\ \hline
\all{die Übung, -en} & l'exercice (d'entraînement) \\ \hline
\all{das Niveau, -s} & le niveau \\ \hline
\all{der Fehler, -} & la faute, l'erreur \\ \hline
\all{die Motivation (Sg.)} & la motivation \\ \hline
\all{die Beziehung, -en} & la relation \\ \hline
\all{das Gespräch, -e} & la conversation \\ \hline
\all{die Daten (Pl. von \all{das Datum})} & les données \\ \hline
\all{anpassen (trennbar)} & adapter \\ \hline
\all{ersetzen} & remplacer \\ \hline
\all{verändern} & changer, modifier \\ \hline
\all{trösten} & consoler \\ \hline
\all{sammeln} & collecter \\ \hline
\all{verbinden (verband, hat verbunden)} & combiner, relier \\ \hline
\all{sofort} & immédiatement \\ \hline
\all{außerdem} & de plus, en outre \\ \hline
\end{tabularx}
\end{center}

\begin{attention}[Faux amis : attention aux traductions automatiques]
\begin{itemize}
\item \all{die App} se traduit bien par « l'application » — aucun piège ici.
\item En revanche, \all{das Lernprogramm} n'est PAS « le programme scolaire » : le programme scolaire officiel se dit \all{der Lehrplan}. \all{Das Lernprogramm} est un \emph{logiciel d'apprentissage}.
\item \all{die Daten} = « les données » (informatiques), et non « les dates » : « les dates » se dit \all{die Termine} ou \all{die Daten} uniquement au sens calendaire — le contexte décide.
\item \all{künstlich} = « artificiel » (et non « artistique », qui se dit \all{künstlerisch}).
\end{itemize}
\end{attention}

\begin{savaistu}[KI ou IA ?]
En Allemagne, on abrège \all{künstliche Intelligenz} en \all{KI}, prononcé « ka-i » (on énonce les deux lettres à l'allemande). Le français dit « IA » en inversant l'ordre des mots. Dans une traduction thème ou version, il faut donc penser à inverser : \all{Die KI entwickelt sich schnell.} = « L'IA se développe rapidement. » Les pays germanophones organisent d'ailleurs de nombreux débats scolaires (\all{die Debatte, -n}) sur la place de la KI dans les examens — un sujet qui pourrait bien tomber au Bac !
\end{savaistu}

\courssub{Repérage grammatical dans le texte}

Avant d'étudier la grammaire en détail (sections 4 et 5), repérons dans le texte les deux structures au programme.

\textbf{1. Les occurrences du Futur I} (\all{werden} + infinitif en fin de phrase) :

\begin{exemplebox}[Le Futur I dans le texte]
\begin{itemize}
\item \all{Die KI wird den Unterricht stark verändern.} « L'IA va fortement transformer les cours. »
\item \all{Aber sie wird den Lehrer nicht ersetzen.} « Mais elle ne remplacera pas le professeur. »
\item \all{Die Schule der Zukunft wird also beides verbinden.} « L'école du futur combinera donc les deux. »
\item \all{Lehrer und Schüler werden gemeinsam entscheiden.} « Professeurs et élèves décideront ensemble. »
\end{itemize}
Observation : le verbe conjugué \all{wird/werden} occupe la \cle{deuxième position} ; l'infinitif est rejeté à la \cle{fin de la phrase}.
\end{exemplebox}

\textbf{2. Les groupes infinitifs} (\all{um ... zu} + infinitif, ou infinitif avec \all{zu}) :

\begin{exemplebox}[Les groupes infinitifs dans le texte]
\begin{itemize}
\item \all{Um gut zu lernen, braucht der Mensch auch Kontakt zu anderen Menschen.} « Pour bien apprendre, l'être humain a aussi besoin de contact avec les autres. »
\item \all{Wer nur am Bildschirm lernt, ohne mit anderen zu sprechen, verliert oft die Motivation.} « Celui qui apprend seulement devant l'écran, sans parler avec les autres, perd souvent sa motivation. »
\end{itemize}
Observation : le groupe infinitif est encadré par une virgule ; l'infinitif précédé de \all{zu} se place à la fin du groupe.
\end{exemplebox}

\courssub{Compréhension globale du texte}

\begin{exoresolu}[Questions de compréhension globale]
Énoncé. Répondez en français, puis formulez une réponse courte en allemand :
\begin{enumerate}
\item De quoi parle le texte ?
\item Qui parle, et d'où vient le texte ?
\item Quelle est la thèse de l'auteur ?
\end{enumerate}
\tcblower
Solution détaillée :
\begin{enumerate}
\item Le texte parle de l'utilisation des outils numériques (tablettes, ordinateurs, applications) et de l'intelligence artificielle dans les écoles des pays germanophones, avec leurs avantages et leurs limites. En allemand : \all{Der Text handelt vom Lernen mit der KI in der Schule.}
\item C'est un article de magazine scolaire (\all{Schülermagazin}) : un journaliste ou un expert s'adresse à des jeunes lecteurs. Le ton est informatif et équilibré.
\item La thèse de l'auteur : la KI va transformer l'école, mais elle ne remplacera pas le professeur ; l'école de demain doit combiner technologie et relations humaines. En allemand : \all{Die KI wird den Unterricht verändern, aber sie wird den Lehrer nicht ersetzen.}
\end{enumerate}
\end{exoresolu}

\courssub{Compréhension détaillée : avantages et dangers}

Construisons le bilan du texte sous forme de tableau — une excellente méthode pour préparer un résumé ou un débat.

\begin{center}
\begin{tabularx}{\linewidth}{|X|X|}\hline
\rowcolor{popgreenL} \textbf{Vorteile (avantages)} & \textbf{Nachteile / Gefahren (dangers)} \\ \hline
\all{individuelles Lernen} : les exercices s'adaptent au niveau de chaque élève & \all{Abhängigkeit} : risque de dépendance à l'écran \\ \hline
\all{sofortige Korrektur} : correction immédiate des fautes & \all{Isolation} : isolement, perte de motivation sans contact humain \\ \hline
\all{eigenes Tempo} : chacun travaille à son rythme, sans stress & \all{persönliche Daten} : les applications collectent des données personnelles \\ \hline
\all{schwierigere Übungen} pour les élèves rapides & \all{teure Geräte} : toutes les familles ne peuvent pas acheter une tablette \\ \hline
\end{tabularx}
\end{center}

\begin{exercice}[Compréhension détaillée — Richtig oder falsch ?]
Dites si les phrases suivantes sont vraies (\all{richtig}) ou fausses (\all{falsch}) d'après le texte, et justifiez par une citation :
\begin{enumerate}
\item \all{Das Lernprogramm korrigiert die Fehler nach einer Woche.}
\item \all{Die KI kann einen traurigen Schüler trösten.}
\item \all{Viele Apps sammeln persönliche Daten.}
\item \all{Die Schule der Zukunft wird nur mit Computern arbeiten.}
\end{enumerate}
\end{exercice}

\begin{corrige}[Exercice : Richtig oder falsch ?]
\begin{enumerate}
\item \all{Falsch.} Le texte dit : \all{... und korrigiert die Fehler sofort} — la correction est immédiate, pas après une semaine.
\item \all{Falsch.} Le texte dit : \all{Ein Computer ... kann keinen Schüler trösten, der traurig ist} — l'ordinateur ne peut pas consoler.
\item \all{Richtig.} Le texte dit : \all{Außerdem sammeln viele Apps persönliche Daten der Nutzer.}
\item \all{Falsch.} Le texte dit : \all{Die Schule der Zukunft wird also beides verbinden: moderne Technik und menschliche Beziehungen.}
\end{enumerate}
\end{corrige}

\courssub{Carte mentale : organiser ses idées}

La carte mentale ci-dessous résume tout le contenu du texte. Reproduisez-la dans votre cahier : c'est un excellent support pour le résumé et pour le débat oral de la section 6.

\begin{popfigure}
\begin{center}
\begin{tikzpicture}[mindmap, grow cyclic,
  every node/.style={concept, align=center, font=\small},
  root concept/.append style={concept color=popblue, text=white, font=\bfseries},
  level 1 concept/.append style={level distance=3.4cm, sibling angle=90, font=\small}]
\node[root concept]{KI in der Schule}
  child[concept color=popgreen] { node[align=center]{Vorteile\\ individuell\\ sofortige\\ Korrektur} }
  child[concept color=poporange] { node[align=center]{Nachteile\\ Abhängigkeit\\ Isolation\\ Daten} }
  child[concept color=poppurple] { node[align=center]{Geräte\\ Tablet\\ Computer\\ App} }
  child[concept color=popgold] { node[align=center]{Zukunft\\ Technik +\\ Menschen\\ verbinden} };
\end{tikzpicture}
\end{center}
\legende{Carte mentale du texte : les quatre axes de réflexion sur la KI à l'école.}
\end{popfigure}

\begin{aretenir}[Ce qu'il faut retenir de cette première lecture]
\begin{itemize}
\item Le texte est un \cle{article d'information} équilibré : il présente les avantages (apprentissage individualisé, correction immédiate) et les dangers (dépendance, isolement, données personnelles) des outils numériques.
\item La thèse de l'auteur : la technologie transforme l'école, mais ne remplace pas le professeur — \all{Die KI wird den Lehrer nicht ersetzen.}
\item Le vocabulaire clé : \all{die künstliche Intelligenz (KI)}, \all{das Lernprogramm}, \all{der Bildschirm}, \all{die Aufgabe}, \all{anpassen}, \all{ersetzen}.
\item Nous avons repéré les deux structures grammaticales de la leçon : le \cle{Futur I} (\all{wird verändern}) et les \cle{groupes infinitifs} (\all{um gut zu lernen}, \all{ohne zu sprechen}) — nous allons maintenant les étudier en détail.
\end{itemize}
\end{aretenir}

\coursec{Les outils numériques pédagogiques et l'intelligence artificielle}

\begin{textebox}[Texte support : Lernen mit der KI – Chance oder Risiko?]
In vielen Schulen in Yaoundé und Douala gehören Tablets und Lern-Apps heute zum Alltag. Die Schülerinnen und Schüler arbeiten mit Programmen, die künstliche Intelligenz (KI) nutzen. Diese Programme analysieren die Fehler der Lernenden und schlagen passende Übungen vor. \textbf{So} kann jeder in seinem eigenen Tempo lernen.

\textbf{Aber} die Digitalisierung hat auch Schattenseiten. Wer nur am Bildschirm sitzt, verliert oft den Kontakt zu Mitschülern und Lehrern. Ein Lehrer kann Emotionen erkennen und motivieren; ein Algorithmus kann das nicht. \textbf{Deshalb} warnen Experten vor einer reinen „Bildschirm-Schule“.

\textbf{Um} gut zu lernen, \textbf{braucht} der Mensch echte Begegnung. Die KI ist ein starkes Werkzeug, \textbf{aber} sie ersetzt nicht den Pädagogen. \textbf{Also} sollten Schulen auf ein hybrides Modell setzen: digitale Werkzeuge \textbf{und} menschliche Begleitung. Nur so profitieren alle vom Fortschritt.
\end{textebox}

\begin{definition}[Glossaire : Mots difficiles du texte]
\begin{itemize}
    \item \all{der Alltag, die Alltage} : le quotidien
    \item \all{das Lernprogramm, die Lernprogramme} : le logiciel éducatif, le programme d'apprentissage
    \item \all{analysieren} : analyser
    \item \all{der Fehler, die Fehler} : l'erreur, la faute
    \item \all{vorschlagen (schlägt vor, hat vorgeschlagen)} : proposer (verbe à particule séparable)
    \item \all{das Tempo, die Tempi} : le rythme, la vitesse
    \item \all{die Schattenseite, die Schattenseiten} : le revers de la médaille, l'inconvénient
    \item \all{der Kontakt, die Kontakte} : le contact
    \item \all{erkennen (erkennt, hat erkannt)} : reconnaître, identifier (voyelle \cle{e $\rightarrow$ ie} à la 2e/3e pers. sg.)
    \item \all{motivieren} : motiver
    \item \all{der Algorithmus, die Algorithmen} : l'algorithme
    \item \all{warnt vor (+ Dat.)} : mettre en garde contre (régence datif)
    \item \all{die Begegnung, die Begegnungen} : la rencontre
    \item \all{das Werkzeug, die Werkzeuge} : l'outil
    \item \all{ersetzen (ersetzt, hat ersetzt)} : remplacer
    \item \all{der Pädagoge, die Pädagogen} : l'éducateur, le pédagogue
    \item \all{das hybride Modell, die hybriden Modelle} : le modèle hybride
    \item \all{profitieren von (+ Dat.)} : profiter de, bénéficier de (régence datif)
    \item \all{der Fortschritt} : le progrès (gén. sans pluriel dans ce sens)
\end{itemize}
\end{definition}

\courssub{Compréhension globale : Worum geht es? Wer spricht?}

Avant de plonger dans les détails, il faut saisir l'intention communicative du texte. Au Baccalauréat, les premières questions portent souvent sur le \textbf{sujet}, la \textbf{source} (qui parle ?), le \textbf{type de texte} et le \textbf{ton}.

\begin{methode}[Répondre à une question de compréhension globale — 3 étapes]
\begin{enumerate}
    \item \textbf{Repérer les indices} : titre, premier paragraphe (amorce), dernier paragraphe (conclusion), mots-clés répétés (\cle{KI}, \cle{Schule}, \cle{Lehrer}, \cle{Experten}).
    \item \textbf{Reformuler avec ses propres mots} (en allemand) : ne recopiez jamais une phrase entière du texte. Utilisez des synonymes et changez la structure (ex. : passage du discours direct au rapporté, nominalisation).
    \item \textbf{Répondre par une phrase complète} : Sujet + Verbe conjugué en position 2 + Compléments. En compréhension écrite, on répond généralement au \textbf{Präsens} (actualité du texte) ou au \textbf{Perfekt} (résumé d'un événement passé).
\end{enumerate}
\end{methode}

\begin{exemplebox}[Modèles de réponses globales]
\textbf{Frage 1: Worum geht es im Text?} \\
\all{Der Text handelt von der Nutzung von künstlicher Intelligenz (KI) im Schulunterricht in Kamerun und diskutiert Vor- und Nachteile digitaler Lernwerkzeuge.} \\
« Le texte traite de l'utilisation de l'intelligence artificielle dans l'enseignement au Cameroun et discute les avantages et inconvénients des outils numériques d'apprentissage. »

\textbf{Frage 2: Wer spricht im Text? (Wer ist der Autor / die Autorin?)} \\
\all{Ein Beobachter der schulischen Digitalisierung (Experte / Journalist) analysiert die Situation objektiv und gibt eine Empfehlung.} \\
« Un observateur de la numérisation scolaire (expert / journaliste) analyse la situation objectivement et donne une recommandation. »

\textbf{Frage 3: Welche Textsorte liegt vor?} \\
\all{Es ist ein argumentativer Kommentar / eine Stellungnahme zu einem aktuellen Bildungsthema.} \\
« C'est un commentaire argumentatif / une prise de position sur un thème éducatif actuel. »
\end{exemplebox}

\begin{exercice}[Compréhension globale — À vous de jouer]
Répondez en allemand par des phrases complètes.
\begin{enumerate}
    \item Worum geht es im Text? (Thema + Kontext Kamerun)
    \item Welche These vertritt der Autor am Ende des Textes?
    \item Nennen Sie zwei Orte, die im ersten Satz genannt werden.
\end{enumerate}
\end{exercice}

\begin{corrige}[Exercice Compréhension globale]
\begin{enumerate}
    \item \all{Der Text thematisiert den Einsatz von KI-gestützten Lernprogrammen an Schulen in Yaoundé und Douala und wägt Chancen und Risiken ab.}
    \item \all{Der Autor plädiert für ein hybrides Lernmodell, das digitale Werkzeuge und menschliche Pädagogik verbindet.}
    \item \all{Yaoundé und Douala.}
\end{enumerate}
\end{corrige}

\courssub{Compréhension détaillée : Geräte, Funktionen, Meinungen}

Les questions de détail exigent une lecture fine. Il faut localiser l'information précise (paragraphe, ligne), la comprendre et la reformuler.

\begin{exemplebox}[Questions de détail — Modèles de réponses]
\textbf{1. Welche Geräte benutzen die Schülerinnen und Schüler?} \\
\all{Sie benutzen Tablets und Laptops (bzw. Computer).} \\
(Repère : §1, „Tablets und Lern-Apps“ / „Lernprogramme“)

\textbf{2. Was macht das Lernprogramm mit der KI genau?} \\
\all{Es analysiert die Fehler der Lernenden und schlägt passende Übungen vor, damit jeder in seinem eigenen Tempo lernen kann.} \\
(Repère : §1, phrases 2-3. Notez la reformulation : „schlägt ... vor“ $\rightarrow$ „propose ...“ ; „in seinem eigenen Tempo“ gardé.)

\textbf{3. Was glauben die Experten? / Wovor warnen die Experten?} \\
\all{Experten warnen vor einer reinen „Bildschirm-Schule“, weil ein Algorithmus keine Emotionen erkennen und nicht motivieren kann.} \\
(Repère : §2. Structure \cle{warnen vor + Dativ} + justification par \cle{weil} + verbe à la fin.)
\end{exemplebox}

\begin{attention}[Piège classique : Recopier vs Reformuler]
\begin{itemize}
    \item \textbf{Interdit} : Recopier le segment exact du texte sans guillemets ni adaptation. Ex. : \all{„Diese Programme analysieren die Fehler der Lernenden und schlagen passende Übungen vor.“} $\leftarrow$ \textbf{Zéro point} au Bac (pas de preuve de compréhension).
    \item \textbf{Attendu} : Reformulation (changement de voix, synonymes, structure). Ex. : \all{Die KI-Programme werten Fehler aus und empfehlen individuelle Aufgaben.}
    \item \textbf{Citation autorisée} : Uniquement si la consigne demande : \all{Zitieren Sie die Stelle im Text, die zeigt, dass...} $\rightarrow$ \all{„Ein Lehrer kann Emotionen erkennen und motivieren; ein Algorithmus kann das nicht.“}
\end{itemize}
\end{attention}

\courssub{Exercice type Bac : Vrai ou Faux justifié (Richtig / Falsch)}

C'est l'exercice roi du Baccalauréat camerounais. La note ne vient pas du « R » ou « F », mais de la \textbf{justification par citation exacte} (entre guillemets allemands) + \textbf{explication en propre mots}.

\begin{exoresolu}[Vrai ou Faux justifié]
\textbf{Consigne :} Entscheiden Sie: Richtig oder Falsch? Begründen Sie mit einem Zitat aus dem Text und einer kurzen Erklärung auf Deutsch.

\begin{enumerate}
    \item \textbf{Die KI-Programme ersetzen die Lehrer vollständig.}
    \item \textbf{Schüler in Yaoundé lernen schneller mit Tablets als ohne.}
    \item \textbf{Ein Algorithmus kann die Motivation der Schüler nicht wie ein Mensch fördern.}
    \item \textbf{Der Autor ist gegen jede Digitalisierung der Schule.}
\end{enumerate}

\tcblower

\textbf{Solution détaillée :}

\begin{enumerate}
    \item \textbf{Falsch.} \\
    \textbf{Zitat:} „Die KI ist ein starkes Werkzeug, \textbf{aber} sie ersetzt nicht den Pädagogen.“ \\
    \textbf{Erklärung:} Der Text sagt explizit, dass die KI den Pädagogen \textbf{nicht} ersetzt, sondern ein Werkzeug ist. Das Adverb „komplett“ macht die Aussage falsch.

    \item \textbf{Falsch.} \\
    \textbf{Zitat:} \emph{(Kein Zitat möglich, Information nicht im Text)} \\
    \textbf{Erklärung:} Der Text erwähnt, dass Schüler \emph{in ihrem eigenen Tempo} lernen können („So kann jeder in seinem eigenen Tempo lernen“), aber nicht, dass sie \emph{schneller} lernen. „Eigenes Tempo“ $\neq$ „schneller“. Vorsicht vor Leseprojektion!

    \item \textbf{Richtig.} \\
    \textbf{Zitat:} „Ein Lehrer kann Emotionen erkennen und motivieren; \textbf{ein Algorithmus kann das nicht}.“ \\
    \textbf{Erklärung:} Der Text stellt einen klaren Kontrast her: Der Mensch (Lehrer) besitzt emotionale Kompetenzen (Empathie, Motivation), die der Algorithmus nicht hat.

    \item \textbf{Falsch.} \\
    \textbf{Zitat:} „\textbf{Also} sollten Schulen auf ein hybrides Modell setzen: digitale Werkzeuge \textbf{und} menschliche Begleitung.“ \\
    \textbf{Erklärung:} Der Autor plädiert für eine \textbf{Kombination} (hybrides Modell), nicht für eine Ablehnung der Digitalisierung. Er sieht die KI als „starkes Werkzeug“.
\end{enumerate}
\end{exoresolu}

\courssub{Grammaire 1 : Das Futur I (Le futur)}

Le Futur I exprime une action future, une supposition ou une promesse. C'est un temps composé : auxiliaire \cle{werden} + infinitif du verbe principal.

\begin{propriete}[Formation du Futur I]
\begin{center}
\begin{tabular}{|l|l|l|l|}
\hline
\rowcolor{popblueL}
\textbf{Person} & \textbf{werden (Präsens)} & \textbf{Infinitif (ex. lernen)} & \textbf{Traduction} \\
\hline
ich & werde & lernen & j'apprendrai / j'apprendrai \\
du & wirst & lernen & tu apprendras \\
er / sie / es & wird & lernen & il/elle apprendra \\
wir & werden & lernen & nous apprendrons \\
ihr & werdet & lernen & vous apprendrez \\
sie / Sie & werden & lernen & ils/elles / Vous apprendrez \\
\hline
\end{tabular}
\end{center}
\textbf{Règle d'ordre :} Verbe conjugué \cle{werden} en \textbf{position 2} ; Infinitif à la \textbf{fin de phrase} (position finale).
\all{Morgen \textbf{werde} ich Deutsch \textbf{lernen}.} (Demain j'apprendrai l'allemand.)
\end{propriete}

\begin{propriete}[Emplois du Futur I]
\begin{enumerate}
    \item \textbf{Futur chronologique} (action future, souvent avec marqueur temporel : \cle{morgen}, \cle{nächstes Jahr}, \cle{bald}, \cle{in Zukunft}).
    \all{Im nächsten Jahr \textbf{werden} mehr Schulen Tablets \textbf{nutzen}.}
    \item \textbf{Supposition / Probabilité} (présent ou futur), souvent avec \cle{wohl}, \cle{schon}, \cle{bestimmt}, \cle{sicher}.
    \all{Das \textbf{wird} wohl \textbf{klappen}.} (Ça marchera probablement / Ça va marcher.)
    \item \textbf{Promesse / Menace / Détermination} (style oral ou solennel).
    \all{Ich \textbf{werde} das \textbf{tun}!} (Je le ferai !)
\end{enumerate}
\end{propriete}

\begin{attention}[Pièges fréquents pour francophones — Futur I]
\begin{itemize}
    \item \textbf{\cle{werden} $\neq$ \cle{bekommen}} : \all{Ich werde ein Geschenk bekommen.} (Je \textbf{recevrai} un cadeau). \cle{werden} = auxiliaire futur ; \cle{bekommen} = recevoir. \textbf{Faux ami} : \all{Ich werde krank} = Je \textbf{deviens} malade (devenir), pas « je recevrai malade ».
    \item \textbf{Präsens vs Futur I} : L'allemand utilise souvent le \textbf{Präsens} + indicateur de temps pour le futur proche ou certain. \all{Morgen fahre ich nach Douala.} (Demain je vais à Douala — certain/planifié). Le Futur I marque plus la prédiction ou l'intention non planifiée.
    \item \textbf{Négation} : \cle{nicht} se place \textbf{avant l'infinitif} (ou avant le complément d'objet si pronominal). \all{Ich werde das \textbf{nicht} machen.} / \all{Ich werde es \textbf{nicht} tun.}
    \item \textbf{Verbes à particule séparable} : La particule reste collée à l'infinitif à la fin. \all{Er wird das Licht \textbf{anmachen}.} (Il allumera la lumière).
\end{itemize}
\end{attention}

\begin{exercice}[Futur I — Entraînement]
\begin{enumerate}
    \item Conjuguez au Futur I : \all{schreiben} (écrire), \all{arbeiten} (travailler), \all{anfangen} (commencer).
    \item Traduisez en allemand :
    \begin{enumerate}
        \item Demain, les élèves utiliseront des tablettes.
        \item L'intelligence artificielle remplacera probablement pas le professeur.
        \item Nous allons apprendre l'allemand l'année prochaine. (Intention/Plan $\rightarrow$ Präsens ou Futur ? Justifiez.)
    \end{enumerate}
\end{enumerate}
\end{exercice}

\begin{corrige}[Exercice Futur I]
\begin{enumerate}
    \item \all{ich werde schreiben, du wirst schreiben, er wird schreiben, wir werden schreiben, ihr werdet schreiben, sie werden schreiben} ; \all{ich werde arbeiten...} ; \all{ich werde anfangen...} (particule \cle{an-} à la fin : \all{... anfangen}).
    \item 
    \begin{enumerate}
        \item \all{Morgen werden die Schüler Tablets nutzen.} (ou \all{... Tablets benutzen.})
        \item \all{Die künstliche Intelligenz wird den Lehrer wohl nicht ersetzen.} (Supposition $\rightarrow$ \cle{wohl} + Futur I).
        \item \all{Nächstes Jahr lernen wir Deutsch.} (Präsens + marqueur temporel = plan fixé). \all{Nächstes Jahr werden wir Deutsch lernen.} (Prédiction/Intention moins ferme). Les deux sont corrects selon la nuance.
    \end{enumerate}
\end{enumerate}
\end{corrige}

\courssub{Grammaire 2 : Die Infinitivgruppen (Groupes infinitifs)}

Une \textbf{Infinitivgruppe} est une construction réduite (sans sujet propre, verbe à l'infinitif + \cle{zu}) qui remplace souvent une proposition subordonnée (\cle{dass}, \cle{damit}, \cle{um ... zu}). Elle est très fréquente en allemand écrit (presse, essais, Bac).

\begin{propriete}[Structure et Ponctuation]
\begin{itemize}
    \item \textbf{Forme de base :} \cle{zu} + \textbf{Infinitif} (à la fin de la phrase / du groupe).
    \item \textbf{Sujet :} Identique au sujet de la principale (contrôle de sujet). \all{Ich versuche, \textbf{zu verstehen}.} (J'essaie de comprendre — moi je comprends).
    \item \textbf{Virgule :} \textbf{Recommandée} (nouvelle orthographe) devant \cle{zu} / \cle{um ... zu} / \cle{ohne ... zu} / \cle{statt ... zu} / \cle{anstatt ... zu} pour la lisibilité. \textbf{Obligatoire} si le groupe dépend d'un nom (\cle{das Vorhaben, zu gehen}), d'un adjectif (\cle{schön, zu sehen}) ou d'un verbe corrélé (\cle{es geht um, zu lernen}).
    \item \textbf{Verbes à particule séparable :} \cle{zu} s'intercale entre la particule et le radical. \all{aufhören $\rightarrow$ auf\textbf{zu}hören} ; \all{vorschlagen $\rightarrow$ vor\textbf{zu}schlagen}.
    \item \textbf{Verbes en \cle{-ern}, \cle{-eln}} : \cle{e} muet tombe souvent. \all{handeln $\rightarrow$ zu handeln} ; \all{lächeln $\rightarrow$ zu lächeln}.
\end{itemize}
\end{propriete}

\begin{propriete}[Types principaux d'Infinitivgruppen]
\begin{center}
\begin{tabularx}{\linewidth}{|X|X|X|}
\hline
\rowcolor{popgreenL}
\textbf{Type / Préposition} & \textbf{Fonction / Sens} & \textbf{Exemple (du texte / construit)} \\
\hline
\cle{zu} + Inf. (seul) & Complément objet / sujet (verbes : \cle{versuchen}, \cle{hoffen}, \cle{versprechen}, \cle{vergessen}, \cle{planen}) & \all{Die Schule plant, \textbf{Tablets anzuschaffen}.} \\
\hline
\cle{um ... zu} + Inf. & \textbf{But final} (pour + infinitif) — \textbf{Sujet identique} & \all{\textbf{Um} gut \textbf{zu lernen}, braucht der Mensch Begegnung.} (Texte §3) \\
\hline
\cle{ohne ... zu} + Inf. & \textbf{Manière / Moyen négatif} (sans + infinitif) & \all{Er ging, \textbf{ohne sich zu verabschieden}.} \\
\hline
\cle{(an)statt ... zu} + Inf. & \textbf{Alternative} (au lieu de + infinitif) & \all{\textbf{Statt zu lernen}, spielt er.} \\
\hline
Infinitivgruppe comme \textbf{Sujet} & En tête de phrase (Position 1) & \all{\textbf{Deutsch zu lernen} ist wichtig.} \\
\hline
\end{tabularx}
\end{center}
\end{propriete}

\begin{methode}[Transformer une subordonnée en Infinitivgruppe (Bac : Umformung)]
\textbf{Condition sine qua non :} Le sujet de la subordonnée = Sujet de la principale.
\begin{enumerate}
    \item Supprimer la conjonction (\cle{dass}, \cle{weil}, \cle{damit}, \cle{obwohl}...).
    \item Mettre le verbe conjugué à l'\textbf{infinitif} + \cle{zu} (ou \cle{um ... zu} si but).
    \item Placer le groupe à la fin (ou au début si thème).
    \item Adapter la ponctuation (virgule).
\end{enumerate}
\all{Weil er gut lernen will, nutzt er die App.} $\rightarrow$ \all{\textbf{Um} gut \textbf{zu lernen}, nutzt er die App.} (But)
\all{Er hofft, dass er besteht.} $\rightarrow$ \all{Er hofft, \textbf{zu bestehen}.} (Verbe \cle{hoffen} + \cle{zu})
\end{methode}

\begin{attention}[Erreurs fatales — Infinitivgruppen]
\begin{itemize}
    \item \textbf{Sujets différents} $\rightarrow$ \textbf{Impossible} de faire une Infinitivgruppe. \all{Ich möchte, \textbf{dass du kommst}.} (Pas *\all{Ich möchte, zu kommen} — cela voudrait dire « Je veux venir moi-même »).
    \item \textbf{\cle{für ... zu} n'existe PAS} pour le but. Faux ami de \emph{pour + infinitif}. \textbf{Toujours} \cle{um ... zu}.
    \item \textbf{Verbe modal + Infinitif} : Pas de \cle{zu}. \all{Er kann gut deutsch sprechen.} (Pas *\all{zu sprechen}).
    \item \textbf{Ordre avec verbes à particule} : \all{Er verspricht, \textbf{auf\textbf{zu}hören}.} (Particule \cle{auf} + \cle{zu} + radical \cle{hören}).
\end{itemize}
\end{attention}

\begin{exercice}[Infinitivgruppen — Transformation \& Complétion]
\begin{enumerate}
    \item Transformez en Infinitivgruppe (\cle{um ... zu} ou \cle{zu}) :
    \begin{enumerate}
        \item \all{Die Schüler benutzen Apps, damit sie besser lernen.}
        \item \all{Der Lehrer hofft, dass die KI hilft.}
        \item \all{Man lernt nicht gut, wenn man nur am Bildschirm sitzt.} (Indice : \cle{ohne ... zu})
    \end{enumerate}
    \item Complétez avec la forme correcte (infinitif + \cle{zu} / \cle{um ... zu} / \cle{ohne ... zu}) :
    \begin{enumerate}
        \item Die KI hilft, Fehler \_\_\_\_\_\_\_\_ (korrigieren).
        \item \_\_\_\_\_\_\_\_ (verstehen), muss man zuhören.
        \item Er ging weg, \_\_\_\_\_\_\_\_ (grüßen).
    \end{enumerate}
\end{enumerate}
\end{exercice}

\begin{corrige}[Exercice Infinitivgruppen]
\begin{enumerate}
    \item 
    \begin{enumerate}
        \item \all{Die Schüler benutzen Apps, \textbf{um besser zu lernen}.} (But $\rightarrow$ \cle{um ... zu})
        \item \all{Der Lehrer hofft, \textbf{dass die KI hilft}.} $\rightarrow$ \textbf{Impossible} (sujets différents : \cle{Lehrer} vs \cle{KI}). Reste : \all{Der Lehrer hofft, dass die KI hilft.} (Pas d'infinitif possible).
        \item \all{Man lernt nicht gut, \textbf{ohne nur am Bildschirm zu sitzen}.} (\cle{ohne ... zu} + verbe \cle{sitzen}).
    \end{enumerate}
    \item 
    \begin{enumerate}
        \item \all{zu korrigieren} (verbe \cle{helfen} + \cle{zu} + infinitif — \textit{helfen} admet les deux constructions : \cle{helfen + Inf.} ou \cle{helfen, zu + Inf.} ; ici structure nominale \cle{Fehler zu korrigieren} comme objet de \cle{helfen} ou but implicite).
        \item \all{Um zu verstehen} (But en position 1 $\rightarrow$ inversion : \all{muss man zuhören}).
        \item \all{ohne zu grüßen} (\cle{ohne ... zu}).
    \end{enumerate}
\end{enumerate}
\end{corrige}

\courssub{Analyse linguistique : Les connecteurs logiques (Konnektoren)}

Le texte est cohérent grâce à des connecteurs (Konjunktionen, Adverbes conjonctifs) qui structurent l'argumentation. Les identifier permet de suivre la logique de l'auteur et de les réutiliser en expression écrite.

\begin{propriete}[Les connecteurs du texte : Forme, Sens, Place du verbe]
\begin{center}
\begin{tabularx}{\linewidth}{|X|X|X|X|}
\hline
\rowcolor{popblueL}
\textbf{Connecteur (allemand)} & \textbf{Fonction / Sens (français)} & \textbf{Structure syntaxique} & \textbf{Exemple du texte} \\
\hline
\cle{aber} & Opposition, concession (mais) & Conjonction de coordination $\rightarrow$ \textbf{Position 0} (verbe en 2) & „Die KI ist ein starkes Werkzeug, \textbf{aber} sie ersetzt nicht den Pädagogen.“ \\
\hline
\cle{also} & Conséquence, conclusion (donc, ainsi) & Adverbe conjonctif $\rightarrow$ \textbf{Position 1} (verbe en 2, inversion) & „\textbf{Also} \textbf{sollten} Schulen auf ein hybrides Modell setzen.“ \\
\hline
\cle{deshalb} & Conséquence forte (c'est pourquoi) & Adverbe conjonctif $\rightarrow$ \textbf{Position 1} (inversion) & „\textbf{Deshalb} \textbf{warnen} Experten vor einer reinen Bildschirm-Schule.“ \\
\hline
\cle{so} & Conséquence / Résultat (ainsi, du coup) & Adverbe conjonctif $\rightarrow$ \textbf{Position 1} (inversion) & „\textbf{So} \textbf{kann} jeder in seinem eigenen Tempo lernen.“ \\
\hline
\cle{um ... zu + Infinitiv} & But final (pour + infinitif) & Groupe infinitif \textbf{en fin de phrase} (ou Pos 1) ; sujet identique à la principale & „\textbf{Um} gut \textbf{zu lernen}, \textbf{braucht} der Mensch echte Begegnung.“ \\
\hline
\end{tabularx}
\end{center}
\end{propriete}

\begin{attention}[Erreurs fréquentes des francophones sur les connecteurs]
\begin{itemize}
    \item \textbf{\cle{aber} vs \cle{sondern}} : \cle{aber} = « mais » (après affirmation ou négation). \cle{sondern} = « mais » (uniquement après une négation pour corriger). Ex. : \all{Nicht die KI entscheidet, \textbf{sondern} der Lehrer.}
    \item \textbf{\cle{also} / \cle{deshalb} / \cle{so} + inversion} : En français « Donc, je vais... » (sujet-verbe). En allemand : \all{Also \textbf{gehe} ich...} (Verbe-Sujet). \textbf{Jamais} : \all{Also ich gehe...} (sauf langage oral très familier).
    \item \textbf{\cle{um ... zu}} : Le verbe conjugué de la principale est souvent au début (position 1) si le groupe \cle{um...zu} est en position 1 (Thème). Ex. : \all{Um Deutsch zu lernen, \textbf{gehe} ich in den Kurs.} Pas de « für ... zu » ! (\cle{für} + nom = pour [nom] ; \cle{um ... zu} = pour [verbe]).
\end{itemize}
\end{attention}

\courssub{Recherche lexicale : Champs lexicaux (Wortfelder)}

Travailler par champs lexicaux (Wortfelder) est la meilleure méthode pour mémoriser le vocabulaire thématique et réussir l'expression écrite. Relevez dans le texte les mots appartenant aux trois catégories ci-dessous.

\begin{exoresolu}[Champs lexicaux du texte — Tableau récapitulatif]
\begin{center}
\begin{tabularx}{\linewidth}{|X|X|X|}
\hline
\rowcolor{popgreenL}
\textbf{Appareils \& Outils (Geräte \& Werkzeuge)} & \textbf{Verbes d'apprentissage \& action (Lernverben)} & \textbf{Opinions \& Jugements (Meinungen \& Urteile)} \\
\hline
\all{das Tablet, die Tablets} & \all{analysieren} (analyser) & \all{stark} (fort, puissant — adj.) \\
\all{der Laptop, die Laptops} & \all{vorschlagen} (proposer) & \all{warnt vor} (met en garde contre) \\
\all{das Lernprogramm, die -e} & \all{lernen} (apprendre) & \all{ersetzen} (remplacer) \\
\all{die App, die Apps} & \all{erkennen} (reconnaître) & \all{plädieren für} (plaider pour) * \\
\all{die künstliche Intelligenz (KI)} & \all{motivieren} (motiver) & \all{profitieren von} (profiter de) \\
\all{der Algorithmus, die -men} & \all{setzen} (mettre en place / opter pour) & \all{Chance / Risiko} (chance / risque) \\
\all{das Werkzeug, die -e} & \all{brauchen} (avoir besoin de) & \all{Vor- / Nachteil} (avantage / inconvénient) \\
\hline
\end{tabularx}
\end{center}
\emph{* \all{plädieren für} n'est pas dans le texte mais est le verbe exact pour « plaider pour / recommander » implicite dans „Also sollten Schulen... setzen“.}
\end{exoresolu}

\begin{exercice}[Entraînement lexical — Complétez en allemand (conjugaison correcte)]
\begin{enumerate}
    \item Le programme \_\_\_\_\_\_\_\_ (analyse) les erreurs.
    \item Il \_\_\_\_\_\_\_\_ (propose) des exercices adaptés.
    \item L'enseignant \_\_\_\_\_\_\_\_ (reconnaît) les émotions.
    \item L'algorithme ne peut pas \_\_\_\_\_\_\_\_ (motiver).
    \item Les experts \_\_\_\_\_\_\_\_ (mettent en garde) contre l'école-écran.
    \item L'école doit \_\_\_\_\_\_\_\_ (opter pour / mettre en place) un modèle hybride.
\end{enumerate}
\end{exercice}

\begin{corrige}[Exercice Entraînement lexical]
\begin{enumerate}
    \item \all{analysiert} (Präsens, 3e pers. sg.)
    \item \all{schlägt vor} (verbe à particule séparable : \cle{vorschlagen} $\rightarrow$ \cle{schlägt ... vor})
    \item \all{erkennt} (changement de voyelle \cle{e $\rightarrow$ ie} à la 2e/3e pers. sg.)
    \item \all{motivieren} (infinitif après modal \cle{kann})
    \item \all{warnt vor} (Präsens 3e pl. + préposition \cle{vor} + Dativ)
    \item \all{setzen} (infinitif après modal \cle{sollten}) — \all{auf ein hybrides Modell setzen}
\end{enumerate}
\end{corrige}

\courssub{Synthèse et production : Le résumé (Zusammenfassung)}

Le résumé (souvent 3-4 phrases, max 5 lignes) est un exercice d'examen fréquent. Il ne s'agit pas de « couper-coller » mais de \textbf{restructurer l'information} en suivant la logique argumentative du texte.

\begin{methode}[Rédiger un résumé en 4 phrases types (Modèle Bac)]
Suivez ce canevas rigoureux. Utilisez le \textbf{Präsens} (texte d'actualité) ou le \textbf{Präteritum} (récit). Reliez les phrases par des connecteurs (\cle{aber}, \cle{also}, \cle{weil}, \cle{indem}).

\begin{enumerate}
    \item \textbf{Phrase d'accroche (Thème + Contexte) :} \all{Der Text „Lernen mit der KI“ thematisiert den Einsatz digitaler Lernwerkzeuge an kamerunischen Schulen.}
    \item \textbf{Phrase sur les avantages (Pro-Argument) :} \all{KI-Programme ermöglichen individuelles Lernen im eigenen Tempo, indem sie Fehler analysieren und passende Übungen vorschlagen.}
    \item \textbf{Phrase sur les limites / contre-argument (Contra) :} \all{Experten warnen jedoch davor, dass reine Bildschirm-Schule den sozialen Kontakt und die emotionale Motivation durch Lehrer gefährdet.}
    \item \textbf{Phrase de conclusion (Position de l'auteur / Solution) :} \all{Daher plädiert der Autor für ein hybrides Modell, das digitale Tools und pädagogische Begleitung verbindet.}
\end{enumerate}
\end{methode}

\begin{exemplebox}[Résumé modèle (Version finale pour copie)]
\all{Der Text „Lernen mit der KI“ thematisiert den Einsatz digitaler Lernwerkzeuge an kamerunischen Schulen. KI-Programme ermöglichen individuelles Lernen im eigenen Tempo, indem sie Fehler analysieren und passende Übungen vorschlagen. Experten warnen jedoch davor, dass reine Bildschirm-Schule den sozialen Kontakt und die emotionale Motivation durch Lehrer gefährdet. Daher plädiert der Autor für ein hybrides Modell, das digitale Tools und pädagogische Begleitung verbindet.}
\end{exemplebox}

\begin{attention}[Critères de notation du résumé (Grille Bac type)]
\begin{itemize}
    \item \textbf{Contenu (6 pts) :} Les 4 idées forces présentes ? (Contexte, Avantages, Risques, Solution hybride).
    \item \textbf{Langue (4 pts) :} Allemand correct (cas, ordre des mots, orthographe, connecteurs). Phrases complètes (pas de style télégraphique).
    \item \textbf{Forme (2 pts) :} Longueur respectée (4-5 lignes), pas de citation brute, rédaction soignée.
    \item \textbf{Piège :} Ne pas donner \emph{votre} avis (« Ich finde... »). Résumer = rendre compte de l'avis \emph{de l'auteur}.
\end{itemize}
\end{attention}

\courssub{Illustration : Avantages et Inconvénients de la KI (à compléter)}

La figure ci-dessous synthétise visuellement l'argumentation du texte. Complétez les cases vides en vous basant \textbf{strictement} sur le texte (et votre connaissance du thème).

\begin{popfigure}
\begin{tikzpicture}[
    node distance=0.8cm and 1.2cm,
    box/.style={draw=popblue, thick, rounded corners, minimum height=1.2cm, minimum width=5.5cm, align=center, fill=popblueL, font=\small},
    title/.style={draw=popdark, thick, rounded corners, minimum height=0.9cm, minimum width=5.5cm, align=center, fill=popblue, font=\bfseries\color{white}\small},
    arrow/.style={-Stealth, thick, poporange}
]
% Titres
\node[title] (t1) at (0,0) {Vorteile der KI (Avantages)};
\node[title] (t2) at (8,0) {Nachteile der KI (Inconvénients)};

% Avantages (colonne gauche)
\node[box, below=of t1, align=center] (a1){Individuelles Lerntempo\\„jeder in seinem eigenen Tempo“};
\node[box, below=of a1, align=center] (a2){Sofortige Fehleranalyse\\„analysieren die Fehler“};
\node[box, below=of a2, align=center] (a3){Übungen passend zum Niveau\\„schlagen passende Übungen vor“};
\node[box, below=of a3, align=center] (a4){Verfügbarkeit überall / flexibel\\(Tablets, Apps)};

% Inconvénients (colonne droite)
\node[box, below=of t2, align=center] (n1){Verlust von sozialem Kontakt\\„verliert oft den Kontakt“};
\node[box, below=of n1, align=center] (n2){Fehlende emotionale Unterstützung\\„Algorithmus kann [Emotionen] nicht erkennen“};
\node[box, below=of n2, align=center] (n3){Keine echte Motivation\\„kann ... nicht motivieren“};
\node[box, below=of n3, align=center] (n4){Risiko: Passivität / Kritikfähigkeit\\(implizit: „nur am Bildschirm“)};

% Flèches de synthèse vers le bas
\node[title, below=1.5cm of a4, xshift=4cm, fill=poppurple, minimum width=11.5cm] (synth) {Synthèse du texte : L'hybridation comme solution (Hybrides Modell)};
\draw[arrow] (a4.south) -- ++(0,-0.5) -| (synth.north west);
\draw[arrow] (n4.south) -- ++(0,-0.5) -| (synth.north east);
\end{tikzpicture}
\legende{Tableau comparatif extrait du texte : à apprendre par cœur pour l'expression orale/écrite.}
\end{popfigure}

\begin{experience}[Activité orale guidée : Debate « KI im Klassenzimmer »]
\textbf{Consigne :} En binôme. L'un joue le rôle d'un \textbf{proviseur favorable à la tablette pour tous} (Rolle: Schulleiter/in), l'autre celui d'un \textbf{professeur sceptique} (Rolle: kritische Lehrkraft).
\textbf{Durée :} 3 min préparation (notes en allemand), 4 min débat.
\textbf{Contraintes linguistiques :} Utiliser au moins 3 connecteurs (\cle{aber}, \cle{deshalb}, \cle{um...zu}, \cle{obwohl}) + 5 mots du champ lexical (Appareils, Verbes, Opinions).
\textbf{Exemple d'amorce :} \all{„Als Schulleiterin finde ich, dass Tablets die Chancengerechtigkeit erhöhen, weil...“} / \all{„Aber als Lehrer sehe ich die Gefahr, dass...“}
\end{experience}

\begin{savaistu}[Le Cameroun et la GermanTech : Le projet « Digital Schools »]
Le Cameroun collabore avec la coopération allemande (GIZ) et des entreprises comme \textbf{SAP} ou \textbf{Siemens} pour équiper des lycées en laboratoires numériques. Le programme \emph{eLearning Africa} (souvent accueilli à Yaoundé) présente des startups germanophones (ex. \emph{Simpleclub}, \emph{StudySmarter}) qui adaptent leurs apps aux programmes africains. La maîtrise de l'allemand ouvre donc des portes directes vers ces métiers de l'\emph{EdTech} (Technologie éducative) !
\end{savaistu}

\coursec{Lexique thématique : digitale Medien und KI}

Après avoir travaillé la compréhension globale et détaillée du texte \emph{Lernen mit der KI}, nous allons maintenant fixer et structurer le vocabulaire essentiel du chapitre. Cette étape est indispensable pour pouvoir ensuite manipuler la grammaire (\cle{Futur I}, \cle{Infinitivgruppen}) et produire à l'oral comme à l'écrit. Nous procéderons par champs lexicaux, en observant d'abord les mots dans leur contexte, puis en mémorisant systématiquement \textbf{article, pluriel, genre} et \textbf{traduction}. Chaque sous-section se termine par des exemples authentiques et des pièges à éviter pour les francophones.

\courssub{Champ lexical 1 : Les appareils (die Geräte)}

Commençons par le matériel physique. Dans le texte support, vous avez rencontré \all{der Computer}, \all{das Tablet} et \all{das Handy}. Observons leurs particularités morphologiques.

\begin{exemplebox}[Observation dans le texte]
\all{„Die Schüler arbeiten am Computer.“} (« Les élèves travaillent sur l'ordinateur. ») \\
\all{„Jeder hat ein Tablet.“} (« Chacun a une tablette. ») \\
\all{„Das Handy liegt auf dem Tisch.“} (« Le smartphone est sur la table. »)
\end{exemplebox}

\begin{propriete}[Règle : Genre et pluriel des noms d'appareils d'origine anglaise]
Beaucoup de termes techniques récents sont des emprunts à l'anglais. En allemand, ils sont souvent neutres (\textbf{das}) et forment leur pluriel en \textbf{-s} (sans Umlaut), exactement comme en anglais.
\begin{itemize}
    \item \textbf{das Tablet, die Tablets} (la tablette)
    \item \textbf{das Handy, die Handys} (le smartphone / portable) — \emph{Attention : en allemand, « Handy » est un nom, pas un adjectif !}
    \item \textbf{der Computer, die Computer} (l'ordinateur) — \emph{Pluriel inchangé (comme der Lehrer / die Lehrer). Ne jamais dire „die Computers“ !}
\end{itemize}
Les mots d'origine germanique suivent les règles habituelles (Umlaut + -e, -e, -n, etc.) :
\begin{itemize}
    \item \textbf{der Bildschirm, die Bildschirme} (l'écran)
    \item \textbf{die Tastatur, die Tastaturen} (le clavier)
    \item \textbf{der Drucker, die Drucker} (l'imprimante — pluriel inchangé)
\end{itemize}
\end{propriete}

\begin{center}
\begin{tabularx}{\linewidth}{|X|X|X|}
\hline
\rowcolor{popblueL} \textbf{Deutsch (Artikel + Plural)} & \textbf{Français} & \textbf{Remarque / Exemple} \\
\hline
der Computer (-) & l'ordinateur & Pluriel identique. \all{Die Computer sind neu.} \\
das Tablet (-s) & la tablette & Emprunt anglais, pluriel en -s. \all{Wir haben zwei Tablets.} \\
das Handy (-s) & le smartphone & \emph{Faux ami :} ne pas confondre avec l'adj. français « handy ». \all{Mein Handy ist kaputt.} \\
der Bildschirm (-e) & l'écran & Umlaut au pluriel. \all{Der Bildschirm ist groß.} \\
die Tastatur (-en) & le clavier & Féminin. \all{Die Tastatur funktioniert nicht.} \\
der Drucker (-) & l'imprimante & Masculin, pluriel inchangé. \all{Der Drucker druckt schnell.} \\
\hline
\end{tabularx}
\end{center}

\begin{attention}[Piège fréquent : « die Computer » et « das Handy »]
\begin{itemize}
    \item \textbf{Ne jamais écrire « die Computers ».} Le mot \all{der Computer} fait partie des noms en \textbf{-er} dont le pluriel est inchangé (comme \all{der Lehrer}, \all{der Schüler}, \all{der Vater}).
    \item \textbf{« Handy » est un nom neutre :} \textbf{das Handy}. En français, « un handy » n'existe pas ; on dit « un portable » ou « un smartphone ». Ne traduisez pas « handy » par « pratique » (c'est l'adjectif anglais) !
    \item \textbf{Préposition :} On travaille \textbf{am} (= an dem) Computer / Tablet / Bildschirm (lieu/d'appareil), mais \textbf{mit dem} Handy (instrument).
\end{itemize}
\end{attention}

\begin{exoresolu}[Entraînement : Articles et pluriels]
\textbf{Énoncé :} Mettez les noms suivants au pluriel avec l'article défini, puis traduisez en français.
\begin{enumerate}
    \item der Bildschirm
    \item das Tablet
    \item die Tastatur
    \item der Drucker
    \item das Handy
    \item der Computer
\end{enumerate}
\tcblower
\textbf{Solution détaillée :}
\begin{enumerate}
    \item \textbf{die Bildschirme} — les écrans (Umlaut + -e).
    \item \textbf{die Tablets} — les tablettes (pluriel anglais en -s).
    \item \textbf{die Tastaturen} — les claviers (féminin en -ur $\rightarrow$ -en).
    \item \textbf{die Drucker} — les imprimantes (pluriel inchangé, masculin en -er).
    \item \textbf{die Handys} — les smartphones (pluriel anglais en -s, \emph{y} $\rightarrow$ \emph{ys}).
    \item \textbf{die Computer} — les ordinateurs (pluriel inchangé).
\end{enumerate}
\end{exoresolu}

\courssub{Champ lexical 2 : Les outils numériques et l'intelligence artificielle}

Passons maintenant aux logiciels, plateformes et au concept central du chapitre : l'intelligence artificielle.

\begin{textebox}[Mini-texte d'observation : Die digitale Werkzeugkiste]
\all{In der modernen Schule benutzen wir viele digitale Werkzeuge. Die wichtigste Software ist das Lernprogramm. Es hilft beim Üben. Die künstliche Intelligenz (KI) analysiert die Fehler der Schüler und passt die Aufgaben an. Man lädt die App aus dem Internet herunter. Die Plattform verbindet Lehrer und Schüler. Man kann Dateien speichern und teilen.}
\end{textebox}

\begin{definition}[\cle{das Lernprogramm}, die Lernprogramme]
Un \textbf{Lernprogramm} est un logiciel éducatif (logiciel d'apprentissage) qui propose des exercices interactifs, corrige automatiquement les réponses de l'élève et adapte souvent le niveau de difficulté (das Niveau) en fonction des résultats. C'est un mot composé typique : \all{Lernen} (l'apprentissage) + \all{Programm} (le programme).
\end{definition}

\begin{propriete}[Vocabulaire : Outils, IA et Internet]
\begin{itemize}
    \item \textbf{die App, die Apps} (l'application) — féminin, pluriel en -s (emprunt anglais court pour \emph{Application}).
    \item \textbf{das Lernprogramm, die Lernprogramme} (le logiciel éducatif / le programme d'apprentissage).
    \item \textbf{die Software} (le logiciel / les logiciels) — \emph{Singulier seulement (Singularetantum), féminin}. Pas de pluriel \*Softwares.
    \item \textbf{das Internet} (internet) — neutre, pas de pluriel, souvent sans article : \all{im Internet} (sur internet).
    \item \textbf{die Plattform, die Plattformen} (la plateforme).
    \item \textbf{die künstliche Intelligenz (KI)} (l'intelligence artificielle / IA) — féminin. Abréviation courante : \all{die KI} (féminin aussi : \all{die KI hilft}).
    \item \textbf{der Algorithmus, die Algorithmen} (l'algorithme) — utile pour comprendre le fonctionnement de la KI.
\end{itemize}
\end{propriete}

\begin{center}
\begin{tabularx}{\linewidth}{|X|X|X|}
\hline
\rowcolor{poppurpleL} \textbf{Deutsch (Artikel + Plural)} & \textbf{Français} & \textbf{Contexte / Collocation} \\
\hline
die Apps & l'application & \all{eine App herunterladen / benutzen} \\
das Lernprogramm, die Lernprogramme & le logiciel éducatif & \all{mit dem Lernprogramm üben} \\
die Software (-) & le(s) logiciel(s) & \all{neue Software installieren} (pas de pluriel) \\
das Internet (-) & internet & \all{im Internet surfen / recherchieren} \\
die Plattform, die Plattformen & la plateforme & \all{auf der Plattform arbeiten} \\
die künstliche Intelligenz / die KI & l'intelligence artificielle (IA) & \all{die KI korrigiert Fehler} \\
der Algorithmus, die Algorithmen & l'algorithme & \all{der Algorithmus lernt} \\
\hline
\end{tabularx}
\end{center}

\begin{savaistu}[L'allemand, champion des mots composés (Komposita)]
L'allemand adore coller des mots existants pour en créer de nouveaux. C'est une \textbf{stratégie majeure pour le Bac} : si vous savez découper le mot, vous devinez son sens sans dictionnaire !
\begin{itemize}
    \item \textbf{künstlich} (artificiel) + \textbf{Intelligenz} (intelligence) = \textbf{künstliche Intelligenz} (IA).
    \item \textbf{Lernen} (apprentissage / infinitif substantivé) + \textbf{Programm} = \textbf{Lernprogramm} (logiciel éducatif).
    \item \textbf{Fehler} (erreurs) + \textbf{Korrektur} (correction) = \textbf{Fehlerkorrektur} (correction d'erreurs).
    \item \textbf{Bild} (image) + \textbf{Schirm} (écran/protection) = \textbf{Bildschirm} (écran d'ordinateur).
\end{itemize}
\textbf{Règle d'or :} Le genre et le pluriel du composé sont ceux du \textbf{dernier élément} (le noyau). Ex: \emph{das Programm} $\rightarrow$ \emph{das Lernprogramm, die Lernprogramme}.
\end{savaistu}

\begin{savaistu}[Le Cameroun et la transition numérique éducative]
Au Cameroun, le projet \textbf{« École Numérique »} (appuyé par le MINESEC et des partenaires comme l'UNESCO ou la GIZ) équipe progressivement les lycées en \all{Tablets} et \all{Lernplattformen}. Des applications comme \all{« MINESEC E-Learning »} ou des plateformes type \all{Moodle} permettent aux élèves de Yaoundé, Douala ou Bertoua d'accéder à des \all{Lernprogramme} hors connexion. Le vocabulaire de cette leçon est donc directement applicable à votre réalité quotidienne !
\end{savaistu}

\begin{exemplebox}[Exemples traduits : Outils et IA]
\all{„Die App passt das Niveau an.“} — « L'application adapte le niveau. » \\
\all{„Wir lernen mit künstlicher Intelligenz.“} — « Nous apprenons avec de l'intelligence artificielle. » \\
\all{„Der Lehrer lädt das Programm herunter.“} — « Le professeur télécharge le programme. » \\
\all{„Die Plattform verbindet alle Teilnehmer.“} — « La plateforme relie tous les participants. »
\end{exemplebox}

\courssub{Champ lexical 3 : Les verbes d'action (Tätigkeitsverben)}

Ces verbes sont le moteur de vos phrases. Certains sont à particule séparable (\all{herunterladen}, \all{anpassen}), d'autres sont simples. Maîtrisez leurs régimes (cas) et leurs formes du \textbf{Perfekt} (auxiliaire \all{haben} sauf verbes de déplacement/état).

\begin{propriete}[Conjugaison et régime des verbes clés]
\begin{center}
\begin{tabular}{|l|l|l|l|l|}
\hline
\rowcolor{popgreenL} \textbf{Infinitif} & \textbf{Präsens (ich/er/sie/es)} & \textbf{Perfekt (ich habe ...)} & \textbf{Régime / Cas} & \textbf{Français} \\
\hline
lernen & lernt & gelernt & Akk (etwas lernen) & apprendre \\
üben & übt & geübt & Akk (etwas üben) & s'exercer, réviser \\
korrigieren & korrigiert & korrigiert & Akk (etwas korrigieren) & corriger \\
anpassen & passt an & angepasst & Akk (etwas anpassen) & adapter \\
benutzen & benutzt & benutzt & Akk (etwas benutzen) & utiliser \\
\hline
herunterladen & lädt herunter & heruntergeladen & Akk (etwas herunterladen) & télécharger \\
speichern & speichert & gespeichert & Akk (etwas speichern) & sauvegarder \\
teilen & teilt & geteilt & Akk (etwas teilen) & partager \\
ersetzen & ersetzt & ersetzt & Akk (etwas ersetzen) & remplacer \\
verbinden & verbindet & verbunden & Akk (etwas mit etwas verbinden) & relier, connecter \\
\hline
\end{tabular}
\end{center}
\emph{Note :} \all{herunterladen} et \all{anpassen} sont à particule séparable $\rightarrow$ \all{ich lade herunter}, \all{er hat heruntergeladen} (participe : \textbf{ge-hauptwort-en} $\rightarrow$ heruntergeladen) ; \all{er passt an}, \all{er hat angepasst}.
\end{propriete}

\begin{attention}[Piège majeur : \textbf{lernen} vs \textbf{studieren}]
C'est l'erreur n°1 des francophones !
\begin{itemize}
    \item \textbf{lernen} = \textbf{apprendre} (acquérir un savoir, une langue, une leçon), à tout âge, n'importe où. \all{Ich lerne Deutsch.} (J'apprends l'allemand.) \all{Er lernt für die Prüfung.} (Il révise pour l'examen.)
    \item \textbf{studieren} = \textbf{étudier à l'université} (être inscrit comme étudiant, suivre un cursus supérieur). \all{Sie studiert Medizin in Berlin.} (Elle étudie la médecine à Berlin / Elle est étudiante en médecine.)
    \item \textbf{Au lycée, on dit toujours :} \all{Deutsch lernen}, \all{Englisch lernen}, \all{Mathe lernen}. \textbf{JAMAIS} \all{Deutsch studieren} (sauf si vous êtes déjà à la fac !).
\end{itemize}
\end{attention}

\begin{exemplebox}[Verbes en contexte (phrases types)]
\all{„Die Schüler \textbf{lernen} mit der App.“} — Les élèves \textbf{apprennent} avec l'appli. \\
\all{„Der Computer \textbf{korrigiert} die Fehler.“} — L'ordinateur \textbf{corrige} les erreurs. \\
\all{„Ich \textbf{lade} die Datei \textbf{herunter}.“} — Je \textbf{télécharge} le fichier. \\
\all{„Das Programm \textbf{passt} die Aufgaben \textbf{an}.“} — Le programme \textbf{adapte} les exercices. \\
\all{„Wir \textbf{teilen} die Ergebnisse.“} — Nous \textbf{partageons} les résultats.
\end{exemplebox}

\begin{exoresolu}[Entraînement : Conjugaison et choix du verbe]
\textbf{Énoncé :} Conjuguez le verbe entre parenthèses au Präsens (3e pers. sg. ou pl.) ou au Perfekt. Choisissez le bon verbe pour la phrase 5.
\begin{enumerate}
    \item Die KI \_\_\_\_\_\_\_\_\_\_ (korrigieren) die Fehler sofort.
    \item Wir \_\_\_\_\_\_\_\_\_\_ (herunterladen) die neue App gestern.
    \item Der Lehrer \_\_\_\_\_\_\_\_\_\_ (anpassen) das Niveau für jeden Schüler.
    \item Ich habe die Hausaufgaben \_\_\_\_\_\_\_\_\_\_ (speichern).
    \item Mein Bruder \_\_\_\_\_\_\_\_\_\_ (lernen / studieren) Informatik an der Uni. (Choisir et conjuguer)
\end{enumerate}
\tcblower
\textbf{Solution :}
\begin{enumerate}
    \item \textbf{korrigiert} (Präsens, 3e sg., verbe régulier).
    \item \textbf{haben heruntergeladen} (Perfekt, particule séparable $\rightarrow$ ge- au milieu : herunter-geladen). \emph{Attention :} « Wir haben die App heruntergeladen. »
    \item \textbf{passt ... an} (Präsens, particule séparable : \all{er passt das Niveau an}).
    \item \textbf{gespeichert} (Perfekt, participe passé régulier ge-t-en).
    \item \textbf{studiert} — Car c'est à l'université (\all{an der Uni}) $\rightarrow$ \all{studieren}. \all{Mein Bruder studiert Informatik.}
\end{enumerate}
\end{exoresolu}

\courssub{Champ lexical 4 : Personnes et contexte scolaire (Personen und Schule)}

Pour parler de qui fait quoi, il faut maîtriser les paires masculin/féminin (Bewegungsformen) et les noms d'institution.

\begin{propriete}[Noms d'agents : Dérivation en \textbf{-er} / \textbf{-in}]
Le suffixe \textbf{-er} (masculin) / \textbf{-in} (féminin, pluriel \textbf{-nen}) forme le nom d'agent à partir du verbe ou du nom de fonction.
\begin{center}
\begin{tabularx}{\linewidth}{|X|X|X|X|}
\hline
\rowcolor{poporangeL} \textbf{Maskulin (der)} & \textbf{Féminin (die)} & \textbf{Pluriel (die)} & \textbf{Français} \\
\hline
der Schüler & die Schülerin & die Schüler / die Schülerinnen & l'élève \\
der Lehrer & die Lehrerin & die Lehrer / die Lehrerinnen & le/la prof \\
der Experte & die Expertin & die Experten / die Expertinnen & l'expert(e) \\
\hline
\end{tabularx}
\end{center}
\emph{Autres noms du contexte :}
\begin{itemize}
    \item \textbf{der Unterricht (-e)} — le cours, l'enseignement (ex: \all{im Unterricht} = en classe / pendant le cours).
    \item \textbf{die Aufgabe (-n)} — la tâche, l'exercice, le devoir.
    \item \textbf{das Niveau (-s)} — le niveau (neutre, pluriel \textbf{-s}).
    \item \textbf{die Klasse (-n)} — la classe (groupe d'élèves).
    \item \textbf{die Schule (-n)} — l'école (établissement).
\end{itemize}
\end{propriete}

\begin{exemplebox}[Personnes et école en phrases]
\all{„Die Schülerin stellt eine Frage.“} — L'élève (fille) pose une question. \\
\all{„Der Experte erklärt den Algorithmus.“} — L'expert explique l'algorithme. \\
\all{„Im Unterricht benutzen wir Tablets.“} — En classe, nous utilisons des tablettes. \\
\all{„Die Aufgabe ist schwierig.“} — L'exercice est difficile.
\end{exemplebox}

\courssub{Familles de mots et dérivation (Wortfamilien)}

Travailler par familles permet de multiplier votre vocabulaire sans effort. Observez la racine commune.

\begin{center}
\begin{tabularx}{\linewidth}{|X|X|X|}
\hline
\rowcolor{poppinkL} \textbf{Racine / Verbe} & \textbf{Nom (der/die/das + Plural)} & \textbf{Adjectif / Autre} \\
\hline
\textbf{lernen} (apprendre) & das Lernen (l'apprentissage, infinitif subst.) & lernend (apprenant) \\
& der Lerner / die Lernerin (l'apprenant) & \\
& \textbf{das Lernprogramm} (logiciel éducatif) & \\
\hline
\textbf{korrigieren} (corriger) & \textbf{die Korrektur (-en)} (la correction) & korrekt (correct) \\
& der Korrektor / die Korrektorin (le correcteur) & \\
\hline
\textbf{intelligent} (adj.) & \textbf{die Intelligenz} (l'intelligence) & intelligent \\
& \textbf{die künstliche Intelligenz (KI)} & künstlich (artificiel) \\
\hline
\textbf{benutzen} (utiliser) & die Benutzung (l'utilisation) & benutzbar (utilisable) \\
& der Benutzer / die Benutzerin (l'utilisateur) & \\
\hline
\end{tabularx}
\end{center}

\begin{methode}[Comment enrichir son vocabulaire par les familles ?]
\begin{enumerate}
    \item Repérez le \textbf{mot de base} (verbe, nom ou adjectif) dans le texte.
    \item Cherchez les \textbf{dérivés} : noms en \textbf{-ung, -tion, -heit, -keit, -er/-in}, adjectifs en \textbf{-bar, -lich, -los}.
    \item Notez le \textbf{genre} et le \textbf{pluriel} du nom dérivé.
    \item Écrivez une \textbf{phrase personnelle} avec chaque nouveau mot.
\end{enumerate}
\end{methode}

\courssub{Expressions utiles (Redemittel) pour l'oral et l'écrit}

Ces blocs figés (chunks) vous permettront de parler fluidement du numérique en classe.

\begin{center}
\begin{tabularx}{\linewidth}{|X|X|}
\hline
\rowcolor{popgoldL} \textbf{Deutsch} & \textbf{Français} \\
\hline
\all{am Bildschirm arbeiten} & travailler sur écran / à l'ordinateur \\
\all{im Internet surfen / recherchieren} & naviguer / faire des recherches sur internet \\
\all{eine App benutzen / herunterladen} & utiliser / télécharger une appli \\
\all{Fehler korrigieren / verbessern} & corriger / améliorer des erreurs \\
\all{den Unterricht verändern / bereichern} & transformer / enrichir le cours \\
\all{das Niveau anpassen} & adapter le niveau \\
\all{mit künstlicher Intelligenz lernen} & apprendre avec l'IA \\
\all{Dateien speichern und teilen} & sauvegarder et partager des fichiers \\
\all{eine Plattform nutzen} & utiliser une plateforme \\
\hline
\end{tabularx}
\end{center}

\begin{attention}[Pièges de collocation (Verben + Prépositions / Cas)]
\begin{itemize}
    \item \textbf{arbeiten \textbf{an} + Dativ} : \all{arbeiten am (an dem) Computer / Tablet / Bildschirm}. Pas *arbeiten mit dem Computer* (possible mais moins précis pour « travailler \emph{sur} l'appareil »).
    \item \textbf{surfen \textbf{im} (= in dem) Internet} (Dativ locatif). Pas *surfen auf dem Internet*.
    \item \textbf{teilnehmen \textbf{an} + Dativ} : \all{am Unterricht teilnehmen} (participer au cours).
    \item \textbf{helfen \textbf{bei} + Dativ} : \all{Die KI hilft beim (bei dem) Lernen.} (L'IA aide pour l'apprentissage).
    \item \textbf{ersetzen \textbf{durch} + Akkusativ} : \all{Der Lehrer wird nicht durch die KI ersetzt.}
\end{itemize}
\end{attention}

\courssub{Texte d'entraînement : « Der digitale Klassenraum »}

Le texte suivant réemploie tout le vocabulaire de la section. Lisez-le à haute voix, soulignez les mots connus, déduisez les inconnus par la composition, puis répondez aux questions.

\begin{textebox}[Der digitale Klassenraum — Texte original (env. 230 mots)]
\all{In vielen Schulen in Deutschland, Österreich und der Schweiz hat sich der Unterricht in den letzten Jahren stark verändert. Früher gab es nur Tafel, Kreide und Bücher. Heute stehen in den Klassenzimmern moderne Geräte: der Computer, das Tablet und manchmal auch das Handy der Schüler. Die Geräte sind mit dem Internet verbunden. 

Die Lehrer benutzen spezielle Lernprogramme und Apps. Ein bekanntes Programm analysiert die Antworten der Schüler. Die künstliche Intelligenz (KI) erkennt sofort, wo die Schwierigkeiten liegen. Dann passt das Programm die Aufgaben automatisch an das Niveau jedes Kindes an. So kann jeder in seinem eigenen Tempo lernen und üben. 

Die Schüler laden die Apps aus dem Internet herunter. Sie speichern ihre Ergebnisse auf der Lernplattform. Der Lehrer sieht den Fortschritt auf seinem Bildschirm. Er kann die Fehler korrigieren und neue Aufgaben teilen. Die digitale Plattform verbindet die Klasse auch zu Hause. 

Natürlich ersetzt die Technik nicht den Lehrer. Der Experte entscheidet, welche Software sinnvoll ist. Er erklärt die schwierigen Themen im Unterricht. Die Geräte sind nur Werkzeuge. Aber sie machen das Lernen oft motivierender und effektiver.}
\end{textebox}

\begin{definition}[Glossaire des mots difficiles du texte]
\begin{itemize}
    \item \textbf{sich verändern} (verändert, hat sich verändert) — changer, évoluer (réfléchi).
    \item \textbf{früher} — autrefois, avant.
    \item \textbf{die Tafel (-n)} — le tableau (noir/blanc).
    \item \textbf{die Kreide} — la craie (féminin, pas de pluriel usuel).
    \item \textbf{stehen (steht, stand, hat gestanden)} — être placé, se trouver (emplacement).
    \item \textbf{verbinden (verbindet, hat verbunden)} — relier, connecter.
    \item \textbf{spezielle} — spécial(e)(s) (adj. décliné).
    \item \textbf{analysieren} — analyser.
    \item \textbf{erkennen (erkennt, erkannte, hat erkannt)} — reconnaître, identifier.
    \item \textbf{automatisch} — automatiquement.
    \item \textbf{der Fortschritt (-e)} — le progrès.
    \item \textbf{sinnvoll} — utile, judicieux, pertinent.
    \item \textbf{das Werkzeug (-e)} — l'outil.
    \item \textbf{motivierend} — motivant.
    \item \textbf{effektiv} — efficace.
\end{itemize}
\end{definition}

\begin{exercice}[Questions de compréhension et de lexique sur le texte]
\begin{enumerate}
    \item Citez trois appareils (\textbf{Geräte}) mentionnés dans le premier paragraphe avec leurs articles.
    \item Que fait la \textbf{KI} avec les réponses des élèves ? (Verbe : \all{analysieren / erkennen / anpassen})
    \item Quel est l'avantage pour l'élève : « in seinem eigenen Tempo ... » ? Traduisez.
    \item Où les élèves sauvegardent-ils leurs résultats ? (Nom composé : \textbf{Lernplattform}).
    \item Quel rôle le texte attribue-t-il au professeur vis-à-vis de la technique ? (Verbe : \all{entscheiden / erklären}).
    \item Trouvez dans le texte le participe passé de \textbf{herunterladen} et de \textbf{verbinden}.
    \item Traduisez en allemand : « Les outils numériques changent l'école. » (Utilisez \all{sich verändern} ou \all{verändern} + sujet approprié).
\end{enumerate}
\end{exercice}

\begin{corrige}[Exercice : Der digitale Klassenraum]
\begin{enumerate}
    \item \textbf{der Computer, das Tablet, das Handy}.
    \item La KI \textbf{analyse} (\all{analysiert}) les réponses, \textbf{reconnaît} (\all{erkennt}) les difficultés et \textbf{adapte} (\all{passt ... an}) les tâches.
    \item « ... à son propre rythme, apprendre et s'exercer. » (\all{in seinem eigenen Tempo lernen und üben}).
    \item Sur la \textbf{Lernplattform} (die Lernplattform, die Lernplattformen).
    \item Le professeur \textbf{décide} (\all{entscheidet}) quel logiciel est pertinent et \textbf{explique} (\all{erklärt}) les thèmes difficiles.
    \item \textbf{heruntergeladen} (l. 10), \textbf{verbunden} (l. 4, partic. passé comme adjectif : \all{sind ... verbunden} ; l. 16 : \all{verbindet} = présent).
    \item \textbf{Digitale Werkzeuge verändern die Schule.} ou \textbf{Der digitale Unterricht verändert die Schule.}
\end{enumerate}
\end{corrige}

\begin{popfigure}
\begin{tikzpicture}[
    mindmap,
    concept color=popblue,
    level 1 concept/.append style={level distance=3.5cm, sibling angle=90},
    level 2 concept/.append style={level distance=2.8cm},
    every node/.style={font=\small},
    text=white
]
\node[concept, align=center]{Digitale\\Medien}
    child[concept color=poporange] { node[concept] {Geräte}
        child { node[concept, scale=0.7] {Computer} }
        child { node[concept, scale=0.7] {Tablet} }
        child { node[concept, scale=0.7] {Handy} }
        child { node[concept, scale=0.7] {Bildschirm} }
        child { node[concept, scale=0.7] {Tastatur} }
        child { node[concept, scale=0.7] {Drucker} }
    }
    child[concept color=popgreen] { node[concept] {Programme}
        child { node[concept, scale=0.7] {App} }
        child { node[concept, scale=0.7] {Lernprogramm} }
        child { node[concept, scale=0.7] {Software} }
        child { node[concept, scale=0.7] {Internet} }
        child { node[concept, scale=0.7] {Plattform} }
        child { node[concept, scale=0.7] {KI} }
    }
    child[concept color=poppurple] { node[concept] {Verben}
        child { node[concept, scale=0.7] {lernen} }
        child { node[concept, scale=0.7] {üben} }
        child { node[concept, scale=0.7] {korrigieren} }
        child { node[concept, scale=0.7] {anpassen} }
        child { node[concept, scale=0.7] {herunterladen} }
        child { node[concept, scale=0.7] {speichern} }
        child { node[concept, scale=0.7] {teilen} }
    }
    child[concept color=poppink] { node[concept] {Personen}
        child { node[concept, scale=0.7] {Schüler/in} }
        child { node[concept, scale=0.7] {Lehrer/in} }
        child { node[concept, scale=0.7] {Experte/in} }
        child { node[concept, scale=0.7] {Unterricht} }
        child { node[concept, scale=0.7] {Aufgabe} }
        child { node[concept, scale=0.7] {Niveau} }
    };
\end{tikzpicture}
\legende{Carte mentale récapitulative : Digitale Medien — 4 branches principales (Geräte, Programme, Verben, Personen) avec vocabulaire clé de la leçon.}
\end{popfigure}

\begin{aretenir}[Check-list de révision — Lexique « Digitale Medien und KI »]
\begin{itemize}
    \item \textbf{Noms :} Article + Pluriel appris par cœur (\cle{Tablets}, \cle{Handys}, \cle{Computer}, \cle{Bildschirme}, \cle{Tastaturen}, \cle{Drucker}, \cle{Apps}, \cle{Lernprogramme}, \cle{Plattformen}, \cle{KI}, \cle{Schüler/innen}, \cle{Lehrer/innen}, \cle{Experten/innen}, \cle{Aufgaben}, \cle{Niveaus}).
    \item \textbf{Verbes :} Présens (3e sg/pl), Perfekt (auxiliaire + participe), Régime (Akk/Dat/Prép). Séparables : \cle{herunterladen}, \cle{anpassen}.
    \item \textbf{Familles :} \cle{lernen} $\rightarrow$ \cle{Lernprogramm}, \cle{Lerner}, \cle{Lernen} ; \cle{korrigieren} $\rightarrow$ \cle{Korrektur} ; \cle{intelligent} $\rightarrow$ \cle{Intelligenz}, \cle{KI}.
    \item \textbf{Expressions :} \cle{am Bildschirm arbeiten}, \cle{im Internet surfen}, \cle{KI benutzen}, \cle{Fehler korrigieren}, \cle{Niveau anpassen}.
    \item \textbf{Pièges :} \cle{lernen} (lycée) vs \cle{studieren} (univ) ; \cle{die Computer} (pas -s) ; \cle{das Handy} (nom, pas adj.) ; \cle{arbeiten an + Dat}.
\end{itemize}
\end{aretenir}

\coursec{Grammaire 1 : das Futur I}

Après avoir enrichi notre vocabulaire du numérique et de l'intelligence artificielle (\all{die künstliche Intelligenz, das Lernprogramm, die App, das Tablet, lernen}), nous allons maintenant observer la grammaire qui travaille dans notre texte \emph{Lernen mit der KI}. En le lisant, vous avez remarqué que l'auteur parle constamment de l'avenir de l'école. Comment l'allemand exprime-t-il le futur ? C'est l'objet de cette section : le \cle{Futur I}.

\courssub{Observation : trois phrases du texte}

Reprenons trois phrases-clés du texte support et observons-les attentivement :

\begin{exemplebox}[Phrases du texte \emph{Lernen mit der KI}]
\begin{enumerate}
\item \all{Die KI wird den Unterricht verändern.} « L'IA va transformer / transformera l'enseignement. »
\item \all{Sie wird den Lehrer nicht ersetzen.} « Elle ne remplacera pas le professeur. »
\item \all{Die Schule der Zukunft wird beides verbinden.} « L'école du futur combinera les deux. »
\end{enumerate}
\end{exemplebox}

Que constatons-nous ? Dans chaque phrase, il y a \textbf{deux verbes} :

\begin{itemize}
\item une forme de \all{werden} conjuguée au présent (\all{wird}), placée en \textbf{position 2}, comme tout verbe conjugué dans une phrase déclarative allemande ;
\item un \textbf{infinitif} (\all{verändern, ersetzen, verbinden}), rejeté tout à la \textbf{fin} de la phrase.
\end{itemize}

C'est exactement la même construction que le \all{Perfekt} (\all{Ich habe eine App benutzt.}), sauf que l'auxiliaire est \all{werden} et que le verbe principal reste à l'infinitif au lieu d'être un participe passé.

\begin{popfigure}
\begin{center}
\begin{tikzpicture}[
box/.style={draw=popblue, fill=popblueL, rounded corners=3pt, align=center, font=\small, inner sep=5pt},
box2/.style={draw=poporange, fill=poporangeL, rounded corners=3pt, align=center, font=\small, inner sep=5pt},
box3/.style={draw=popgreen, fill=popgreenL, rounded corners=3pt, align=center, font=\small, inner sep=5pt},
fleche/.style={-{Stealth[length=2.5mm]}, thick, popdark}
]
\node[box, align=center] (suj) at (0,0){Sujet\\ \all{Die KI}};
\node[box2, right=0.9cm of suj, align=center] (werd){\all{werden} conjugué\\ position 2\\ \all{wird}};
\node[box, right=0.9cm of werd, align=center] (mil){compléments\\ \all{den Unterricht}};
\node[box3, right=0.9cm of mil, align=center] (inf){Infinitif\\ fin de phrase\\ \all{verändern}};
\draw[fleche] (suj) -- (werd);
\draw[fleche] (werd) -- (mil);
\draw[fleche] (mil) -- (inf);
\draw[fleche, poppurple, dashed, bend left=35] (werd.north) to node[above, font=\small, text=poppurple]{cadre verbal} (inf.north);
\end{tikzpicture}
\end{center}
\legende{Structure du Futur I : \all{werden} en position 2 et l'infinitif à la fin forment le cadre verbal.}
\end{popfigure}

\courssub{Formation du Futur I}

\begin{propriete}[Règle de formation du Futur I]
Le \cle{Futur I} se forme avec :
\begin{center}
\textbf{\all{werden} conjugué au présent (position 2) + infinitif du verbe principal (fin de phrase)}
\end{center}
Attention : ici, \all{werden} est un \textbf{auxiliaire de temps} (comme \all{haben} ou \all{sein} au Perfekt). Il ne signifie \emph{pas} « devenir » dans cette construction. Il porte uniquement la marque de la personne et du futur.
\end{propriete}

Voici la conjugaison complète de \all{werden} au présent, avec le Futur I du verbe \all{lernen}. Notez que \all{werden} est irrégulier au présent : \all{du wirst}, \all{er/sie/es wird} (voyelle \all{e} $\rightarrow$ \all{i}).

\begin{center}
\begin{tabular}{|l|l|l|l|}
\hline
\rowcolor{popblueL} \textbf{Personne} & \textbf{\all{werden} (présent)} & \textbf{Futur I de \all{lernen}} & \textbf{Français} \\
\hline
ich & werde & ich werde lernen & j'apprendrai \\
\hline
du & wirst & du wirst lernen & tu apprendras \\
\hline
er / sie / es & wird & er / sie / es wird lernen & il / elle apprendra \\
\hline
wir & werden & wir werden lernen & nous apprendrons \\
\hline
ihr & werdet & ihr werdet lernen & vous apprendrez \\
\hline
sie / Sie & werden & sie / Sie werden lernen & ils / vous apprendrez \\
\hline
\end{tabular}
\end{center}

\begin{attention}[Ne confondez pas les deux emplois de \all{werden} !]
Le verbe \all{werden} a deux vies :
\begin{itemize}
\item \textbf{Verbe plein} = « devenir », suivi d'un nom ou d'un adjectif : \all{Er wird Lehrer.} « Il devient professeur. » / \all{Es wird kalt.} « Il fait de plus en plus froid. »
\item \textbf{Auxiliaire du futur}, suivi d'un \textbf{infinitif en fin de phrase} : \all{Er wird lernen.} « Il apprendra. »
\end{itemize}
Le test est simple : s'il y a un infinitif à la fin de la phrase, \all{werden} est l'auxiliaire du futur.
\end{attention}

\begin{exemplebox}[Exemples traduits]
\begin{itemize}
\item \all{Ich werde eine App benutzen.} « J'utiliserai une application. »
\item \all{Wir werden morgen üben.} « Nous nous entraînerons demain. »
\item \all{Du wirst die Prüfung bestehen.} « Tu réussiras l'examen. »
\item \all{Ihr werdet mit dem Tablet arbeiten.} « Vous travaillerez avec la tablette. »
\item \all{Die Schüler werden mit einem Lernprogramm lesen lernen.} « Les élèves apprendront à lire avec un programme d'apprentissage. »
\end{itemize}
\end{exemplebox}

\courssub{Les trois emplois du Futur I}

\textbf{1. La prévision, la prédiction.} C'est l'emploi principal, celui de notre texte : on annonce ce qui arrivera. On trouve souvent des adverbes comme \all{sicher} (sûrement), \all{wahrscheinlich} (probablement), \all{bestimmt} (certainement), \all{bald} (bientôt), \all{in der Zukunft} (à l'avenir).

\begin{exemplebox}[Prévision]
\begin{itemize}
\item \all{Die KI wird sicher wichtiger werden.} « L'IA deviendra sûrement plus importante. »
\item \all{In der Zukunft werden viele Schüler mit Apps lernen.} « À l'avenir, beaucoup d'élèves apprendront avec des applications. »
\item \all{Der Computer wird das Lernen wahrscheinlich leichter machen.} « L'ordinateur rendra probablement l'apprentissage plus facile. »
\end{itemize}
\end{exemplebox}

\textbf{2. Le projet, l'intention.} On annonce ce qu'on a décidé de faire.

\begin{exemplebox}[Projet, intention]
\begin{itemize}
\item \all{Ich werde nächstes Jahr Deutsch lernen.} « J'apprendrai l'allemand l'année prochaine. »
\item \all{Wir werden nach dem Abitur in Deutschland studieren.} « Nous étudierons en Allemagne après le baccalauréat. »
\end{itemize}
\end{exemplebox}

\textbf{3. La supposition sur le présent.} Avec \all{wohl} ou \all{schon}, le Futur I exprime une probabilité sur le moment présent ; le français traduit par « probablement, sans doute » + présent.

\begin{exemplebox}[Supposition]
\begin{itemize}
\item \all{Er wird jetzt wohl zu Hause sein.} « Il est probablement à la maison maintenant. »
\item \all{Sie wird schon Recht haben.} « Elle a sans doute raison. »
\end{itemize}
\end{exemplebox}

\begin{popfigure}
\begin{center}
\begin{tikzpicture}[font=\small]
\draw[-{Stealth[length=3mm]}, very thick, popdark] (0,0) -- (10.5,0);
\foreach \x/\txt in {1.5/Vergangenheit, 5.2/Gegenwart, 8.8/Zukunft}{
  \draw[thick, popdark] (\x,-0.15) -- (\x,0.15);
  \node[below, align=center] at (\x,-0.2) {\txt};
}
\node[below, popmuted] at (1.5,-0.9) {passé};
\node[below, popmuted] at (5.2,-0.9) {présent (\all{jetzt})};
\node[below, popmuted] at (8.8,-0.9) {futur};
\draw[fill=poporange, draw=poporange] (8.8,0.35) circle (0.12);
\node[above, text=poporange, align=center] at (8.8,0.5) {\textbf{Futur I}\\ \all{Die KI wird lernen.}};
\draw[fill=popgreen, draw=popgreen] (1.5,0.35) circle (0.12);
\node[above, text=popgreen, align=center] at (1.5,0.5) {Perfekt / Präteritum\\ \all{Sie hat gelernt.}};
\draw[fill=popblue, draw=popblue] (5.2,0.35) circle (0.12);
\node[above, text=popblue, align=center] at (5.2,0.5) {Présent\\ \all{Sie lernt.}};
\end{tikzpicture}
\end{center}
\legende{Frise des temps : le Futur I se situe après \all{jetzt} (maintenant), dans la zone du futur.}
\end{popfigure}

\courssub{Comparaison avec le français : un piège majeur}

\begin{propriete}[Le présent peut exprimer le futur en allemand]
Contrairement au français, qui exige le futur simple après « quand » ou pour un événement à venir, l'allemand utilise très souvent le \textbf{présent accompagné d'un adverbe de temps} (\all{morgen, bald, nächstes Jahr}) quand le futur est évident par le contexte :
\begin{itemize}
\item \all{Morgen fahre ich nach Berlin.} « Demain, j'irai à Berlin. » (et non \all{werde ich fahren})
\item \all{Nächste Woche beginnen die Ferien.} « Les vacances commenceront la semaine prochaine. »
\end{itemize}
Le Futur I n'est donc pas obligatoire : il sert surtout à la \textbf{prédiction}, à la \textbf{promesse} ou à la \textbf{supposition}.
\end{propriete}

\begin{center}
\begin{tabularx}{\linewidth}{|X|X|X|}
\hline
\rowcolor{popblueL} \textbf{Français} & \textbf{Deutsch} & \textbf{Remarque} \\
\hline
Demain, j'irai à Berlin. & \all{Morgen fahre ich nach Berlin.} & présent + adverbe : très fréquent \\
\hline
Quand il viendra, nous mangerons. & \all{Wenn er kommt, essen wir.} & pas de futur après \all{wenn} \\
\hline
L'IA transformera l'école. (prédiction) & \all{Die KI wird die Schule verändern.} & Futur I : prédiction, opinion \\
\hline
Il est probablement malade. & \all{Er wird wohl krank sein.} & Futur I + \all{wohl} : supposition \\
\hline
\end{tabularx}
\end{center}

\courssub{Négation, interrogation et subordonnée}

\textbf{La négation.} \all{nicht} se place devant l'infinitif ou devant le complément nié, jamais juste après \all{werden} :

\begin{exemplebox}[Négation au Futur I]
\begin{itemize}
\item \all{Sie wird den Lehrer nicht ersetzen.} « Elle ne remplacera pas le professeur. »
\item \all{Ich werde das Tablet nicht kaufen.} « Je n'achèterai pas la tablette. »
\item \all{Wir werden morgen nicht arbeiten.} « Nous ne travaillerons pas demain. »
\end{itemize}
\end{exemplebox}

\textbf{La question.} Comme toujours en allemand, le verbe conjugué passe en position 1 dans les questions sans mot interrogatif ; l'infinitif reste à la fin :

\begin{exemplebox}[Interrogation au Futur I]
\begin{itemize}
\item \all{Wird die KI den Lehrer ersetzen?} « L'IA remplacera-t-elle le professeur ? »
\item \all{Wirst du eine App benutzen?} « Utiliseras-tu une application ? »
\item \all{Wann werden wir mit dem Lernprogramm arbeiten?} « Quand travaillerons-nous avec le programme d'apprentissage ? »
\end{itemize}
\end{exemplebox}

\textbf{La subordonnée.} Dans une subordonnée introduite par \all{dass}, \all{weil}, \all{wenn}, le verbe conjugué (\all{werden}) est rejeté à la \textbf{toute fin} de la subordonnée, \emph{après} l'infinitif :

\begin{exemplebox}[Le Futur I en subordonnée]
\begin{itemize}
\item \all{Ich denke, dass er die Prüfung bestehen wird.} « Je pense qu'il réussira l'examen. »
\item \all{Ich glaube, dass die KI den Unterricht verändern wird.} « Je crois que l'IA transformera l'enseignement. »
\item \all{Sie hofft, dass wir das Lernprogramm bald benutzen werden.} « Elle espère que nous utiliserons bientôt le programme d'apprentissage. »
\end{itemize}
Ordre en fin de subordonnée : \textbf{infinitif + \all{werden}} (\all{bestehen wird}), jamais l'inverse.
\end{exemplebox}

\begin{attention}[Erreurs fréquentes des francophones]
\begin{enumerate}
\item \textbf{Ne placez jamais l'infinitif juste après \all{werden}} : \all{Ich werde lernen Deutsch} ✗. En français, « j'apprendrai l'allemand » a le complément après le verbe ; en allemand, l'infinitif va \textbf{toujours en fin de phrase} : \all{Ich werde Deutsch lernen.} ✓
\item \textbf{N'oubliez pas l'infinitif} : \all{Ich werde morgen nach Yaoundé} ✗ n'a pas de sens ; il faut \all{Ich werde morgen nach Yaoundé fahren.} ✓
\item \textbf{Pas de futur après \all{wenn}} (quand) pour un fait futur : \all{Wenn er kommt, ...} et non \all{wenn er kommen wird}.
\item \textbf{En subordonnée}, \all{werden} va après l'infinitif : \all{..., dass sie lernen wird} ✓ ; \all{..., dass sie wird lernen} ✗.
\end{enumerate}
\end{attention}

\courssub{Texte d'application : \emph{Meine Zukunft mit der KI}}

Lisons ce court texte original et soulignons mentalement tous les Futur I.

\begin{textebox}[Meine Zukunft mit der KI]
\all{Nächstes Jahr werde ich mein Abitur machen, und danach werde ich Informatik studieren. Ich bin sicher: Die künstliche Intelligenz wird mein Leben stark verändern. Ich werde jeden Tag mit Apps und Lernprogrammen arbeiten, aber ich werde die Bücher nicht vergessen. Mein Bruder wird wahrscheinlich Lehrer werden. Er sagt: „Die KI wird mir helfen, aber sie wird mich nicht ersetzen.“ Ich glaube, dass er Recht haben wird. Die Schule der Zukunft wird Mensch und Maschine verbinden. Werden wir eines Tages nur noch mit dem Computer lernen? Ich denke, dass wir den Kontakt mit den Lehrern und mit den Freunden immer brauchen werden. Die Technik wird uns unterstützen, aber sie wird nicht alles machen. Und du? Wirst du auch mit der KI lernen?}

\textbf{Glossaire} : das Abitur (le baccalauréat) ; die Informatik (l'informatique) ; stark (fortement) ; vergessen (oublier) ; ersetzen (remplacer) ; Recht haben (avoir raison) ; verbinden (combiner, relier) ; der Kontakt, -e (le contact) ; unterstützen (soutenir, aider) ; eines Tages (un jour).
\end{textebox}

Questions de compréhension : 1. \all{Was wird der Autor nach dem Abitur machen?} 2. \all{Was wird sein Bruder wahrscheinlich werden?} 3. \all{Was wird die KI nach dem Bruder nicht machen?} 4. Relevez trois Futur I du texte et traduisez-les.

\courssub{Exercice résolu et entraînement}

\begin{exoresolu}[Conjugaison au Futur I]
Mettez les phrases suivantes au Futur I : 1. \all{Die KI (verändern) den Unterricht.} 2. \all{Ich (lernen) nächstes Jahr Deutsch.} 3. \all{Wir (benutzen) morgen ein Tablet.} 4. \all{Die Schüler (arbeiten) mit einer App.}
\tcblower
\textbf{Solution détaillée.} Méthode : on conjugue \all{werden} à la personne du sujet en position 2, et on place l'infinitif à la fin.
\begin{enumerate}
\item \all{Die KI wird den Unterricht verändern.} (3\textsuperscript{e} personne du singulier : \all{wird})
\item \all{Ich werde nächstes Jahr Deutsch lernen.} (1\textsuperscript{re} personne : \all{werde} ; l'infinitif \all{lernen} à la fin, après le complément \all{Deutsch})
\item \all{Wir werden morgen ein Tablet benutzen.} (\all{werden} ; attention à ne pas coller \all{benutzen} à \all{werden})
\item \all{Die Schüler werden mit einer App arbeiten.} (\all{die Schüler} = 3\textsuperscript{e} personne du pluriel : \all{werden})
\end{enumerate}
\end{exoresolu}

\begin{exercice}[À vous de jouer]
\begin{enumerate}
\item Traduisez en allemand : a) « J'utiliserai une application demain. » b) « Tu réussiras l'examen. » c) « L'IA ne remplacera pas le professeur. » d) « Je pense que nous apprendrons avec des tablettes. »
\item Transformez en question : \all{Die KI wird den Unterricht verändern.}
\item Mettez à la forme négative : \all{Er wird das Lernprogramm kaufen.}
\end{enumerate}
\end{exercice}

\begin{corrige}[Exercice : À vous de jouer]
\begin{enumerate}
\item a) \all{Ich werde morgen eine App benutzen.} b) \all{Du wirst die Prüfung bestehen.} c) \all{Die KI wird den Lehrer nicht ersetzen.} d) \all{Ich denke, dass wir mit Tablets lernen werden.} (subordonnée : \all{werden} à la toute fin, après l'infinitif)
\item \all{Wird die KI den Unterricht verändern?} (verbe conjugué en position 1)
\item \all{Er wird das Lernprogramm nicht kaufen.} (\all{nicht} devant l'infinitif)
\end{enumerate}
\end{corrige}

\begin{aretenir}[L'essentiel du Futur I]
\begin{itemize}
\item Formation : \all{werden} au présent (position 2) + infinitif en fin de phrase : \all{ich werde, du wirst, er wird, wir werden, ihr werdet, sie/Sie werden}.
\item Emplois : prédiction (\all{sicher, wahrscheinlich, bald}), projet, supposition sur le présent (\all{wohl, schon}).
\item Pour un futur évident, l'allemand préfère souvent le présent + adverbe : \all{Morgen fahre ich nach Berlin.}
\item Négation : \all{nicht} devant l'infinitif ; question : \all{werden} en position 1 ; subordonnée : infinitif puis \all{werden} à la fin (\all{..., dass er bestehen wird}).
\item \all{werden} + infinitif = futur ; \all{werden} + nom/adjectif = « devenir ».
\end{itemize}
\end{aretenir}

\coursec{Grammaire 2 : die Infinitivgruppen}

Après avoir étudié le \emph{Futur I} (section précédente), qui exprime une action future ou une supposition, nous abordons maintenant une structure fondamentale de la syntaxe allemande : les \textbf{groupes infinitifs} (\cle{Infinitivgruppen}). Ils permettent d'exprimer le but, la négation d'action, la substitution ou de compléter certains verbes, tout en allégeant la phrase (pas de sujet répété, pas de conjonction de subordination comme \all{dass} ou \all{weil}). Vous les retrouverez constamment dans les textes journalistiques, les consignes d'exercices et les productions écrites au Bac.

\courssub{Observation dans le texte support}

Reprenons deux phrases extraites du texte \emph{Lernen mit der KI} (section 1) :

\begin{textebox}[Extrait du texte support]
Um gut zu lernen, braucht man Kontakt zu anderen Menschen. Wer lernt, ohne mit anderen zu sprechen, verliert die Motivation.
\end{textebox}

\textbf{Analyse :}
\begin{itemize}
    \item \all{Um gut zu lernen} : groupe infinitif introduit par \all{um ... zu} placé \emph{en tête de phrase} (position 1), suivi du verbe conjugué \all{braucht} (position 2). Sens : \og\ \emph{pour bien apprendre}\ \fg{} (but).
    \item \all{ohne mit anderen zu sprechen} : groupe infinitif introduit par \all{ohne ... zu} à l'intérieur de la proposition relative \all{Wer lernt, ...}. Sens : \og\ \emph{sans parler avec d'autres}\ \fg{} (négation d'action).
    \item Dans les deux cas, l'infinitif (\all{lernen}, \all{sprechen}) est en \textbf{position finale} du groupe, précédé de \all{zu}.
    \item Une \textbf{virgule} sépare le groupe infinitif du reste de la phrase (obligatoire devant \all{um ... zu}, \all{ohne ... zu}, \all{anstatt ... zu} et tout groupe infinitif développé).
\end{itemize}

\courssub{Structure générale : \texorpdfstring{\cle{zu + Infinitiv}}{zu + Infinitif} en fin de groupe}

\begin{propriete}[Structure de base du groupe infinitif]
Un groupe infinitif (Infinitivgruppe) se compose de :
\begin{enumerate}
    \item une \textbf{particule introductrice} (souvent \all{um}, \all{ohne}, \all{anstatt}, \all{statt}) — facultative après certains verbes ;
    \item le mot \cle{zu} ;
    \item l'\textbf{infinitif du verbe} (éventuellement avec ses compléments) — \textbf{toujours en position finale du groupe}.
\end{enumerate}
Le groupe infinitif n'a \textbf{pas de sujet propre} : son sujet est celui de la proposition principale (ou de la proposition dont il dépend).
\end{propriete}

\begin{exemplebox}[Exemples de base]
\all{Ich habe vor, morgen zu lernen.} — J'ai l'intention d'apprendre demain. \\
\all{Sie vergisst, die App herunterzuladen.} — Elle oublie de télécharger l'appli. \\
\all{Wir beginnen, die Aufgaben zu lösen.} — Nous commençons à résoudre les exercices.
\end{exemplebox}

\courssub{Cas particulier : verbes à particule séparable (trennbare Verben)}

Avec un verbe séparable (\all{herunterladen}, \all{anpassen}, \all{aufgeben}, \all{mitmachen}...), \cle{zu s'insère ENTRE la particule et le radical verbal}.

\begin{propriete}[Verbes séparables dans les groupes infinitifs]
\begin{center}
\begin{tabular}{|l|l|l|}
\hline
\rowcolor{popblueL}
Verbe séparable & Décomposition & Groupe infinitif correct \\
\hline
herunterladen & herunter + laden & herunter\cle{zu}laden \\
anpassen & an + passen & an\cle{zu}passen \\
aufgeben & auf + geben & auf\cle{zu}geben \\
mitmachen & mit + machen & mit\cle{zu}machen \\
teilnehmen & teil + nehmen & teil\cle{zu}nehmen \\
\hline
\end{tabular}
\end{center}
\emph{Règle mnémotechnique :} \og\ La particule et le verbe s'embrassent, \cle{zu} se glisse entre eux.\ \fg{}
\end{propriete}

\begin{popfigure}
\begin{tikzpicture}[
    node distance=0.8cm and 1.2cm,
    box/.style={draw=popblue, thick, rounded corners, fill=popblueL, minimum height=1cm, text width=2.2cm, align=center, font=\small},
    arrow/.style={-Stealth, thick, poporange}
]
\node[box, align=center] (part){Particule\\ \textbf{herunter}};
\node[box, right=of part] (zu) {\textbf{zu}};
\node[box, right=of zu, align=center] (verb){Radical\\ \textbf{laden}};
\node[box, below=1.2cm of zu, fill=popgreenL, draw=popgreen, text width=5cm] (result) {Résultat : \textbf{herunterzuladen}};

\draw[arrow] (part) -- (zu);
\draw[arrow] (zu) -- (verb);
\draw[arrow, poppurple] (part.south) |- (result.west);
\draw[arrow, poppurple] (zu.south) -- (result.north);
\draw[arrow, poppurple] (verb.south) |- (result.east);
\end{tikzpicture}
\legende{Insertion de \textit{zu} dans un verbe à particule séparable}
\end{popfigure}

\begin{attention}[Piège fréquent : verbes séparables]
\og\ \all{Das Programm passt die Aufgaben an, um besser zu helfen.}\ \fg{} (verbe principal \all{anpassen} conjugué : particule \all{an} en fin de proposition).
\og\ \all{..., um die Aufgaben anzupassen.}\ \fg{} (groupe infinitif : \cle{zu} entre \all{an} et \all{passen} $\to$ \all{anzupassen}).
\textbf{Ne jamais écrire :} \all{*um anzupassen die Aufgaben} (ordre incorrect) ni \all{*um anpassen zu} (zu mal placé).
\end{attention}

\courssub{Les cinq grands cas d'emploi}

\begin{definition}[Cas 1 — \texorpdfstring{\all{um ... zu + Infinitiv}}{um ... zu + Infinitif} : le BUT ( = \og pour\ \fg{})]
Exprime la finalité, l'objectif visé. Équivalent français : \og\ pour + infinitif\ \fg{}, \og\ afin de + infinitif\ \fg{}.
\begin{itemize}
    \item Sujet du groupe infinitif = sujet de la proposition principale.
    \item Virgule obligatoire devant \all{um}.
\end{itemize}
\end{definition}

\begin{exemplebox}[um ... zu — But]
\all{Ich lerne Deutsch, \textbf{um in Deutschland zu studieren}.} — J'apprends l'allemand \textbf{pour étudier en Allemagne}. \\
\all{\textbf{Um die Prüfung zu bestehen}, muss man viel üben.} — \textbf{Pour réussir l'examen}, il faut beaucoup s'entraîner. (Groupe en position 1) \\
\all{Wir nutzen die App, \textbf{um den Stoff besser zu verstehen}.} — Nous utilisons l'appli \textbf{pour mieux comprendre la matière}.
\end{exemplebox}

\begin{definition}[Cas 2 — \texorpdfstring{\all{ohne ... zu + Infinitiv}}{ohne ... zu + Infinitif} : la NÉGATION d'action ( = \og sans\ \fg{})]
Exprime qu'une action se fait \emph{sans} qu'une autre action ne soit accomplie. Équivalent français : \og\ sans + infinitif\ \fg{} (souvent gérondif en \emph{-ant}).
\begin{itemize}
    \item Sujet identique dans les deux propositions.
    \item Virgule obligatoire devant \all{ohne}.
\end{itemize}
\end{definition}

\begin{exemplebox}[ohne ... zu — Négation]
\all{Er lernt, \textbf{ohne zu üben}.} — Il apprend \textbf{sans s'entraîner} / \textbf{sans étudier}. \\
\all{Sie geht nach Hause, \textbf{ohne ihre Hausaufgaben zu machen}.} — Elle rentre à la maison \textbf{sans faire ses devoirs}. \\
\all{\textbf{Ohne das Buch zu lesen}, hat er eine Meinung dazu.} — \textbf{Sans avoir lu le livre}, il a un avis là-dessus.
\end{exemplebox}

\begin{definition}[Cas 3 — \texorpdfstring{\all{anstatt ... zu + Infinitiv} / \all{statt ... zu + Infinitiv}}{anstatt ... zu + Infinitiv / statt ... zu + Infinitif} : la SUBSTITUTION ( = \og au lieu de\ \fg{})]
Exprime qu'une action remplace une autre action attendue ou logique. \all{statt} est plus familier, \all{anstatt} plus formel (préféré à l'écrit).
\begin{itemize}
    \item Sujet identique.
    \item Virgule obligatoire devant \all{anstatt} / \all{statt}.
\end{itemize}
\end{definition}

\begin{exemplebox}[anstatt ... zu — Substitution]
\all{Sie spielt am Handy, \textbf{anstatt zu lernen}.} — Elle joue sur son portable \textbf{au lieu d'étudier}. \\
\all{\textbf{Anstatt am Bildschirm zu spielen}, solltest du lernen.} — \textbf{Au lieu de jouer devant l'écran}, tu devrais travailler. \\
\all{Er schreibt die Arbeit \textbf{statt sie nur zu kopieren}.} — Il rédige le devoir \textbf{au lieu de le copier seulement}.
\end{exemplebox}

\begin{definition}[Cas 4 — \texorpdfstring{\all{zu + Infinitiv}}{zu + Infinitif} après certains verbes (Verben mit \all{zu}-Infinitiv)]
Certains verbes appellent obligatoirement un complément infinitif introduit par \cle{zu} (sans \all{um}, \all{ohne}, \all{anstatt}). Les plus fréquents au programme :
\begin{center}
\begin{tabular}{|l|l|l|}
\hline
\rowcolor{popblueL}
Verbe allemand & Français & Exemple \\
\hline
\all{hoffen} & espérer & \all{Ich hoffe, die Prüfung zu bestehen.} \\
\all{versuchen} & essayer & \all{Er versucht, das Problem zu lösen.} \\
\all{beginnen / anfangen} & commencer & \all{Wir beginnen, den Text zu lesen.} \\
\all{vergessen} & oublier & \all{Vergiss nicht, die App herunterzuladen!} \\
\all{planen} & planifier & \all{Wir planen, nächstes Jahr nach Berlin zu fahren.} \\
\all{versprechen} & promettre & \all{Er verspricht, pünktlich zu kommen.} \\
\all{ablehnen} & refuser & \all{Sie lehnt ab, mitzumachen.} \\
\all{beabsichtigen} & avoir l'intention & \all{Ich beabsichtige, Medizin zu studieren.} \\
\hline
\end{tabular}
\end{center}
\textbf{Remarque :} Ces verbes sont souvent suivis d'une virgule si le groupe infinitif est développé (compléments présents).
\end{definition}

\begin{exemplebox}[Verbes appelant \texorpdfstring{\all{zu}}{zu} + infinitif]
\all{Ich hoffe, \textbf{bald flüssig sprechen zu können}.} — J'espère \textbf{pouvoir parler couramment bientôt}. \\
\all{Versuch bitte, \textbf{ruhig zu bleiben}!} — Essaie s'il te plaît de \textbf{rester calme} ! \\
\all{Wir haben vergessen, \textbf{die Datei zu speichern}.} — Nous avons oublié \textbf{d'enregistrer le fichier}.
\end{exemplebox}

\begin{definition}[Cas 5 — Verbes SANS \texorpdfstring{\all{zu}}{zu} (Infinitiv ohne \all{zu})]
Les verbes suivants s'emploient \textbf{directement} avec l'infinitif, \textbf{sans \all{zu}} :
\begin{itemize}
    \item \textbf{Verbes modaux} : \all{können, müssen, wollen, sollen, dürfen, möchten} \\
          \all{Ich kann gut lernen.} (pas de \all{zu} !)
    \item \textbf{Verbes de perception} : \all{sehen, hören, fühlen, spüren} \\
          \all{Ich höre ihn singen.} — Je l'entends chanter.
    \item \textbf{Verbes causatifs / permissifs} : \all{lassen} \\
          \all{Der Lehrer lässt die Schüler arbeiten.} — Le prof laisse les élèves travailler.
    \item \textbf{Verbes de mouvement} (surtout \all{gehen, kommen, fahren, laufen} + infinitif de but) : \\
          \all{Wir gehen schwimmen.} — Nous allons nager. / \all{Er kommt helfen.} — Il vient aider.
\end{itemize}
\end{definition}

\begin{attention}[Erreur TRÈS fréquente des francophones : \texorpdfstring{\all{*können zu}}{*können zu}]
\og\ En français, on dit \emph{je \textbf{peux} faire} (pas de préposition). En allemand aussi : \textbf{pas de \all{zu} après un modal}.\ \fg{}
\begin{center}
\begin{tabular}{|l|l|}
\hline
\rowcolor{poporangeL}
\all{\textbf{Faux}} & \all{\textbf{Juste}} \\
\hline
\all{*Ich kann zu lernen.} & \all{Ich kann lernen.} \\
\all{*Er muss zu gehen.} & \all{Er muss gehen.} \\
\all{*Wir wollen zu helfen.} & \all{Wir wollen helfen.} \\
\all{*Sie darf zu bleiben.} & \all{Sie darf bleiben.} \\
\hline
\end{tabular}
\end{center}
\emph{Astuce :} Les modaux sont des \og\ verbes auxiliaires de modalité\ \fg{} — ils se comportent comme \all{werden} au Futur I : verbe conjugué en position 2, infinitif nu en fin de phrase.
\end{attention}

\courssub{Tableau récapitulatif comparatif Français / Allemand}

\begin{center}
\begin{tabularx}{\linewidth}{|X|X|X|}
\hline
\rowcolor{popblueL}
\textbf{Français} & \textbf{Deutsch} & \textbf{Remarque clé} \\
\hline
\og\ Pour + infinitif\ \fg{} (but) & \all{um ... zu + Infinitiv} & \cle{Virgule devant \all{um}} ; \all{zu} + infinitif en \textbf{fin de groupe}. \\
\hline
\og\ Sans + infinitif\ \fg{} (négation) & \all{ohne ... zu + Infinitiv} & \cle{Virgule devant \all{ohne}} ; même sujet. \\
\hline
\og\ Au lieu de + infinitif\ \fg{} (substitution) & \all{anstatt ... zu + Infinitiv} / \all{statt ... zu + Infinitiv} & \cle{Virgule devant \all{anstatt}/\all{statt}} ; \all{anstatt} plus formel. \\
\hline
Verbe + \og\ à/de + infinitif\ \fg{} & Verbe + \all{zu + Infinitiv} & Liste fermée : \all{hoffen, versuchen, beginnen, vergessen, planen, versprechen, ablehnen, beabsichtigen...} \\
\hline
Verbe modal (\emph{pouvoir, devoir, vouloir...}) + infinitif & Modalverben (\all{können, müssen, wollen...}) + \textbf{Infinitif SANS \all{zu}} & \textbf{JAMAIS de \all{zu} après un modal !} \\
\hline
Verbe de perception (\emph{voir, entendre}) + infinitif & \all{sehen, hören, fühlen} + \textbf{Infinitif SANS \all{zu}} & Construction \og\ Accusatif + Infinitif\ \fg{} (\all{Ich sehe ihn kommen}). \\
\hline
\og\ Laisser\ \fg{} + infinitif & \all{lassen} + \textbf{Infinitif SANS \all{zu}} & \all{Er lässt mich warten.} \\
\hline
\end{tabularx}
\end{center}

\courssub{La virgule : règle d'or}

\begin{aretenir}[La virgule devant les groupes infinitifs]
Une virgule \textbf{sépare obligatoirement} le groupe infinitif de la proposition principale (ou de la proposition dont il dépend) dans les cas suivants :
\begin{enumerate}
    \item Devant \all{um ... zu}, \all{ohne ... zu}, \all{anstatt ... zu}, \all{statt ... zu}.
    \item Devant \textbf{tout groupe infinitif développé} (avec compléments) après les verbes du Cas 4 (\all{hoffen, versuchen, vergessen...}).
    \item \textbf{Pas de virgule} si le groupe infinitif est \textbf{court} (seulement \all{zu} + infinitif, sans complément) après ces verbes — mais la virgule reste \emph{correcte et recommandée} pour la clarté.
\end{enumerate}
\end{aretenir}

\begin{exemplebox}[Virgule : oui / non]
\all{Ich hoffe, die Prüfung zu bestehen.} (virgule recommandée, groupe développé) \\
\all{Ich hoffe zu bestehen.} (virgule facultative, groupe court) \\
\all{Er lernt, ohne zu üben.} (virgule \textbf{obligatoire} devant \all{ohne}) \\
\all{Um zu lernen, braucht man Zeit.} (virgule \textbf{obligatoire} après groupe en position 1)
\end{exemplebox}

\courssub{Exercice résolu : transformation de phrases}

\begin{exoresolu}[Transformation : propositions subordonnées $\to$ groupes infinitifs]
\textbf{Énoncé :} Transformez chaque phrase en remplaçant la proposition subordonnée (introduite par \all{dass}, \all{weil}, \all{obwohl}) par un groupe infinitif approprié (\all{um ... zu}, \all{ohne ... zu}, \all{anstatt ... zu} ou \all{zu} seul). Attention aux verbes séparables !

\begin{enumerate}
    \item \all{Ich gehe in die Bibliothek, weil ich lernen will.}
    \item \all{Er macht die Hausaufgaben nicht, obwohl er sie machen soll.}
    \item \all{Sie hat vor, dass sie die App herunterlädt.}
    \item \all{Wir versuchen, dass wir das Problem lösen.}
    \item \all{Man kann nicht gut lernen, wenn man nicht übt.}
\end{enumerate}

\tcblower
\textbf{Solution détaillée :}

\begin{enumerate}
    \item \all{Ich gehe in die Bibliothek, \textbf{um zu lernen}.} \\
          (But : \all{um ... zu} ; sujet identique \all{ich} ; \all{wollen} $\to$ infinitif \all{lernen} ; pas de \all{zu} après \all{wollen} dans la version originale, mais ici structure \all{um ... zu} impose \all{zu}.)
    \item \all{Er macht die Hausaufgaben nicht, \textbf{anstatt sie zu machen}.} \quad (\textit{ou :} \all{... ohne sie zu machen.}) \\
          (Substitution / Négation : \all{anstatt ... zu} ou \all{ohne ... zu} ; \all{sollen} disparaît, infinitif \all{machen} avec \all{zu}.)
    \item \all{Sie hat vor, \textbf{die App herunterzuladen}.} \\
          (Verbe \all{vorhaben} appelle \all{zu} + infinitif ; verbe séparable \all{herunterladen} $\to$ \all{herunterzuladen}.)
    \item \all{Wir versuchen, \textbf{das Problem zu lösen}.} \\
          (Verbe \all{versuchen} appelle \all{zu} + infinitif.)
    \item \all{Man kann nicht gut lernen, \textbf{ohne zu üben}.} \\
          (Négation d'action : \all{ohne ... zu} ; \all{können} est modal $\to$ pas de \all{zu} après \all{kann} ; \all{üben} avec \all{zu} dans le groupe \all{ohne ... zu}.)
\end{enumerate}
\end{exoresolu}

\courssub{Entraînement : exercices à faire}

\begin{exercice}[Entraînement : groupes infinitifs]
\textbf{A.} Complétez par la structure infinitivale correcte (\all{um ... zu}, \all{ohne ... zu}, \all{anstatt ... zu}, \all{zu} seul, ou \textbf{rien} si verbe modal/perception/laisser).

\begin{enumerate}
    \item Ich lerne Deutsch, \_\_\_\_\_\_\_\_\_\_\_\_ in Deutschland zu studieren.
    \item Er geht nach Hause, \_\_\_\_\_\_\_\_\_\_\_\_ die Hausaufgaben zu machen.
    \item Wir versuchen, das Problem \_\_\_\_\_\_\_\_\_\_\_\_ (lösen).
    \item Sie kann \_\_\_\_\_\_\_\_\_\_\_\_ sehr gut schwimmen.
    \item Ich höre ihn \_\_\_\_\_\_\_\_\_\_\_\_ Klavier spielen.
    \item Der Lehrer lässt die Schüler \_\_\_\_\_\_\_\_\_\_\_\_ selbst arbeiten.
    \item \_\_\_\_\_\_\_\_\_\_\_\_ die Prüfung zu bestehen, muss man viel üben.
    \item Ich habe vergessen, \_\_\_\_\_\_\_\_\_\_\_\_ die Geräte \_\_\_\_\_\_\_\_\_\_\_\_ (aufladen).
    \item \_\_\_\_\_\_\_\_\_\_\_\_ am Handy zu spielen, solltest du lernen.
    \item Wir planen, \_\_\_\_\_\_\_\_\_\_\_\_ nächstes Jahr nach Österreich \_\_\_\_\_\_\_\_\_\_\_\_ (fahren).
\end{enumerate}

\textbf{B.} Réécrivez les phrases en utilisant un groupe infinitif (le sujet est le même dans les deux propositions).

\begin{enumerate}
    \item Ich möchte, dass ich die App benutze. $\to$ Ich möchte \_\_\_\_\_\_\_\_\_\_\_\_.
    \item Er lernt nicht, weil er faul ist. (Utilisez \all{ohne ... zu}) $\to$ Er lernt \_\_\_\_\_\_\_\_\_\_\_\_.
    \item Wir essen zu Mittag, danach gehen wir spazieren. (Utilisez \all{um ... zu}) $\to$ Wir essen zu Mittag, \_\_\_\_\_\_\_\_\_\_\_\_.
    \item Sie liest das Buch nicht, sie schaut nur den Film. (Utilisez \all{anstatt ... zu}) $\to$ Sie schaut nur den Film, \_\_\_\_\_\_\_\_\_\_\_\_.
\end{enumerate}
\end{exercice}

\begin{corrige}[Exercice : Groupes infinitifs]
\textbf{A.}
\begin{enumerate}
    \item \all{um} (but)
    \item \all{ohne} (négation d'action)
    \item \all{zu lösen} (verbe \all{versuchen} appelle \all{zu} devant l'infinitif final)
    \item \all{} (modal \all{kann} $\to$ \textbf{pas de \all{zu}})
    \item \all{} (verbe de perception \all{höre} $\to$ \textbf{pas de \all{zu}})
    \item \all{} (verbe \all{lässt} $\to$ \textbf{pas de \all{zu}})
    \item \all{Um} (but, groupe en position 1)
    \item \all{zu} / \all{aufzuladen} (verbe \all{vergessen} + \all{zu} ; verbe séparable \all{aufladen} $\to$ \all{aufzuladen})
    \item \all{Anstatt} (substitution)
    \item \all{zu} / \all{zu fahren} (verbe \all{planen} + \all{zu} ; \all{fahren} n'est pas séparable ici)
\end{enumerate}

\textbf{B.}
\begin{enumerate}
    \item \all{die App benutzen.} (\all{möchten} est modal $\to$ \textbf{pas de \all{zu}} ! La phrase naturelle : \all{Ich möchte die App benutzen.})
    \item \all{ohne fleißig zu sein.} (ou \all{ohne zu lernen.} — sens proche)
    \item \all{um danach spazieren zu gehen.} (but)
    \item \all{anstatt das Buch zu lesen.} (substitution)
\end{enumerate}
\end{corrige}

\courssub{Production guidée : mini-texte avec groupes infinitifs}

\begin{methode}[Rédiger un court paragraphe (5–6 phrases) sur « Lernen mit digitalen Medien »]
\textbf{Consigne :} Écrivez un paragraphe cohérent (niveau B1/B2) en utilisant \textbf{au moins 5 groupes infinitifs différents} parmi : \all{um ... zu}, \all{ohne ... zu}, \all{anstatt ... zu}, \all{zu} après \all{hoffen/versuchen/vergessen/planen}, infinitif sans \all{zu} après modal/\all{lassen}/\all{sehen}.

\textbf{Étapes :}
\begin{enumerate}
    \item Notez 3–4 idées clés (avantages, risques, conseils).
    \item Choisissez les structures infinitivales adaptées à chaque idée.
    \item Rédigez en allemand : verbe conjugué position 2, groupes infinitifs en fin de phrase (ou position 1 pour \all{um ... zu}).
    \item Vérifiez : virgules ? \all{zu} bien placé (verbes séparables) ? Pas de \all{zu} après modaux ?
    \item Relisez à voix haute.
\end{enumerate}
\end{methode}

\begin{exemplebox}[Modèle de production attendue]
\all{Digitale Medien bieten viele Chancen. \textbf{Um besser zu lernen}, nutzen Schüler Apps und Lernprogramme. \textbf{Ohne selbst zu denken}, kopieren manche nur die Lösungen. \textbf{Anstatt nur zu konsumieren}, sollte man aktiv mitmachen. Ich \textbf{hoffe}, dass die KI \textbf{helfen wird}, die Aufgaben \textbf{anzupassen}. Man \textbf{muss} aber \textbf{aufpassen}, \textbf{nicht alles zu glauben, ohne den Lehrer zu fragen}.}
\end{exemplebox}

\courssub{Texte d'entraînement complet : « Die digitale Klasse »}

\begin{textebox}[Die digitale Klasse — Texte original (env. 230 mots)]
An vielen Schulen in Deutschland, Österreich und der Schweiz gehört der digitale Unterricht heute zum Alltag. Lehrer planen, Tablets und Lern-Apps fest in den Stundenplan einzubauen. Die Programme passen die Aufgaben automatisch an das Niveau jedes Schülers an, um niemanden zu überfordern oder zu unterfordern. Wer die Übungen ohne Hilfe zu lösen versucht, bekommt sofort ein Feedback. Das motiviert viele Jugendliche, weil sie ihren Fortschritt direkt sehen.

Doch es gibt auch Risiken. Manche Schüler spielen anstatt zu lernen oder chatten, ohne dem Lehrer zuzuhören. Manche vergessen, die Geräte aufzuladen, oder lassen sich von Benachrichtigungen ablenken. Experten warnen: „Wer nur am Bildschirm liest, ohne handschriftlich zu schreiben, vergisst schneller.“ Auch der Datenschutz ist ein Thema: Schulen müssen sicherstellen, dass keine sensiblen Daten an Dritte weitergegeben werden.

Trotzdem hoffen Pädagogen, die digitale Bildung weiter zu verbessern. Sie beginnen, KI-Tools einzusetzen, um individuelle Lernpfade zu erstellen. Dabei darf der Mensch nicht fehlen: Ein Algorithmus kann nicht lächeln, nicht ermutigen und nicht verstehen, warum ein Kind traurig ist. Deshalb sollen digitale Werkzeuge den Lehrer unterstützen, nicht ersetzen.
\end{textebox}

\begin{definition}[Glossaire — Mots difficiles du texte]
\begin{itemize}
    \item \all{der Stundenplan, die Stundenpläne} — emploi du temps
    \item \all{einbauen} (trennbar: \all{baut ein}) — intégrer, incorporer
    \item \all{automatisch} — automatiquement
    \item \all{überfordern / unterfordern} — surcharger / sous-charger (niveau trop haut / trop bas)
    \item \all{das Feedback, -s} — retour, rétroaction
    \item \all{der Fortschritt, -e} — progrès
    \item \all{chatten} — tchatter, discuter en ligne
    \item \all{zuhören} (trennbar: \all{hört zu}) — écouter
    \item \all{aufladen} (trennbar: \all{lädt auf}) — recharger (batterie)
    \item \all{abgelenkt sein} — être distrait
    \item \all{handschriftlich} — manuscrit, à la main
    \item \all{der Datenschutz} — protection des données
    \item \all{sicherstellen} (trennbar: \all{stellt sicher}) — garantir, s'assurer que
    \item \all{sensibel} — sensible
    \item \all{weitergeben} (trennbar: \all{gibt weiter}) — transmettre
    \item \all{der Lernpfad, die Lernpfade} — parcours d'apprentissage
    \item \all{ermutigen} — encourager
    \item \all{unterstützen} — soutenir, appuyer
    \item \all{ersetzen} — remplacer
\end{itemize}
\end{definition}

\begin{exercice}[Compréhension et grammaire sur « Die digitale Klasse »]
\textbf{A. Compréhension (répondez en allemand, phrases complètes).}
\begin{enumerate}
    \item Was machen die Lernprogramme automatisch? (Utilisez \all{um ... zu} dans votre réponse.)
    \item Was machen manche Schüler anstatt zu lernen?
    \item Wovor warnen Experten? (Deux points.)
    \item Was hoffen Pädagogen?
    \item Was kann ein Algorithmus nicht?
\end{enumerate}

\textbf{B. Grammaire : Repérage dans le texte.}
Soulignez ou citez dans le texte :
\begin{enumerate}
    \item 3 groupes \all{um ... zu} (but).
    \item 2 groupes \all{ohne ... zu} (négation).
    \item 1 groupe \all{anstatt ... zu} (substitution).
    \item 3 verbes appelant \all{zu} + infinitif (\all{planen, versuchen, hoffen, beginnen, vergessen...}).
    \item 2 verbes modaux + infinitif sans \all{zu}.
    \item 1 verbe séparable dans un groupe infinitif (\all{zu} inséré).
\end{enumerate}

\textbf{C. Transformation.}
Reformulez les phrases suivantes en utilisant un groupe infinitif (le sujet reste le même).
\begin{enumerate}
    \item Die Programme passen die Aufgaben an, weil sie niemanden überfordern wollen.
    \item Die Schüler chatten, obwohl sie dem Lehrer zuhören sollen.
    \item Ich habe vor, dass ich nächstes Jahr Informatik studiere.
    \item Man kann die Daten schützen, wenn man sie nicht weitergibt.
\end{enumerate}
\end{exercice}

\begin{corrige}[Exercice : « Die digitale Klasse »]
\textbf{A. Compréhension (propositions de réponses).}
\begin{enumerate}
    \item \all{Die Programme passen die Aufgaben automatisch an, um niemanden zu überfordern oder zu unterfordern.}
    \item \all{Manche Schüler spielen oder chatten, anstatt zu lernen.}
    \item \all{Experten warnen: Wer nur am Bildschirm liest, ohne handschriftlich zu schreiben, vergisst schneller. Auch der Datenschutz ist ein Risiko.}
    \item \all{Pädagogen hoffen, die digitale Bildung weiter zu verbessern.}
    \item \all{Ein Algorithmus kann nicht lächeln, nicht ermutigen und nicht verstehen, warum ein Kind traurig ist.}
\end{enumerate}

\textbf{B. Repérage dans le texte.}
\begin{enumerate}
    \item \all{um niemanden zu überfordern oder zu unterfordern} ; \all{um individuelle Lernpfade zu erstellen} ; \all{um den Lehrer zu unterstützen} (aussi \all{um niemanden zu überfordern} compte comme un seul groupe avec deux infinitifs coordonnés).
    \item \all{ohne dem Lehrer zuzuhören} ; \all{ohne handschriftlich zu schreiben}.
    \item \all{anstatt zu lernen}.
    \item \all{planen, ... einzubauen} ; \all{versuchen, ... zu lösen} ; \all{hoffen, ... zu verbessern} ; \all{beginnen, ... einzusetzen} ; \all{vergessen, ... aufzuladen} — (5 au total, en citer 3 suffit).
    \item \all{muss ... unterstützen} ; \all{kann nicht lächeln / nicht ermutigen / nicht verstehen} (modaux \all{müssen, können}).
    \item \all{anzupassen} (de \all{anpassen}) ; \all{einzubauen} (de \all{einbauen}) ; \all{aufzuladen} (de \all{aufladen}) ; \all{weiterzugeben} (de \all{weitergeben}) ; \all{zuzuhören} (de \all{zuhören}) — (en citer 1 suffit).
\end{enumerate}

\textbf{C. Transformation.}
\begin{enumerate}
    \item \all{Die Programme passen die Aufgaben an, um niemanden zu überfordern.} (déjà dans le texte)
    \item \all{Die Schüler chatten, ohne dem Lehrer zuzuhören.} (ou \all{... anstatt dem Lehrer zuzuhören.})
    \item \all{Ich habe vor, nächstes Jahr Informatik zu studieren.}
    \item \all{Man kann die Daten schützen, ohne sie weiterzugeben.}
\end{enumerate}
\end{corrige}

\begin{savaistu}[Le point « garanti » au Bac — Thème : \texorpdfstring{\og\ pour + infinitif\ \fg{} $\to$ \all{um ... zu}}{pour + infinitif -> um ... zu}]
Dans les épreuves de \textbf{thème} (français $\to$ allemand) du Baccalauréat camerounais (série A), la traduction de \og\ \textbf{pour + infinitif}\ \fg{} par \textbf{\all{um ... zu + Infinitiv}} est un classique quasi systématique.
\begin{itemize}
    \item \og\ Il vient \textbf{pour apprendre} l'allemand.\ \fg{} $\to$ \all{Er kommt, \textbf{um Deutsch zu lernen}.}
    \item \og\ Nous travaillons dur \textbf{pour réussir}.\ \fg{} $\to$ \all{Wir arbeiten hart, \textbf{um zu bestehen}.}
    \item \og\ \textbf{Pour comprendre}, il faut écouter.\ \fg{} $\to$ \all{\textbf{Um zu verstehen}, muss man zuhören.}
\end{itemize}
\textbf{Check-list mentale avant de rendre sa copie :}
\begin{enumerate}
    \item Ai-je mis \all{um} au début du groupe ?
    \item Ai-je mis \all{zu} devant l'infinitif (à la fin du groupe) ?
    \item Ai-je mis la \textbf{virgule} devant \all{um} (ou après le groupe s'il est en position 1) ?
    \item Si le verbe est séparable : \all{zu} est-il \textbf{entre la particule et le radical} ? (\all{anzupassen}, pas \all{*anpassen zu})
    \item Le sujet du groupe infinitif est-il bien le même que celui de la principale ?
\end{enumerate}
\emph{Maîtrisez cette structure : elle rapporte des points faciles à chaque fois !}
\end{savaistu}

\begin{experience}[Activité orale : « KI-Debatte — Für und Wider »]
\textbf{Objectif :} Produire des arguments structurés avec groupes infinitifs (\all{um ... zu}, \all{ohne ... zu}, \all{anstatt ... zu}, \all{zu} + infinitif).

\textbf{Déroulement (15 min) :}
\begin{enumerate}
    \item \textbf{Préparation (5 min) :} Par paires. Chaque duo tire un carton : \all{PRO KI im Unterricht} ou \all{CONTRA KI im Unterricht}. Noter 4 arguments en utilisant \textbf{au moins 3 structures infinitivales différentes} par argument.
    \item \textbf{Débat (7 min) :} Les duos \all{PRO} et \all{CONTRA} s'affrontent. Un argument = une phrase complète. L'autre duo doit répondre en reprenant la structure (ex. : \all{Du sagst: „Um Zeit zu sparen, nutzt man KI.“ — Aber: „Anstatt Zeit zu sparen, verliert man den Kontakt zur Klasse.“})
    \item \textbf{Synthèse (3 min) :} La classe vote les 3 meilleurs arguments (forme + fond). Le prof note au tableau les structures utilisées.
\end{enumerate}

\textbf{Exemples d'arguments attendus :}
\begin{itemize}
    \item \all{Um individuell zu fördern, braucht man KI-Tools.}
    \item \all{Ohne digitale Medien sind wir nicht zukunftsfähig.}
    \item \all{Anstatt selbst zu denken, kopieren Schüler die KI-Antworten.}
    \item \all{Wir versuchen, den Datenschutz zu garantieren.}
    \item \all{Lehrer müssen lernen, die Tools richtig einzusetzen.}
\end{itemize}
\end{experience}

\coursec{Méthode du commentaire et production guidée}

Dans la section précédente, tu as appris à construire des groupes infinitifs avec \all{um ... zu}, \all{ohne ... zu} et \all{anstatt ... zu}. Ces structures, combinées au \cle{Futur I}, sont exactement les outils grammaticaux dont tu auras besoin pour réussir l'épreuve reine du Bac d'allemand : le \cle{commentaire de texte} suivi d'une courte production personnelle. Cette dernière section te donne la méthode complète, pas à pas, avec un modèle rédigé et une grille d'auto-évaluation.

\courssub{Le commentaire de texte au Bac : la méthode en trois parties}

\begin{methode}[Comment réussir un commentaire de texte]
\begin{enumerate}
\item \textbf{Lire le texte deux fois.} La première lecture donne le sens global (\all{das Thema}) ; la deuxième permet de souligner 2 ou 3 idées principales (\all{die Hauptideen}).
\item \textbf{Rédiger l'introduction} (\all{die Einleitung}) : présenter le texte en 1 ou 2 phrases — le thème, le type de texte, éventuellement la source. Utiliser : \all{Der Text handelt von ...} (« Le texte traite de ... »), \all{In dem Text geht es um ...} (« Dans ce texte, il s'agit de ... »).
\item \textbf{Rédiger le développement} (\all{der Hauptteil}) : dégager les idées principales et les arguments de l'auteur, \emph{avec tes propres mots}. Utiliser : \all{Der Autor erklärt, dass ...}, \all{Der Autor ist der Meinung, dass ...} (« L'auteur est d'avis que ... »), \all{Er meint, dass ...}.
\item \textbf{Donner ton avis personnel} (\all{die eigene Meinung}) : c'est la partie la plus valorisée ! Utiliser : \all{Meiner Meinung nach ...}, \all{Ich finde, dass ...}, \all{Ich meine, dass ...}, et justifier avec \all{weil} ou \all{denn}.
\item \textbf{Rédiger la conclusion} (\all{der Schluss}) : faire un bilan en une phrase (\all{Zusammenfassend kann man sagen, dass ...} — « En résumé, on peut dire que ... ») et éventuellement ouvrir sur une question ou une prévision au Futur I.
\end{enumerate}
\end{methode}

\begin{attention}[L'erreur capitale : la traduction mot à mot]
Ne traduis \textbf{jamais} le texte phrase par phrase : ce n'est pas un commentaire, c'est une version, et elle ne rapporte presque aucun point en expression écrite. Le correcteur attend que tu \emph{reformules} les idées avec ton propre vocabulaire. Mauvais réflexe : recopier \all{„Lernprogramme passen die Aufgaben an das Niveau an.“} Bon réflexe : reformuler — \all{Der Autor erklärt, dass der Computer für jeden Schüler andere Aufgaben findet.} (« L'auteur explique que l'ordinateur trouve des exercices différents pour chaque élève. »)
\end{attention}

\begin{popfigure}
\begin{tikzpicture}[
  node distance=0.55cm,
  box/.style={draw=popblue, fill=popblueL, rounded corners=3pt, align=center, text width=5.6cm, inner sep=5pt, font=\small},
  arr/.style={-{Stealth[length=2.5mm]}, thick, popdark}
]
\node[box, align=center] (intro){\textbf{Einleitung}\\ Der Text handelt von ...};
\node[box, below=of intro, align=center] (analyse){\textbf{Analyse}\\ 2--3 Hauptideen + Argumente des Autors\\ \emph{Der Autor erklärt / meint, dass ...}};
\node[box, below=of analyse, fill=poporangeL, draw=poporange, align=center] (meinung){\textbf{Eigene Meinung}\\ Meiner Meinung nach ... / Ich finde, dass ...\\ + Begründung mit \emph{weil / denn}};
\node[box, below=of meinung, fill=popgreenL, draw=popgreen, align=center] (schluss){\textbf{Schluss}\\ Zusammenfassend ... + Ausblick (Futur I)};
\draw[arr] (intro) -- (analyse);
\draw[arr] (analyse) -- (meinung);
\draw[arr] (meinung) -- (schluss);
\end{tikzpicture}
\legende{Organigramme du commentaire de texte : quatre étapes obligatoires}
\end{popfigure}

\courssub{Les expressions pour introduire et pour argumenter}

Pour structurer ton commentaire, tu dois impérativement connaître par cœur un stock de \cle{Redemittel} (expressions toutes faites). Les voici, classées par fonction :

\begin{center}
\begin{tabularx}{\linewidth}{|X|X|}
\hline
\rowcolor{popblueL} \textbf{Deutsch} & \textbf{Français} \\
\hline
Der Text handelt von der künstlichen Intelligenz. & Le texte traite de l'intelligence artificielle. \\
\hline
In dem Text geht es um digitale Medien in der Schule. & Dans ce texte, il s'agit des médias numériques à l'école. \\
\hline
Der Autor erklärt, dass Lernprogramme sehr nützlich sind. & L'auteur explique que les programmes d'apprentissage sont très utiles. \\
\hline
Der Autor ist der Meinung, dass die KI den Lehrer nicht ersetzen wird. & L'auteur est d'avis que l'IA ne remplacera pas le professeur. \\
\hline
Meiner Meinung nach ist \textbf{die} Technik eine große Hilfe. & À mon avis, la technique est une grande aide. \\
\hline
Ich finde, dass Tablets den Unterricht interessanter machen. & Je trouve que les tablettes rendent le cours plus intéressant. \\
\hline
\end{tabularx}
\end{center}

Pour argumenter, la structure la plus efficace est le couple \all{einerseits ... andererseits ...} (« d'un côté ... de l'autre ... »), idéal pour présenter les avantages (\all{die Vorteile}) et les inconvénients (\all{die Nachteile}) :

\begin{exemplebox}[Argumenter en allemand]
\all{Einerseits hilft die Technik den Schülern, andererseits braucht man Menschen.} « D'un côté la technique aide les élèves, de l'autre on a besoin d'êtres humains. »

\all{Ein Vorteil ist, dass die App die Fehler sofort korrigiert.} « Un avantage est que l'application corrige les erreurs immédiatement. »

\all{Ein Nachteil ist, dass nicht alle Familien ein Tablet kaufen können.} « Un inconvénient est que toutes les familles ne peuvent pas acheter une tablette. »

\all{Meiner Meinung nach wird die KI den Lehrer nie ersetzen.} « À mon avis, l'IA ne remplacera jamais le professeur. »
\end{exemplebox}

\begin{attention}[La place du verbe après « Meiner Meinung nach »]
\all{Meiner Meinung nach} occupe la \textbf{première position} de la phrase : le verbe conjugué vient donc \textbf{immédiatement après}, en position 2. On dit : \all{Meiner Meinung nach \textbf{ist} die KI nützlich.} et jamais \emph{« Meiner Meinung nach die KI ist nützlich »}. Même règle avec \all{einerseits} et \all{andererseits} : \all{Einerseits \textbf{hilft} die Technik, andererseits \textbf{braucht} man Menschen.} C'est l'erreur d'inversion la plus fréquente des francophones, car en français le sujet suit normalement l'expression (« à mon avis, l'IA est utile »).
\end{attention}

\courssub{Commentaire modèle rédigé collectivement}

Voici un commentaire modèle construit sur le texte support \all{Lernen mit der KI}. Observe comment chaque phrase correspond à une étape de l'organigramme, et repère le \cle{Futur I} et le \cle{groupe infinitif} exigés par la consigne.

\begin{textebox}[Kommentar — Modell]
Der Text handelt von der künstlichen Intelligenz in der Schule. Der Autor erklärt, dass Lernprogramme die Aufgaben anpassen und Fehler sofort korrigieren. Er meint aber, dass die KI den Lehrer nicht ersetzen wird. Ich finde, dass \textbf{die} Technik sehr nützlich ist, um schneller zu lernen. Aber die Schule der Zukunft wird auch menschliche Beziehungen brauchen.
\end{textebox}

\textbf{Analyse du modèle :}
\begin{itemize}
\item Phrase 1 = \textbf{introduction} avec \all{Der Text handelt von ...} ;
\item phrases 2 et 3 = \textbf{développement} : reformulation des idées de l'auteur (\all{erklärt, dass ...} / \all{meint, dass ...}) ;
\item phrase 4 = \textbf{avis personnel} avec \all{Ich finde, dass ...} + groupe infinitif \all{um schneller zu lernen} ;
\item phrase 5 = \textbf{conclusion} avec une prévision au Futur I : \all{wird ... brauchen}.
\end{itemize}

\begin{savaistu}[L'avis personnel rapporte des points]
Au Bac d'allemand, l'expression personnelle est toujours valorisée : un avis justifié avec \all{weil} ou \all{denn} rapporte plus de points qu'une simple paraphrase du texte. Écris par exemple : \all{Ich meine, dass Tablets wichtig sind, weil viele Schüler in Kamerun keinen Computer zu Hause haben.} (« Je pense que les tablettes sont importantes, parce que beaucoup d'élèves au Cameroun n'ont pas d'ordinateur à la maison. ») Une phrase comme celle-ci montre au correcteur ta grammaire (subordonnée avec verbe final), ton vocabulaire et ta capacité de réflexion.
\end{savaistu}

\courssub{Production guidée : « Die Schule der Zukunft in Kamerun »}

\begin{methode}[Rédiger un court paragraphe argumenté (6--8 lignes)]
\begin{enumerate}
\item \textbf{Planifier} : note 2 avantages, 1 inconvénient et 1 prévision avant d'écrire.
\item \textbf{Commencer} par une phrase d'introduction : \all{Wie wird die Schule der Zukunft in Kamerun aussehen?} (« À quoi ressemblera l'école du futur au Cameroun ? »)
\item \textbf{Présenter les avantages} avec \all{Ein Vorteil ist, dass ...} et \all{Außerdem ...} (« De plus ... »).
\item \textbf{Présenter l'inconvénient} avec \all{Aber ein Nachteil ist, dass ...} ou \all{Andererseits ...}.
\item \textbf{Terminer} par une prévision au Futur I (\all{wird ... werden / haben / sein}) et, si possible, un groupe infinitif avec \all{um ... zu}.
\item \textbf{Relire} : vérifier la place du verbe (position 2, verbe final dans les subordonnées), la majuscule aux noms, les articles.
\end{enumerate}
\end{methode}

\begin{exoresolu}[Production guidée — sujet de type Bac]
\textbf{Énoncé.} Rédige un paragraphe de 6 à 8 lignes intitulé \all{Die Schule der Zukunft in Kamerun}. Tu dois inclure : 2 avantages des outils numériques, 1 inconvénient, 1 prévision au Futur I et 1 groupe infinitif avec \all{um ... zu}.
\tcblower
\textbf{Solution détaillée.} On suit le plan : question d'introduction, deux avantages, un inconvénient, prévision au Futur I, groupe infinitif.

\all{Wie wird die Schule der Zukunft in Kamerun aussehen? Ein Vorteil ist, dass die Schüler mit Tablets und Apps schneller lernen können. Außerdem korrigiert das Lernprogramm die Fehler sofort, und das hilft besonders den Schülern in den Dörfern. Andererseits ist ein Nachteil, dass der Strom und das Internet nicht überall funktionieren. Meiner Meinung nach werden die Schulen in Yaoundé und Douala in zehn Jahren ganz digital sein. Aber die KI wird den Lehrer nie ersetzen, denn die Schüler brauchen Menschen, um motiviert zu bleiben.}

« À quoi ressemblera l'école du futur au Cameroun ? Un avantage est que les élèves peuvent apprendre plus vite avec des tablettes et des applications. De plus, le programme d'apprentissage corrige les erreurs immédiatement, et cela aide surtout les élèves des villages. D'un autre côté, un inconvénient est que l'électricité et Internet ne fonctionnent pas partout. À mon avis, les écoles de Yaoundé et Douala seront entièrement numériques dans dix ans. Mais l'IA ne remplacera jamais le professeur, car les élèves ont besoin d'êtres humains pour rester motivés. »

\textbf{Vérification des contraintes :} 2 avantages (\all{Tablets und Apps} ; \all{korrigiert die Fehler sofort}) ; 1 inconvénient (\all{Strom und Internet}) ; Futur I (\all{werden ... sein}, \all{wird ... ersetzen}) ; groupe infinitif avec \all{um ... zu} (\all{um motiviert zu bleiben}) ; verbe en position 2 après \all{Meiner Meinung nach} et \all{Andererseits}.
\end{exoresolu}

\courssub{Expression orale guidée : présenter et débattre}

\begin{experience}[Présentation orale et questions-réponses]
\textbf{Déroulement (10 minutes) :}
\begin{enumerate}
\item Chaque élève lit son paragraphe \all{Die Schule der Zukunft in Kamerun} devant la classe, à voix haute et sans se presser.
\item Les camarades posent une ou deux questions avec les structures apprises : \all{Warum meinst du, dass die KI den Lehrer nicht ersetzen wird?} (« Pourquoi penses-tu que l'IA ne remplacera pas le professeur ? »), \all{Was ist für dich der größte Vorteil?} (« Quel est pour toi le plus grand avantage ? »)
\item L'élève interrogé répond en phrase complète : \all{Ich meine das, weil ...} / \all{Der größte Vorteil ist, dass ...}.
\item La classe vote : qui a donné l'argument le plus convaincant (\all{das beste Argument}) ?
\end{enumerate}
\textbf{Conseil de prononciation :} \all{die Zukunft} se prononce « tsou-kounft », \all{die künstliche Intelligenz} « kiunst-li-che in-té-li-guèn-ts ».
\end{experience}

\courssub{Auto-évaluation : la grille de contrôle}

Avant de rendre ta copie ou de passer à l'oral, vérifie chaque point de cette grille. Mets une croix dans la colonne de droite seulement si le critère est vraiment rempli.

\begin{center}
\begin{tabularx}{\linewidth}{|X|l|}
\hline
\rowcolor{popblueL} \textbf{Critère} & \textbf{Oui} \\
\hline
J'ai utilisé au moins 3 mots du vocabulaire du thème (die KI, das Lernprogramm, das Tablet, die App, der Computer, lernen). & \\
\hline
J'ai écrit au moins une phrase au Futur I correct (werden + infinitif à la fin). & \\
\hline
J'ai construit un groupe infinitif correct (um ... zu + infinitif à la fin). & \\
\hline
Mes phrases sont complètes : sujet + verbe conjugué en position 2. & \\
\hline
Après \all{Meiner Meinung nach} et \all{einerseits}, le verbe est bien en position 2. & \\
\hline
J'ai justifié mon avis avec \all{weil} ou \all{denn} (verbe à la fin après \all{weil}). & \\
\hline
Tous les noms ont une majuscule et un article correct. & \\
\hline
\end{tabularx}
\end{center}

\begin{exercice}[À toi de jouer !]
Rédige ton propre paragraphe \all{Die Schule der Zukunft in Kamerun} (6--8 lignes) en respectant les contraintes : 2 avantages, 1 inconvénient, 1 prévision au Futur I, 1 groupe infinitif avec \all{um ... zu}. Puis auto-évalue-toi avec la grille ci-dessus avant de présenter ton texte à la classe.
\end{exercice}

\begin{corrige}[Éléments de correction — production libre]
La production est personnelle, mais le correcteur vérifiera les points suivants : présence d'au moins 3 mots du lexique du chapitre ; un Futur I correctement formé (\all{werden} conjugué en position 2 + infinitif en fin de phrase : \all{Die Schulen werden mehr Tablets benutzen.}) ; un groupe infinitif avec \all{um ... zu} sans sujet répété (\all{..., um besser zu lernen.}) ; l'inversion du verbe après \all{Meiner Meinung nach} ; la justification de l'avis avec \all{weil} + verbe final. Un paragraphe qui remplit ces cinq critères obtient la totalité des points de grammaire et de vocabulaire.
\end{corrige}

\begin{aretenir}[L'essentiel de la section]
\begin{itemize}
\item Un commentaire de texte suit toujours le plan : \textbf{Einleitung} (\all{Der Text handelt von ...}) → \textbf{Analyse} (idées et arguments reformulés) → \textbf{eigene Meinung} (\all{Meiner Meinung nach ...} + \all{weil}) → \textbf{Schluss} (bilan + prévision au Futur I).
\item On ne traduit jamais le texte : on \emph{reformule} les idées principales.
\item Après \all{Meiner Meinung nach}, \all{einerseits}, \all{andererseits} : verbe conjugué en \textbf{position 2}.
\item Pour argumenter : \all{Ein Vorteil ist, dass ...} / \all{Ein Nachteil ist, dass ...} ; pour nuancer : \all{einerseits ... andererseits ...}.
\item Un avis justifié avec \all{weil} ou \all{denn} rapporte plus de points qu'une paraphrase.
\end{itemize}
\end{aretenir}

\newpage

\coursec{Activité d'intégration}

\coursec{AL19 : Les outils numériques pédagogiques et l'intelligence artificielle}

\courssub{Texte support : Lernen mit der KI}

Le texte suivant est une annonce authentique, rédigée pour le club d'allemand d'un lycée bilingue de Yaoundé. Il sert de déclencheur à la leçon : il présente la situation de communication, introduit le vocabulaire thématique et contextualise les points de grammaire (Futur I, groupes infinitifs).

\begin{textebox}[Ankündigung des Deutschclubs]
\textbf{Große Debatte: Brauchen wir KI in unseren Schulen in Kamerun?}

Liebe Schülerinnen und Schüler,

unser Gymnasium plant, bald Tablets anzuschaffen und ein Lernprogramm mit künstlicher Intelligenz zu testen. Ist das ein Fortschritt für unser Lernen? Die Befürworter argumentieren: „Die KI wird jedem Schüler helfen, schneller und besser zu lernen.“ Die Gegner warnen: „Die Schüler werden die Apps nur benutzen, um Hausaufgaben zu kopieren, ohne selbst zu denken.“

Kommt am Freitag um 14 Uhr in die Aula! Zwei Teams werden diskutieren: die \textbf{Befürworter} und die \textbf{Gegner}. Jeder Schüler darf seine Meinung äußern. Am Ende schreibt jeder einen kurzen Text: \textbf{„Die Schule der Zukunft in Kamerun“}.

Der Deutschclub
\end{textebox}

\textbf{Glossaire du texte :}
\begin{itemize}
\item \all{die Ankündigung, -en} : l'annonce, l'avis
\item \all{das Gymnasium, -en} : le lycée (d'enseignement général)
\item \all{anschaffen (schafft an, schaffte an, hat angeschafft)} : se procurer, acheter
\item \all{das Lernprogramm, -e} : le logiciel éducatif, le programme d'apprentissage
\item \all{die künstliche Intelligenz (die KI)} : l'intelligence artificielle (l'IA)
\item \all{testen} : tester, essayer
\item \all{der Befürworter, - / die Befürworterin, -nen} : le partisan / la partisane
\item \all{der Gegner, - / die Gegnerin, -nen} : l'adversaire, l'opposant
\item \all{argumentieren} : argumenter
\item \all{warnten (warnen)} : avertir, mettre en garde
\item \all{benutzen} : utiliser, se servir de
\item \all{kopieren} : copier
\item \all{ohne ... zu} : sans ... (infinitif)
\item \all{die Aula, -en} : la salle de fêtes, l'amphithéâtre
\item \all{äußern (sich äußern)} : s'exprimer, donner son avis
\item \all{die Meinung, -en} : l'opinion
\item \all{die Zukunft} : l'avenir, le futur
\end{itemize}

\courssub{Compréhension du texte (Global und Detailverständnis)}

\begin{exercice}[Questions de compréhension]
Réponds en allemand, par des phrases complètes, aux questions suivantes :
\begin{enumerate}
\item Was plant das Gymnasium in Yaoundé anzuschaffen?
\item Welche zwei Meinungen gibt es zur Einführung der KI? Nenne je ein Argument der Befürworter und der Gegner.
\item Wann und wo findet die Debatte statt?
\item Was ist die Aufgabe für jeden Schüler am Ende der Debatte?
\end{enumerate}
\end{exercice}

\begin{corrige}[Exercice 1 : Compréhension]
\begin{enumerate}
\item \all{Das Gymnasium plant, Tablets anzuschaffen und ein Lernprogramm mit künstlicher Intelligenz zu testen.}
\item \all{Die Befürworter sagen: Die KI wird jedem Schüler helfen, schneller und besser zu lernen. Die Gegner warnen: Die Schüler werden die Apps nur benutzen, um Hausaufgaben zu kopieren, ohne selbst zu denken.}
\item \all{Die Debatte findet am Freitag um 14 Uhr in der Aula statt.}
\item \all{Jeder Schüler muss einen kurzen Text schreiben: „Die Schule der Zukunft in Kamerun“.}
\end{enumerate}
\end{corrige}

\courssub{Lexique thématique : Digitale Medien und KI}

Le vocabulaire suivant est indispensable pour comprendre le texte, participer au débat et rédiger la production finale. Apprends les noms \textbf{avec leur article et leur pluriel}, les verbes \textbf{avec leur régime} (cas) et les adjectifs utiles pour argumenter.

\begin{popfigure}
\begin{tikzpicture}[
  mindmap, grow cyclic, every node/.style=concept, concept color=popblue,
  level 1/.append style={level distance=4.5cm, sibling angle=90, concept color=popblue!70},
  level 2/.append style={level distance=3.2cm, sibling angle=45, concept color=poporange!70},
  text=white, font=\small\bfseries
]
\node[align=center]{Digitale\\Schule}
  child [concept color=popgreen] { node {Hardware}
    child { node {das Tablet, -s} }
    child { node {der Computer, -} }
    child { node {das Smartphone, -s} }
  }
  child [concept color=poppurple] { node {Software \& KI}
    child { node {die App, -s} }
    child { node {das Lernprogramm, -e} }
    child { node {die künstliche Intelligenz} }
    child { node {der Algorithmus, -men} }
  }
  child [concept color=poppink] { node {Verben \& Aktionen}
    child { node {lernen / üben / wiederholen} }
    child { node {helfen (Dat.)} }
    child { node {benutzen / bedienen} }
    child { node {kopieren / abschreiben} }
    child { node {denken / reflektieren} }
  }
  child [concept color=popgold] { node {Adjektive \& Meinung}
    child { node {nützlich / sinnvoll} }
    child { node {gefährlich / riskant} }
    child { node {teuer / günstig} }
    child { node {modern / veraltet} }
    child { node {unverzichtbar} }
  };
\end{tikzpicture}
\legende{Carte mentale : Le vocabulaire de l'école numérique (Hardware, Software, Actions, Opinions).}
\end{popfigure}

\begin{aretenir}[Noms clés (Genre, Pluriel, Sens)]
\begin{center}
\begin{tabularx}{\linewidth}{|>{\bfseries}l|X|X|}\hline
\rowcolor{popblueL} \textbf{Deutsch (Nom, Plural)} & \textbf{Français} & \textbf{Exemple / Contexte} \\ \hline
\all{das Tablet, -s} & la tablette numérique & \all{Ich lerne mit dem Tablet.} \\ \hline
\all{der Computer, -} & l'ordinateur & \all{Der Computer ist neu.} \\ \hline
\all{die App, -s} & l'application (mobile) & \all{Die App hilft mir beim Vokabeln lernen.} \\ \hline
\all{das Lernprogramm, -e} & le logiciel éducatif & \all{Das Lernprogramm kostet Geld.} \\ \hline
\all{die künstliche Intelligenz} & l'intelligence artificielle & \all{Die KI verändert die Schule.} \\ \hline
\all{der Algorithmus, -men} & l'algorithme & \all{Der Algorithmus wählt die Aufgaben.} \\ \hline
\all{die Hausaufgabe, -n} & le devoir (à la maison) & \all{Ich mache meine Hausaufgaben.} \\ \hline
\all{der Schüler, - / die Schülerin, -nen} & l'élève & \all{Die Schüler diskutieren.} \\ \hline
\all{die Aula, -en} & la salle de fêtes & \all{Wir treffen uns in der Aula.} \\ \hline
\all{die Zukunft} & l'avenir & \all{In der Zukunft wird alles digital sein.} \\ \hline
\end{tabularx}
\end{center}
\end{aretenir}

\begin{aretenir}[Verbes utiles et leur régime (Cas)]
\begin{center}
\begin{tabularx}{\linewidth}{|>{\bfseries}l|l|X|}\hline
\rowcolor{popblueL} \textbf{Verbe (Präsens / Präteritum / Perfekt)} & \textbf{Régime} & \textbf{Exemple} \\ \hline
\all{helfen, hilft, half, hat geholfen} & \textbf{+ Dativ} & \all{Die KI hilft \textbf{dem Schüler} / \textbf{den Schülern}.} \\ \hline
\all{benutzen, benutzt, benutzte, hat benutzt} & + Akkusativ & \all{Er benutzt \textbf{das Tablet}.} \\ \hline
\all{bedienen, bedient, bediente, hat bedient} & + Akkusativ & \all{Kannst du den Computer \textbf{bedienen}?} \\ \hline
\all{unterstützen, unterstützt, unterstützte, hat unterstützt} & + Akkusativ & \all{Das Programm unterstützt \textbf{das Lernen}.} \\ \hline
\all{ersetzen, ersetzt, ersetzte, hat ersetzt} & + Akkusativ & \all{Die KI wird den Lehrer nicht \textbf{ersetzen}.} \\ \hline
\all{sich verlassen auf + Akk.} & Réfléchi + Préposition & \all{Wir verlassen uns auf \textbf{die Technik}.} \\ \hline
\end{tabularx}
\end{center}
\end{aretenir}

\begin{attention}[Faux amis et pièges lexicaux]
\begin{itemize}
\item \all{die Tablette} (féminin, plur. \all{-n}) = le comprimé (médicament). Pour la tablette numérique : \textbf{\all{das Tablet}} (neutre, plur. \all{-s}).
\item \all{das Handy} = le téléphone portable (familier). Forme standard : \all{das Mobiltelefon} ou \all{das Smartphone}.
\item \all{der Computer} (masculin) ne prend pas de \textit{-s} au pluriel : \all{die Computer}.
\item \all{helfen} + \textbf{Datif} ! Jamais \all{*helfen den Schüler} (Akkusatif). Pense : \all{Ich helfe \textbf{dir} / \textbf{ihm} / \textbf{den Schülern}}.
\item \all{lernen} s'emploie souvent avec préposition : \all{lernen \textbf{für} die Prüfung} (réviser pour), \all{etwas \textbf{auswendig} lernen} (apprendre par cœur).
\end{itemize}
\end{attention}

\courssub{Grammaire 1 : Le Futur I (das Futur I)}

Le Futur I exprime une action future, une prévision, une hypothèse ou une promesse. C'est le temps de la \textbf{prédiction} et du \textbf{projet}.

\begin{propriete}[Formation du Futur I]
\textbf{Forme :} \all{werden} (conjugué au présent) + \textbf{Infinitif du verbe principal à la fin de la phrase (position finale)}.
\begin{center}
\begin{tabular}{|l|l|l|l|}\hline
\rowcolor{popblueL} \textbf{Personne} & \textbf{werden (Präsens)} & \textbf{Infinitif (ex. lernen)} & \textbf{Traduction} \\ \hline
\all{ich} & \all{werde} & \all{lernen} & \all{ich werde lernen} (j'apprendrai / je vais apprendre) \\ \hline
\all{du} & \all{wirst} & \all{lernen} & \all{du wirst lernen} \\ \hline
\all{er / sie / es} & \all{wird} & \all{lernen} & \all{er wird lernen} \\ \hline
\all{wir} & \all{werden} & \all{lernen} & \all{wir werden lernen} \\ \hline
\all{ihr} & \all{werdet} & \all{lernen} & \all{ihr werdet lernen} \\ \hline
\all{sie / Sie} & \all{werden} & \all{lernen} & \all{sie werden lernen} \\ \hline
\end{tabular}
\end{center}
\end{propriete}

\begin{exemplebox}[Emplois du Futur I avec exemples traduits]
\begin{itemize}
\item \textbf{Prévision / Prédiction objective :} \all{Die KI \textbf{wird} den Unterricht \textbf{verändern}.} « L'IA changera l'enseignement. »
\item \textbf{Hypothèse / Supposition (souvent avec \all{wohl}, \all{vermutlich}, \all{sicher}) :} \all{Die Tablets \textbf{werden} \textbf{wohl} \textbf{teuer} \textbf{sein}.} « Les tablettes seront probablement chères. »
\item \textbf{Promesse / Engagement personnel :} \all{Ich \textbf{werde} die App \textbf{nur} zum Üben \textbf{benutzen}.} « J'utiliserai l'appli seulement pour m'exercer. »
\item \textbf{Projet / Intention future (souvent avec indicateur de temps) :} \all{Wir \textbf{werden} nächstes Jahr \textbf{mit Tablets arbeiten}.} « Nous travaillerons avec des tablettes l'an prochain. »
\end{itemize}
\end{exemplebox}

\begin{attention}[Erreurs fréquentes des francophones — Futur I]
\begin{itemize}
\item \textbf{Infinitif sans \all{zu} :} Après \all{werden}, l'infinitif est \textbf{nu} (sans \all{zu}). \all{*Ich werde zu lernen} $\rightarrow$ \textbf{\all{Ich werde lernen.}}
\item \textbf{Place de l'infinitif :} L'infinitif va \textbf{à la fin de la phrase} (ou de la proposition), après les compléments. \all{*Ich werde lernen Deutsch} $\rightarrow$ \textbf{\all{Ich werde Deutsch lernen.}}
\item \textbf{Ne pas confondre Futur I et Futur II :} Le Futur II (\all{werden + Partizip II + haben/sein}) exprime l'antériorité future (\all{Ich werde gelernt haben}). Non au programme de Terminale LV2, mais à reconnaître en lecture.
\item \textbf{\all{werden} comme verbe principal :} \all{werden} = devenir. \all{Er wird Lehrer.} (Il devient prof). Pas d'infinitif derrière.
\end{itemize}
\end{attention}

\courssub{Grammaire 2 : Les groupes infinitifs (die Infinitivgruppen)}

Les groupes infinitifs (ou propositions infinitives) permettent d'alléger la phrase en remplaçant une proposition subordonnée (\all{dass}, \all{weil}, \all{obwohl}...) par une construction \textbf{infinitif + \all{zu}}. Ils sont obligatoires après certaines prépositions ou expressions.

\begin{propriete}[Structure générale et Ponctuation]
\textbf{Règle d'or :} L'infinitif avec \all{zu} se place \textbf{à la fin du groupe infinitif}. Une \textbf{virgule} sépare toujours le groupe infinitif de la proposition principale (sauf si le groupe est très court et non ambigu, mais au Bac : \textbf{met la virgule}).
\par\medskip
\textbf{Schéma :} \all{Hauptsatz ,} \textbf{\all{um / ohne / anstatt / statt}} \dots \textbf{\all{zu}} \textbf{Infinitif} \textbf{.}
\end{propriete}

\begin{exemplebox}[Les trois types principaux au programme]
\begin{enumerate}
\item \textbf{\all{um ... zu}} = \textbf{pour / afin de} (But, Finalité). Le sujet du groupe infinitif = sujet de la principale.
\all{Ich benutze die App, \textbf{um} besser Deutsch \textbf{zu lernen}.} (J'utilise l'appli pour mieux apprendre l'allemand.)
\item \textbf{\all{ohne ... zu}} = \textbf{sans} + infinitif passé/présent (Manière, absence d'action).
\all{Er kopiert den Text, \textbf{ohne} ihn \textbf{zu verstehen}.} (Il copie le texte sans le comprendre.)
\item \textbf{\all{anstatt ... zu} / \all{statt ... zu}} = \textbf{au lieu de} (Alternative, opposition).
\all{Sie spielt mit dem Tablet, \textbf{anstatt} ihre Hausaufgaben \textbf{zu machen}.} (Elle joue avec la tablette au lieu de faire ses devoirs.)
\end{enumerate}
\end{exemplebox}

\begin{aretenir}[Autres introductions de groupes infinitifs (Reconnaissance)]
\begin{itemize}
\item \all{haben / sein / brauchen + zu + Infinitif} (Modalité passive / nécessité) : \all{Die Hausaufgaben sind \textbf{zu machen}.} (Les devoirs sont à faire / doivent être faits.) \all{Es gibt viel \textbf{zu tun}.} (Il y a beaucoup à faire.)
\item \all{Adjectif + zu + Infinitif} : \all{Das ist leicht \textbf{zu verstehen}.} (C'est facile à comprendre.) \all{Er ist nicht in der Lage, \textbf{zu kommen}.} (Il n'est pas en mesure de venir.)
\item \all{Verbes de perception / mouvement + Infinitif \textbf{sans zu}} : \all{Ich sehe ihn \textbf{kommen}.} \all{Wir gehen \textbf{spazieren}.} (Voir faire / Aller faire).
\end{itemize}
\end{aretenir}

\begin{attention}[Pièges à éviter — Groupes infinitifs]
\begin{itemize}
\item \textbf{Oubli de la virgule :} \all{*Ich lerne um Deutsch zu sprechen.} $\rightarrow$ \textbf{\all{Ich lerne, um Deutsch zu sprechen.}}
\item \textbf{Mauvaise place de \all{zu} :} \all{*Ich versuche, zu die Hausaufgaben machen.} $\rightarrow$ \textbf{\all{Ich versuche, die Hausaufgaben zu machen.}} (\all{zu} immédiatement devant l'infinitif, à la fin).
\item \textbf{Sujet différent :} Si le sujet du groupe infinitif est \textbf{différent} de celui de la principale, on \textbf{ne peut pas} utiliser de groupe infinitif $\rightarrow$ proposition subordonnée (\all{dass}, \all{weil}...). \all{*Ich will, um er lernt.} $\rightarrow$ \textbf{\all{Ich will, dass er lernt.}}
\item \textbf{\all{anstatt} / \all{statt} + Génitif (nom) vs \all{anstatt ... zu} (verbe) :} \all{Anstatt des Lehrers} (Au lieu du prof - nom) vs \all{Anstatt zu lehren} (Au lieu d'enseigner - verbe).
\end{itemize}
\end{attention}

\courssub{Expression guidée : Débat et production écrite « Die Schule der Zukunft in Kamerun »}

Cette activité d'intégration mobilise toutes les compétences de la leçon : compréhension, lexique, grammaire (Futur I, Infinitivgruppen), expression orale et écrite.

\begin{methode}[Conduite de l'activité : Du débat à la rédaction]
\begin{enumerate}
\item \textbf{Préparation individuelle (10 min) :} Lis l'annonce, choisis ton camp (\all{Befürworter} ou \all{Gegner}), note 3 arguments en respectant les contraintes grammaticales (1 Futur I, 1 Groupe infinitif, 3 mots de vocabulaire).
\item \textbf{Mise en commun par équipe (10 min) :} Regroupe-toi avec ton camp. Sélectionnez les 4 meilleurs arguments. Désignez 2 orateurs principaux.
\item \textbf{Débat modéré (15 min) :} L'enseignant modère. Les équipes s'expriment à tour de rôle. Les autres élèves notent les arguments adverses pour la synthèse.
\item \textbf{Rédaction individuelle (15 min) :} Rédige ton paragraphe \all{„Die Schule der Zukunft in Kamerun“} (6--8 lignes). Structure : Sujet $\rightarrow$ 2 arguments (grammaire obligatoire) $\rightarrow$ Avis personnel nuancé $\rightarrow$ Conclusion.
\item \textbf{Relecture croisée / Affichage (5 min) :} Échange ton texte avec un voisin. Vérifie : Majuscules noms ? Virgules groupes infinitifs ? \all{werden} + Infinitif fin ? \all{helfen} + Datif ? Les meilleurs sont lus et affichés.
\end{enumerate}
\end{methode}

\begin{experience}[Grand débat : Brauchen wir KI in unseren Schulen in Kamerun?]
\textbf{Matériel :} Tableau divisé en deux colonnes (\all{Argumente der Befürworter} / \all{Argumente der Gegner}). Cartes vocabulaire (mots-clés de la leçon) distribuées aux élèves.
\par\medskip
\textbf{Déroulement :}
\begin{enumerate}
\item L'enseignant écrit le sujet au tableau : \all{Brauchen wir KI in unseren Schulen in Kamerun?}
\item \textbf{Tour 1 (Arguments) :} Équipe A (Befürworter) présente un argument. Équipe B (Gegner) réagit avec une formule de politesse/désaccord (\all{Ich stimme nicht zu, denn... / Das sehe ich anders, weil...}) puis présente son argument. Alternance.
\item \textbf{Tour 2 (Réactions libres) :} La parole est libre. L'enseignant note au tableau les structures grammaticales cibles entendues (Futur I, \all{um...zu}, \all{ohne...zu}, \all{helfen + Dat}).
\item \textbf{Vote final :} \all{Hand hoch: Wer ist für KI? Wer dagegen?} Comptage en allemand (\all{eins, zwei, drei...}).
\end{enumerate}
\textbf{Formules de prise de parole (à afficher) :}
\all{Ich bin der Meinung, dass ... ; Meiner Meinung nach ... ; Ich stimme dir zu, aber ... ; Das sehe ich anders, denn ... ; Meinst du wirklich, dass ...? ; Ein Argument dafür / dagegen ist ... ; Außerdem ... ; Allerdings ...}
\end{experience}

\begin{savaistu}[Contexte : L'école numérique en Allemagne et au Cameroun]
\begin{itemize}
\item \textbf{Allemagne :} Le \all{DigitalPakt Schule} (2019--2024) a investi 5 milliards d'euros pour équiper les écoles en WLAN, tablettes et tableaux numériques. La \all{Kultusministerkonferenz} (KMK) a publié une stratégie \all{„Bildung in der digitalen Welt“}. L'IA (\all{KI}) est testée pour l'apprentissage adaptatif (ex. \all{Bettermarks}, \all{Anton App}).
\item \textbf{Autriche / Suisse :} Stratégies similaires (\all{digi.komp} en Autriche, \all{Lehrplan 21} en Suisse romande/alémanique). Focus sur les \all{Medienkompetenz} (compétences médiatiques) dès le primaire.
\item \textbf{Cameroun :} Le MINESEC encourage l'intégration des TICE. Projets \all{« École numérique »}, distribution de tablettes dans certains lycées pilotes (Yaoundé, Douala). Défis : connexion internet instable, coût de la maintenance, formation des enseignants, fracture numérique ville/campagne.
\item \textbf{Débat éthique :} Protection des données (\all{Datenschutz / DSGVO}), biais algorithmiques, dépendance technologique vs autonomie de l'élève, rôle irremplaçable de l'enseignant (\all{der Lehrer / die Lehrerin}) comme guide et éducateur.
\end{itemize}
\end{savaistu}

\begin{exoresolu}[Paragraphe modèle : „Die Schule der Zukunft in Kamerun“]
\textbf{Consigne :} Rédige un paragraphe de 6 à 8 lignes donnant ton avis personnel sur l'IA à l'école. Contraintes : au moins \textbf{un Futur I}, \textbf{un groupe infinitif} (\all{um/ohne/anstatt ... zu}), \textbf{trois mots du lexique} de la leçon.
\tcblower
\textbf{Production modèle :}

\all{In unserem Gymnasium in Yaoundé diskutieren wir derzeit intensiv über den Einsatz von künstlicher Intelligenz im Unterricht. Ich bin der Meinung, dass die KI den Schülern in Kamerun große Chancen bieten wird. Mit einem guten Lernprogramm kann jeder Schüler zu Hause selbstständig üben, \textbf{um} seine Noten \textbf{zu verbessern}. Die Lehrer \textbf{werden} mehr Zeit haben, \textbf{um} schwache Schüler individuell \textbf{zu fördern}. Aber wir dürfen die Technik nicht blind vertrauen: Die Schüler dürfen die Apps nicht \textbf{ohne} kritisch \textbf{zu denken} \textbf{benutzen}. \textbf{Anstatt} nur Antworten \textbf{zu kopieren}, müssen sie die Lösungswege verstehen. Die Schule der Zukunft \textbf{wird} also Tablets und Computer \textbf{haben}, aber der Mensch – der Lehrer und der Schüler – \textbf{wird} immer der Mittelpunkt des Lernens \textbf{bleiben}.}

\textbf{Traduction :}
« Dans notre lycée de Yaoundé, nous discutons actuellement intensément de l'usage de l'intelligence artificielle en classe. Je suis d'avis que l'IA offrira de grandes chances aux élèves du Cameroun. Avec un bon logiciel éducatif, chaque élève peut s'exercer de manière autonome à la maison \textbf{afin d'améliorer} ses notes. Les professeurs \textbf{auront} plus de temps \textbf{pour} soutenir individuellement les élèves faibles. Mais nous ne devons pas faire aveuglément confiance à la technique : les élèves ne doivent pas utiliser les applications \textbf{sans réfléchir} de manière critique. \textbf{Au lieu de} seulement copier les réponses, ils doivent comprendre les chemins de résolution. L'école du futur \textbf{aura} donc des tablettes et des ordinateurs, mais l'être humain – le prof et l'élève – \textbf{restera} toujours le centre de l'apprentissage. »

\textbf{Commentaire linguistique détaillé (Corrigé commenté) :}
\begin{itemize}
\item \textbf{Futur I (3 occurrences correctes) :}
\begin{itemize}
\item \all{die KI ... \textbf{wird} große Chancen \textbf{bieten}} (Prévision, infinitif \all{bieten} en fin de phrase).
\item \all{Die Lehrer \textbf{werden} mehr Zeit \textbf{haben}} (Prévision, infinitif \all{haben} en fin).
\item \all{Die Schule ... \textbf{wird} Tablets und Computer \textbf{haben} / der Mensch ... \textbf{wird} ... \textbf{bleiben}} (Prévision, infinitifs finaux).
\end{itemize}
\item \textbf{Groupes infinitifs (3 types variés, virgules respectées) :}
\begin{itemize}
\item \all{, \textbf{um} seine Noten \textbf{zu verbessern}} (But : \all{um ... zu}).
\item \all{, \textbf{ohne} kritisch \textbf{zu denken} \textbf{benutzen}} (Manière négative : \all{ohne ... zu} — attention : l'infinitif principal \all{benutzen} est à la fin du groupe principal, \all{zu denken} est dans le groupe \all{ohne}). Structure : \all{nicht [ohne kritisch zu denken] benutzen}.
\item \all{\textbf{Anstatt} nur Antworten \textbf{zu kopieren}} (Alternative : \all{Anstatt ... zu} en tête de phrase, virgule après le groupe).
\end{itemize}
\item \textbf{Lexique thématique réemployé (7 items) :} \all{künstliche Intelligenz, Lernprogramm, Schüler, Apps, Tablets, Computer, Lehrer, Noten, Hausaufgaben} (implicite dans \all{Antworten kopieren}), \all{Zukunft} (dans le titre).
\item \textbf{Grammaire accessoire maîtrisée :}
\begin{itemize}
\item \all{helfen} + Datif : \all{den Schülern ... helfen} (dans la phrase source du texte, réutilisé ici implicitement dans \all{Chancen bieten} + Datif \all{den Schülern}).
\item \all{vertrauen} + Datif : \all{der Technik ... vertrauen}.
\item Verbes modaux + infinitif : \all{kann ... üben}, \all{müssen ... verstehen}, \all{dürfen ... benutzen}.
\item Déclinaison adjectivale : \all{gute Lernprogramm} (accusatif neutre indéfini), \all{schwache Schüler} (accusatif pluriel faible), \all{kritisch denken} (adverbe).
\item Majuscules systématiques à tous les noms (\all{Gymnasium, Einsatz, Unterricht, KI, Chancen, Lernprogramm, Schüler, Noten, Lehrer, Zeit, Apps, Mensch, Mittelpunkt, Lernen}).
\end{itemize}
\item \textbf{Structure argumentative :} Introduction contextuelle $\rightarrow$ Thèse personnelle (\all{Ich bin der Meinung, dass...}) $\rightarrow$ Argument 1 (Avantage élèves + Futur I + \all{um...zu}) $\rightarrow$ Argument 2 (Avantage profs + Futur I + \all{um...zu}) $\rightarrow$ Concession / Nuance (\all{Aber... ohne...zu}) $\rightarrow$ Contre-argument constructif (\all{Anstatt...zu}) $\rightarrow$ Conclusion équilibrée (Futur I + \all{aber}).
\end{itemize}
\end{exoresolu}

\begin{aretenir}[Bilan des acquis de la leçon AL19 — Check-list pour le Bac]
À la fin de cette leçon, tu dois être capable de :
\begin{itemize}
\item \textbf{Comprendre} un texte authentique (annonce, article, dialogue) sur le numérique éducatif et l'IA.
\item \textbf{Maîtriser le vocabulaire :} Noms (article, pluriel), verbes (régime casuel : \all{helfen + Dat}, \all{benutzen + Akk}), adjectifs d'opinion, familles de mots (\all{lernen / der Lerner / die Lernapp}).
\item \textbf{Conjuguer et employer le Futur I :} \all{werden} + Infinitif (fin de phrase, sans \all{zu}). Distinguer prévision, hypothèse, promesse.
\item \textbf{Construire des groupes infinitifs corrects :} \all{um ... zu} (but), \all{ohne ... zu} (sans), \all{anstatt ... zu} (au lieu de). \textbf{Virgule obligatoire. Infinitif + \all{zu} en fin de groupe.}
\item \textbf{Argumenter à l'oral :} Utiliser les formules de prise de parole, réagir, contredire poliment.
\item \textbf{Rédiger un paragraphe d'opinion structuré :} 6--8 lignes, cohérent, grammaticalement correct (Futur I, Infinitivgruppen, cas, majuscules), vocabulaire précis.
\end{itemize}
\end{aretenir}

\newpage

\coursec{Méthodes et exercices résolus}

\courssub{Méthode 1 : Comprendre un texte sur les outils numériques et l'IA}

\begin{methode}[Lire et comprendre un texte sur les outils numériques et l'intelligence artificielle]
\begin{enumerate}
\item \textbf{Lecture globale.} Lis le texte une première fois sans dictionnaire. Repère : le thème (\all{Lernen mit der KI}), le type de texte (article, reportage, interview), la situation (école, famille, entreprise) et la thèse générale (avantages ou dangers du numérique ?).
\item \textbf{Repérage des mots-clés du champ lexical.} Souligne les termes du thème : \all{die künstliche Intelligenz}, \all{das Lernprogramm}, \all{die App}, \all{das Tablet}, \all{der Computer}, \all{lernen}. Ils portent le sens du texte et réapparaissent dans les questions.
\item \textbf{Lecture détaillée.} Pour chaque paragraphe, formule en une phrase française l'idée principale. Entoure les connecteurs (\all{denn}, \all{aber}, \all{deshalb}, \all{außerdem}) : ils indiquent la logique de l'argumentation.
\item \textbf{Devine le sens des mots inconnus} grâce au contexte, aux familles de mots (\all{lernen} $\rightarrow$ \all{das Lernprogramm}, \all{der Lerner}) et aux mots transparents (\all{der Computer}, \all{die App}).
\item \textbf{Réponds aux questions} en reformulant avec tes propres mots, en phrases complètes, sans recopier de longs passages. Cite un court extrait entre guillemets „…“ si la question le demande.
\end{enumerate}
\end{methode}

\begin{exoresolu}[Compréhension de l'écrit : questions sur le texte \all{Lernen mit der KI}]
\textbf{Énoncé.} Texte : \all{In vielen Schulen in Deutschland und auch in Kamerun benutzen Schüler heute Tablets und Lernprogramme. Eine App kann zum Beispiel Fehler sofort korrigieren. Die künstliche Intelligenz wird das Lernen in Zukunft verändern, denn sie passt die Übungen an jeden Schüler an. Aber der Lehrer bleibt wichtig: Er motiviert die Klasse und erklärt schwierige Regeln.}\\
Beantworten Sie die Fragen auf Deutsch :
\begin{enumerate}
\item Was benutzen die Schüler heute in vielen Schulen ?
\item Warum wird die KI das Lernen verändern ?
\item Welche Rolle spielt der Lehrer ?
\end{enumerate}
\tcblower
\textbf{Solution détaillée.}
\begin{enumerate}
\item \textbf{Question 1.} La question porte sur \all{was} (quoi). On repère dans la première phrase : \all{benutzen Schüler heute Tablets und Lernprogramme}. On reformule en phrase complète, verbe en deuxième position :\\
\all{Die Schüler benutzen heute Tablets und Lernprogramme.} « Les élèves utilisent aujourd'hui des tablettes et des programmes d'apprentissage. »
\item \textbf{Question 2.} \all{Warum} appelle une cause : on cherche le connecteur \all{denn} dans le texte. Attention : après \all{denn}, le verbe reste en position 2 (conjonction de coordination). Réponse :\\
\all{Die KI wird das Lernen verändern, weil sie die Übungen an jeden Schüler anpasst.} « L'IA va transformer l'apprentissage parce qu'elle adapte les exercices à chaque élève. »\\
Justification : on peut aussi garder \all{denn} : \all{..., denn sie passt die Übungen an jeden Schüler an.} Avec \all{weil}, le verbe conjugué (\all{anpasst}, verbe à particule \all{anpassen}) va à la fin de la subordonnée.
\item \textbf{Question 3.} On repère \all{der Lehrer bleibt wichtig : er motiviert... und erklärt...}. Réponse rédigée :\\
\all{Der Lehrer bleibt wichtig, denn er motiviert die Klasse und erklärt schwierige Regeln.} « Le professeur reste important car il motive la classe et explique les règles difficiles. »
\end{enumerate}
\textbf{Conclusion méthodologique :} chaque réponse est une phrase complète (sujet + verbe en position 2), qui reformule le texte sans le recopier mot à mot.
\end{exoresolu}

\courssub{Méthode 2 : Employer le vocabulaire du thème}

\begin{methode}[Maîtriser le lexique des médias numériques et de l'IA]
\begin{enumerate}
\item \textbf{Apprends chaque nom avec son article et son pluriel} : \all{der Computer, die Computer} ; \all{die App, die Apps} ; \all{das Tablet, die Tablets} ; \all{das Lernprogramm, die Lernprogramme} ; \all{die künstliche Intelligenz} (sans pluriel) ; \all{die Übung, die Übungen}.
\item \textbf{Mémorise les familles de mots} autour de \all{lernen} : \all{der Lerner} (l'apprenant), \all{das Lernprogramm} (le programme d'apprentissage), \all{die Lernplattform} (la plateforme d'apprentissage), \all{lernen} (apprendre).
\item \textbf{Retiens les verbes associés} : \all{benutzen} (utiliser), \all{korrigieren} (corriger), \all{anpassen} (adapter, particule séparable), \all{verändern} (changer), \all{herunterladen} (télécharger, séparable).
\item \textbf{Utilise le mot juste selon le contexte} : on dit \all{eine App benutzen}, \all{ein Lernprogramm installieren}, \all{mit dem Tablet lernen} (datif après \all{mit}).
\item \textbf{En production écrite}, réemploie au moins trois mots du champ lexical pour montrer ta maîtrise du thème.
\end{enumerate}
\end{methode}

\begin{exoresolu}[Vocabulaire : texte à trous sur le numérique à l'école]
\textbf{Énoncé.} Complétez avec le mot qui convient (article compris si nécessaire) : \all{die künstliche Intelligenz -- das Tablet -- die App -- das Lernprogramm -- lernen}\\
\all{1. Die Schüler in Yaoundé benutzen in der Schule ein \dots\ mit vielen Übungen.\\
2. Mit dieser \dots\ kann man Deutschvokabeln wiederholen.\\
3. \dots\ kann die Fehler der Schüler sofort korrigieren.\\
4. Auf dem \dots\ lesen die Kinder digitale Bücher.\\
5. Mit digitalen Medien \dots\ die Schüler schneller.}
\tcblower
\textbf{Solution détaillée.}
\begin{enumerate}
\item \all{ein Lernprogramm} : l'article indéfini \all{ein} (neutre, accusatif après \all{benutzen}) impose un nom neutre ; « avec beaucoup d'exercices » désigne un programme d'apprentissage.
\item \all{App} : \all{dieser} (féminin datif après \all{mit}) exige un nom féminin ; on répète du vocabulaire avec une application.
\item \all{Die künstliche Intelligenz} : sujet en début de phrase, majuscule ; c'est l'IA qui corrige automatiquement les erreurs.
\item \all{Tablet} : \all{dem} est le datif neutre de \all{das} après la préposition \all{auf} (localisation).
\item \all{lernen} : sujet pluriel \all{die Schüler}, verbe conjugué en position 2, à la 3\up{e} personne du pluriel.
\end{enumerate}
Texte reconstitué : \all{Die Schüler in Yaoundé benutzen in der Schule ein Lernprogramm mit vielen Übungen. Mit dieser App kann man Deutschvokabeln wiederholen. Die künstliche Intelligenz kann die Fehler der Schüler sofort korrigieren. Auf dem Tablet lesen die Kinder digitale Bücher. Mit digitalen Medien lernen die Schüler schneller.}\\
\textbf{Conclusion :} le choix du mot dépend du sens, mais aussi de la grammaire (article, cas, place du verbe).
\end{exoresolu}

\courssub{Méthode 3 : Conjuguer et employer le Futur I}

\begin{methode}[Former et utiliser le Futur I]
\begin{enumerate}
\item \textbf{Formation} : auxiliaire \all{werden} conjugué au présent + infinitif du verbe principal \textbf{à la fin de la phrase}.
\begin{center}
\begin{tabular}{|l|l|l|}
\hline
\rowcolor{popblueL} Personne & werden + Infinitiv & Exemple \\
\hline
ich & werde lernen & \all{Ich werde mit einer App lernen.} \\
\hline
du & wirst lernen & \all{Du wirst ein Tablet benutzen.} \\
\hline
er/sie/es & wird lernen & \all{Die KI wird das Lernen verändern.} \\
\hline
wir & werden lernen & \all{Wir werden online üben.} \\
\hline
ihr & werdet lernen & \all{Ihr werdet viel lesen.} \\
\hline
sie/Sie & werden lernen & \all{Sie werden Fortschritte machen.} \\
\hline
\end{tabular}
\end{center}
\item \textbf{Emplois} : exprimer une action future (\all{Morgen werde ich eine App herunterladen.}), une prévision (\all{Die KI wird die Schulen verändern.}) ou une supposition (\all{Er wird jetzt lernen.} « Il est probablement en train d'apprendre. »).
\item \textbf{Verbe à particule séparable} : la particule reste collée à l'infinitif final : \all{Ich werde die App herunterladen.}
\item \textbf{Attention} : en allemand, le présent avec un complément de temps futur est souvent suffisant (\all{Morgen lerne ich mit dem Tablet.}) ; le Futur I insiste sur la prévision ou l'intention.
\end{enumerate}
\end{methode}

\begin{exoresolu}[Grammaire : conjuguer au Futur I]
\textbf{Énoncé.} Mettez les phrases au Futur I :
\begin{enumerate}
\item \all{Die Schüler benutzen Tablets im Unterricht.}
\item \all{Ich lade eine neue App herunter.}
\item \all{Die künstliche Intelligenz verändert das Lernen.}
\item \all{Wir lernen mit einem Lernprogramm.}
\end{enumerate}
\tcblower
\textbf{Solution détaillée.}
\begin{enumerate}
\item Sujet \all{die Schüler} = 3\up{e} personne du pluriel $\rightarrow$ \all{werden} ; infinitif \all{benutzen} en fin de phrase :\\
\all{Die Schüler werden Tablets im Unterricht benutzen.} « Les élèves utiliseront des tablettes en classe. »
\item Sujet \all{ich} $\rightarrow$ \all{werde} ; verbe à particule séparable \all{herunterladen} : la particule reste soudée à l'infinitif final :\\
\all{Ich werde eine neue App herunterladen.} « Je téléchargerai une nouvelle application. »
\item Sujet singulier \all{die künstliche Intelligenz} $\rightarrow$ \all{wird} :\\
\all{Die künstliche Intelligenz wird das Lernen verändern.} « L'intelligence artificielle transformera l'apprentissage. »
\item Sujet \all{wir} $\rightarrow$ \all{werden} :\\
\all{Wir werden mit einem Lernprogramm lernen.} « Nous apprendrons avec un programme d'apprentissage. »
\end{enumerate}
\textbf{Conclusion :} la structure est toujours : sujet + \all{werden} conjugué (position 2) + éléments + infinitif final. Ne jamais conjuguer le verbe principal.
\end{exoresolu}

\courssub{Méthode 4 : Construire les groupes infinitifs}

\begin{methode}[Former et employer les groupes infinitifs (Infinitivgruppen)]
\begin{enumerate}
\item \textbf{Structure} : \all{um \dots\ zu + Infinitiv} (pour + infinitif, but) ; \all{ohne \dots\ zu + Infinitiv} (sans + infinitif) ; \all{zu + Infinitiv} après certains verbes (\all{hoffen}, \all{versuchen}, \all{beginnen}, \all{vergessen}).
\item \textbf{Place} : le groupe infinitif se met généralement après la proposition principale ; l'infinitif avec \all{zu} est toujours \textbf{à la fin} du groupe : \all{Ich benutze die App, um Deutsch zu lernen.}
\item \textbf{Verbe à particule séparable} : le \all{zu} s'insère \textbf{entre la particule et le verbe} : \all{herunterzuladen}, \all{anzupassen} : \all{Sie benutzt das Tablet, um die App herunterzuladen.}
\item \textbf{Virgule obligatoire} devant \all{um \dots\ zu} et \all{ohne \dots\ zu}.
\item \textbf{Condition d'emploi de \all{um \dots\ zu}} : le sujet des deux actions doit être le même. Sinon, utilise \all{damit} + subordonnée : \all{Der Lehrer erklärt die Regel, damit die Schüler sie verstehen.}
\end{enumerate}
\end{methode}

\begin{exoresolu}[Grammaire : groupes infinitifs avec \all{um \dots\ zu}]
\textbf{Énoncé.} Reliez les phrases avec \all{um \dots\ zu} :
\begin{enumerate}
\item \all{Die Schüler benutzen Tablets. Sie lernen Deutsch.}
\item \all{Ich lade die App herunter. Ich wiederhole die Vokabeln.}
\item \all{Die KI passt die Übungen an. Sie hilft jedem Schüler.}
\end{enumerate}
\tcblower
\textbf{Solution détaillée.}
\begin{enumerate}
\item Les deux sujets sont identiques (\all{die Schüler}) $\rightarrow$ \all{um \dots\ zu} possible. Infinitif \all{lernen} à la fin, virgule devant \all{um} :\\
\all{Die Schüler benutzen Tablets, um Deutsch zu lernen.} « Les élèves utilisent des tablettes pour apprendre l'allemand. »
\item Verbe séparable \all{herunterladen} dans la proposition principale : particule détachée en fin de phrase (\all{lade \dots\ herunter}) ; dans le groupe infinitif, \all{wiederholen} prend \all{zu} :\\
\all{Ich lade die App herunter, um die Vokabeln zu wiederholen.} « Je télécharge l'application pour réviser le vocabulaire. »
\item Verbe séparable \all{anpassen} : particule \all{an} en fin de proposition principale ; infinitif \all{helfen} avec \all{zu} en fin de groupe :\\
\all{Die KI passt die Übungen an, um jedem Schüler zu helfen.} « L'IA adapte les exercices pour aider chaque élève. » Remarque : \all{helfen} exige le datif (\all{jedem Schüler}).
\end{enumerate}
\textbf{Conclusion :} trois réflexes : virgule devant \all{um}, \all{zu} + infinitif à la fin, \all{zu} inséré dans les verbes séparables (\all{herunterzuladen}).
\end{exoresolu}

\courssub{Méthode 5 : Rédiger un court paragraphe d'expression guidée}

\begin{methode}[Produire un court texte sur le numérique et l'école]
\begin{enumerate}
\item \textbf{Analyse la consigne} : nombre de phrases demandé, thème, temps verbaux attendus (ici souvent le Futur I pour les prévisions).
\item \textbf{Plan simple en trois temps} : situation actuelle (présent) $\rightarrow$ avantages ou problèmes $\rightarrow$ prévision au Futur I.
\item \textbf{Phrases d'appui à réutiliser} :
\begin{itemize}
\item \all{Heute benutzen viele Schüler Tablets und Apps.} « Aujourd'hui, beaucoup d'élèves utilisent des tablettes et des applications. »
\item \all{Meiner Meinung nach ist die KI sehr nützlich, denn \dots} « À mon avis, l'IA est très utile car \dots »
\item \all{In Zukunft wird das Lernen anders sein.} « À l'avenir, l'apprentissage sera différent. »
\end{itemize}
\item \textbf{Vérifie} : verbe en position 2, infinitif final au Futur I, majuscules aux noms, articles corrects.
\item \textbf{Conclus} par une phrase d'opinion personnelle.
\end{enumerate}
\end{methode}

\begin{exoresolu}[Expression écrite guidée : le numérique dans mon école]
\textbf{Énoncé.} Schreiben Sie einen kurzen Text (5--6 Sätze) : \all{„Digitale Medien in meiner Schule“}. Benutzen Sie mindestens einmal das Futur I und eine Infinitivgruppe mit \all{um \dots\ zu}.
\tcblower
\textbf{Solution détaillée.}\\
\textbf{Étape 1 — Plan.} 1) situation actuelle ; 2) exemple d'outil ; 3) but (\all{um \dots\ zu}) ; 4) prévision (Futur I) ; 5) opinion.\\
\textbf{Étape 2 — Rédaction.}\\
\all{In meiner Schule in Yaoundé benutzen die Schüler heute Tablets und Computer. Mit einer App wiederholen wir jeden Tag Deutschvokabeln. Ich benutze ein Lernprogramm, um meine Fehler zu korrigieren. In Zukunft wird die künstliche Intelligenz das Lernen in unseren Schulen verändern. Meiner Meinung nach sind digitale Medien sehr nützlich, aber der Lehrer bleibt der wichtigste Helfer.}\\
« Dans mon école à Yaoundé, les élèves utilisent aujourd'hui des tablettes et des ordinateurs. Avec une application, nous révisons chaque jour le vocabulaire d'allemand. J'utilise un programme d'apprentissage pour corriger mes erreurs. À l'avenir, l'intelligence artificielle transformera l'apprentissage dans nos écoles. À mon avis, les médias numériques sont très utiles, mais le professeur reste l'aide la plus importante. »\\
\textbf{Étape 3 — Vérification.} Futur I présent : \all{wird \dots\ verändern} (infinitif final). Groupe infinitif : \all{um meine Fehler zu korrigieren} (virgule + \all{zu} + infinitif final). Verbe en position 2 dans chaque phrase (\all{benutzen}, \all{wiederholen} après le complément initial, \all{bleibt}). Majuscules : \all{Tablets}, \all{Computer}, \all{App}, \all{Deutschvokabeln}, \all{Lernprogramm}, \all{Medien}, \all{Lehrer}, \all{Helfer}.\\
\textbf{Conclusion :} un texte court mais complet, avec les deux structures grammaticales exigées et le vocabulaire du thème.
\end{exoresolu}

\begin{attention}[Les 5 erreurs qui coûtent des points au Bac]
\begin{enumerate}
\item \textbf{Oublier l'infinitif final au Futur I.} Faux : \all{Die KI wird verändern das Lernen.} Correct : \all{Die KI wird das Lernen verändern.} L'infinitif se place toujours en fin de phrase.
\item \textbf{Mal placer le \all{zu} dans les verbes séparables.} Faux : \all{um die App zu herunterladen.} Correct : \all{um die App herunterzuladen.} Le \all{zu} s'insère entre la particule et le verbe.
\item \textbf{Oublier la virgule devant \all{um \dots\ zu}.} Faux : \all{Ich benutze das Tablet um Deutsch zu lernen.} Correct : \all{Ich benutze das Tablet, um Deutsch zu lernen.}
\item \textbf{Oublier la majuscule aux noms et l'article.} Faux : \all{die schüler benutzen computer.} Correct : \all{Die Schüler benutzen Computer.} Tout nom allemand prend une majuscule ; apprends chaque nom avec son article (\all{das Tablet}, \all{die App}, \all{der Computer}).
\item \textbf{Calquer le français.} « utiliser » ne se traduit pas par un faux ami : on dit \all{benutzen} ou \all{verwenden}, jamais un verbe inventé. De même, « passer du temps sur l'ordinateur » se dit \all{Zeit am Computer verbringen}, et non un calque mot à mot.
\end{enumerate}
\end{attention}

\newpage

\coursec{Exercices}

\courssub{Je vérifie mes connaissances}

\begin{exercice}[QCM : grammaire et vocabulaire]
Entoure la bonne réponse.
\begin{enumerate}
\item \all{Die Schüler \ldots\ morgen mit der App lernen.}
\begin{itemize}
\item[a)] \all{werden}
\item[b)] \all{würden}
\item[c)] \all{haben}
\end{itemize}
\item Le pluriel de \all{das Tablet} est :
\begin{itemize}
\item[a)] \all{die Tabletten}
\item[b)] \all{die Tablets}
\item[c)] \all{die Tablet}
\end{itemize}
\item \all{Er geht nach Hause, \ldots\ seine Hausaufgaben \ldots\ machen.}
\begin{itemize}
\item[a)] \all{um \ldots\ zu}
\item[b)] \all{ohne \ldots\ zu}
\item[c)] \all{anstatt \ldots\ zu}
\end{itemize}
\item \all{Die künstliche Intelligenz} signifie :
\begin{itemize}
\item[a)] l'intelligence artificielle
\item[b)] l'intelligence naturelle
\item[c)] l'enseignement numérique
\end{itemize}
\end{enumerate}
\end{exercice}

\begin{exercice}[Vrai ou faux ? Justifie ta réponse]
Réponds par \all{richtig} ou \all{falsch} et justifie en une phrase en français.
\begin{enumerate}
\item Au Futur I, le verbe conjugué est \all{haben} et l'infinitif se place à la fin de la phrase.
\item Dans un groupe infinitif avec \all{um \ldots\ zu}, l'infinitif se place à la fin du groupe.
\item \all{das Lernprogramm} est un nom féminin.
\item \all{ohne \ldots\ zu} exprime l'absence d'une action (« sans faire quelque chose »).
\end{enumerate}
\end{exercice}

\begin{exercice}[Vocabulaire : trouve le mot]
Complète chaque phrase avec le mot du thème qui convient (forme correcte) : \all{die App, der Bildschirm, das Tablet, das Lernprogramm, der Computer}.
\begin{enumerate}
\item Ich schreibe meine Hausaufgaben am \ldots\ .
\item Die Kinder schauen zu lange auf den \ldots\ .
\item Ich benutze diese \ldots\ , um jeden Tag zehn neue Wörter zu lernen.
\item Das \ldots\ hilft mir bei der Grammatik.
\item In der Schule benutzen wir ein \ldots\ statt eines Buches.
\end{enumerate}
\end{exercice}

\begin{exercice}[Conjugaison : mets au Futur I]
Mets les phrases suivantes au Futur I.\\
Exemple : \all{Die Schüler lernen mit Tablets.} $\rightarrow$ \all{Die Schüler werden mit Tablets lernen.}
\begin{enumerate}
\item \all{Die KI hilft den Lehrern.}
\item \all{Wir benutzen ein Lernprogramm.}
\item \all{Ich mache meine Hausaufgaben mit dem Computer.}
\item \all{Die Schule kauft neue Tablets.}
\item \all{Du lernst Deutsch mit einer App.}
\end{enumerate}
\end{exercice}

\courssub{Je m'entraîne}

\begin{exercice}[Thème : traduis en allemand]
Traduis les phrases suivantes en allemand.
\begin{enumerate}
\item Pour réussir l'examen, j'utiliserai une application.
\item Il apprend l'allemand sans faire ses devoirs.
\item L'IA ne remplacera jamais le professeur.
\item Au lieu de regarder la télévision, elle travaille avec un programme d'apprentissage.
\item Nous achèterons demain une tablette pour l'école.
\end{enumerate}
\end{exercice}

\begin{exercice}[Complète avec um \ldots\ zu, ohne \ldots\ zu ou anstatt \ldots\ zu]
Complète les phrases avec la structure qui convient.
\begin{enumerate}
\item Er lernt Vokabeln, \ldots\ die Prüfung \ldots\ bestehen.
\item Sie benutzt die App, \ldots\ ein Buch \ldots\ kaufen.
\item Er verlässt das Zimmer, \ldots\ seine Hausaufgaben \ldots\ machen.
\item Wir gehen in die Bibliothek, \ldots\ im Internet \ldots\ arbeiten.
\item \ldots\ den Lehrer \ldots\ fragen, sucht er die Antwort bei der KI.
\end{enumerate}
\end{exercice}

\begin{exercice}[Texte à trous]
Complète le texte suivant avec les mots de la liste : \all{werden, zu, Tablet, ohne, App, lernen, um, Computer}.\\[4pt]
\all{In Zukunft \ldots\ viele Schüler in Kamerun mit einem \ldots\ lernen. Sie benutzen eine \ldots, \ldots\ schneller Deutsch \ldots\ lernen. Manche Schüler gehen in die Schule, \ldots\ ihre Bücher mitzunehmen. Der \ldots\ bleibt aber wichtig, denn nicht jede Familie hat ein Tablet.}
\end{exercice}

\begin{exercice}[Appariements : mot et traduction]
Relie chaque mot allemand à sa traduction française.
\begin{center}
\begin{tabularx}{\linewidth}{|X|X|}
\hline
\rowcolor{popblueL} \textbf{Deutsch} & \textbf{Français} \\
\hline
1. \all{die künstliche Intelligenz} & a. l'écran \\
\hline
2. \all{der Bildschirm} & b. apprendre \\
\hline
3. \all{das Lernprogramm} & c. l'intelligence artificielle \\
\hline
4. \all{lernen} & d. le devoir (à la maison) \\
\hline
5. \all{die Hausaufgabe} & e. le programme d'apprentissage \\
\hline
6. \all{der Unterricht} & f. le cours \\
\hline
\end{tabularx}
\end{center}
\end{exercice}

\courssub{Je me prépare au Bac}

\begin{exercice}[Compréhension écrite : texte type Bac]
Lis le texte suivant, puis réponds aux questions en allemand, par des phrases complètes.

\begin{textebox}[KI in Österreichs Schulen]
\all{In vielen Schulen in Österreich arbeiten die Schüler schon heute mit der künstlichen Intelligenz. Eine Lehrerin in Wien erklärt: „Die KI wird den Unterricht verändern, aber sie wird den Lehrer nie ersetzen.“ Die Schüler benutzen Lernprogramme, um schneller zu lernen. Mit einer App wiederholen sie jeden Tag ihre Vokabeln, ohne ein Buch zu öffnen. Anstatt lange im Wörterbuch zu suchen, fragen sie den Computer. In Zukunft werden mehr Schulen Tablets kaufen, und die Schüler werden ihre Hausaufgaben digital machen. Aber es gibt auch Probleme: Nicht jede Familie hat einen Computer zu Hause, und manche Kinder schauen zu lange auf den Bildschirm. Die Lehrer sagen: „Die KI ist ein Werkzeug. Der Mensch muss denken, nicht die Maschine.“}
\end{textebox}

\textbf{Questions globales :}
\begin{enumerate}
\item Worum geht es in diesem Text? (2--3 Sätze)
\item Was ist die Meinung der Lehrerin aus Wien?
\end{enumerate}
\textbf{Questions de détail :}
\begin{enumerate}
\setcounter{enumi}{2}
\item Wozu benutzen die Schüler die Lernprogramme?
\item Was machen die Schüler, anstatt im Wörterbuch zu suchen?
\item Nenne zwei Probleme, die der Text erwähnt.
\end{enumerate}
\textbf{Vrai ou faux ?} Justifie ta réponse par une citation du texte.
\begin{enumerate}
\setcounter{enumi}{5}
\item \all{Die KI wird den Lehrer ersetzen.}
\end{enumerate}
\textbf{Questions de langue :}
\begin{enumerate}
\setcounter{enumi}{6}
\item Souligne dans le texte deux phrases au Futur I et recopie-les.
\item Relève dans le texte un groupe infinitif avec \all{um \ldots\ zu} et un avec \all{ohne \ldots\ zu}, puis traduis-les en français.
\end{enumerate}
\end{exercice}

\begin{exercice}[Version : traduis en français]
Traduis le paragraphe suivant en français soigné.

\begin{textebox}[Tablets in der Schweiz]
\all{In einigen Schulen in der Schweiz werden die Schüler bald mit Tablets arbeiten. Sie werden ihre Bücher zu Hause lassen und alles digital lernen. Ein Lehrer sagt: „Die Kinder benutzen die Tablets, um besser zu schreiben und zu rechnen.“ Aber manche Eltern sind skeptisch: Sie glauben, dass die Kinder zu viel spielen, anstatt zu lernen. Die Schule will deshalb die Eltern informieren, ohne die neuen Medien zu verbieten.}
\end{textebox}
\end{exercice}

\begin{exercice}[Production écrite type Bac]
Rédige un texte de 8 à 10 lignes (environ 80 à 100 mots) sur le sujet suivant :\\[4pt]
\textbf{\all{Meiner Meinung nach: Vorteile und Nachteile des Handys in der Schule}}\\[4pt]
Consignes obligatoires :
\begin{itemize}
\item Présente au moins deux avantages et deux inconvénients du téléphone portable à l'école.
\item Emploie au moins deux phrases au Futur I.
\item Emploie au moins un groupe infinitif (\all{um \ldots\ zu}, \all{ohne \ldots\ zu} ou \all{anstatt \ldots\ zu}).
\item Termine par ton opinion personnelle (\all{Meiner Meinung nach \ldots}).
\end{itemize}
\end{exercice}



\newpage

\coursec{Corrigés des exercices}

\courssub{Je vérifie mes connaissances}

\begin{corrige}[Exercice 1 — QCM : grammaire et vocabulaire]
\begin{enumerate}
\item \textbf{a) \all{werden}} : \all{Die Schüler werden morgen mit der App lernen.} Au Futur I, l'auxiliaire conjugué est \all{werden} et l'infinitif (\all{lernen}) se place en fin de phrase. \all{würden} relève du Konjunktiv II (conditionnel), \all{haben} est l'auxiliaire du Perfekt.
\item \textbf{b) \all{die Tablets}} : les noms d'origine anglaise prennent généralement le pluriel en \all{-s} (\all{das Tablet, die Tablets} ; \all{der Computer, die Computer}). Attention : \all{die Tabletten} signifie « les comprimés » (médicaments) — faux ami classique.
\item \textbf{a) \all{um \ldots\ zu}} : \all{Er geht nach Hause, um seine Hausaufgaben zu machen.} exprime le but (« pour faire ses devoirs »). \all{ohne \ldots\ zu} = « sans faire », \all{anstatt \ldots\ zu} = « au lieu de faire » : le sens ne convient pas ici.
\item \textbf{a) l'intelligence artificielle} : \all{künstlich} = artificiel ; \all{die Intelligenz} = l'intelligence (nom féminin en \all{-enz}).
\end{enumerate}
\textbf{Points de vigilance :} ne pas confondre \all{werden} (Futur I) et \all{würden} (Konjunktiv II) ; mémoriser le genre de \all{die Intelligenz} grâce à la terminaison \all{-enz}, toujours féminine (comme \all{die Konferenz}, \all{die Differenz}).
\end{corrige}

\begin{corrige}[Exercice 2 — Vrai ou faux ?]
\begin{enumerate}
\item \textbf{Falsch.} Au Futur I, l'auxiliaire conjugué est \all{werden} et non \all{haben} ; seul l'infinitif du verbe principal se place à la fin. \all{haben} sert d'auxiliaire au Perfekt (\all{Ich habe gelernt}).
\item \textbf{Richtig.} Dans tout groupe infinitif (\all{um \ldots\ zu}, \all{ohne \ldots\ zu}, \all{anstatt \ldots\ zu}), l'infinitif précédé de \all{zu} occupe la position finale du groupe.
\item \textbf{Falsch.} \all{das Lernprogramm} est un nom \textbf{neutre} (article \all{das}), pluriel \all{die Lernprogramme}. Le genre ne se devine pas au sens : il faut l'apprendre avec le mot.
\item \textbf{Richtig.} \all{ohne \ldots\ zu} + infinitif exprime l'absence de l'action : \all{ohne zu zögern} = « sans hésiter ».
\end{enumerate}
\textbf{Points de vigilance :} les francophones oublient souvent la virgule avant le groupe infinitif et placent mal l'infinitif ; en allemand, il est toujours en dernière position du groupe.
\end{corrige}

\begin{corrige}[Exercice 3 — Vocabulaire : trouve le mot]
\begin{enumerate}
\item \all{Ich schreibe meine Hausaufgaben am \textbf{Computer}.} (\all{am} = \all{an dem}, datif après \all{an} pour une position.)
\item \all{Die Kinder schauen zu lange auf den \textbf{Bildschirm}.} (accusatif après \all{auf} dans cette construction.)
\item \all{Ich benutze diese \textbf{App}, um jeden Tag zehn neue Wörter zu lernen.} (\all{benutzen} + accusatif ; \all{die App} est féminin.)
\item \all{Das \textbf{Lernprogramm} hilft mir bei der Grammatik.} (sujet au nominatif, nom neutre.)
\item \all{In der Schule benutzen wir ein \textbf{Tablet} statt eines Buches.} (accusatif neutre : \all{ein Tablet}.)
\end{enumerate}
\textbf{Points de vigilance :} vérifier à chaque fois le cas exigé par le verbe ou la préposition ; \all{helfen} exige le datif (\all{hilft mir}).
\end{corrige}

\begin{corrige}[Exercice 4 — Conjugaison : mets au Futur I]
\begin{enumerate}
\item \all{Die KI wird den Lehrern helfen.} (\all{die KI} = 3\textsuperscript{e} personne du singulier $\rightarrow$ \all{wird} ; le datif \all{den Lehrern} exigé par \all{helfen} est conservé.)
\item \all{Wir werden ein Lernprogramm benutzen.}
\item \all{Ich werde meine Hausaufgaben mit dem Computer machen.}
\item \all{Die Schule wird neue Tablets kaufen.}
\item \all{Du wirst Deutsch mit einer App lernen.} (2\textsuperscript{e} personne : \all{wirst}, avec \all{-st}.)
\end{enumerate}
\textbf{Points de vigilance :} l'infinitif passe obligatoirement en fin de phrase ; ne pas conjuguer le verbe principal (\all{*wird hilft} est impossible). Rappel de la conjugaison de \all{werden} au présent : \all{ich werde, du wirst, er/sie/es wird, wir werden, ihr werdet, sie/Sie werden}.
\end{corrige}

\courssub{Je m'entraîne}

\begin{corrige}[Exercice 5 — Thème : traduis en allemand]
\begin{enumerate}
\item \all{Um die Prüfung zu bestehen, werde ich eine App benutzen.} Le groupe infinitif en tête de phrase occupe la position 1 : le verbe conjugué suit immédiatement (inversion sujet–verbe). \all{die Prüfung bestehen} = « réussir l'examen ».
\item \all{Er lernt Deutsch, ohne seine Hausaufgaben zu machen.} Virgule obligatoire avant \all{ohne \ldots\ zu}.
\item \all{Die KI wird den Lehrer nie ersetzen.} \all{ersetzen} = « remplacer » ; négation du verbe : \all{nie} (jamais).
\item \all{Anstatt fernzusehen, arbeitet sie mit einem Lernprogramm.} Le verbe à particule séparable \all{fernsehen} donne \all{fernzusehen} (\all{zu} s'insère entre la particule et le verbe).
\item \all{Wir werden morgen ein Tablet für die Schule kaufen.}
\end{enumerate}
\textbf{Points de vigilance :} dans les groupes infinitifs, le \all{zu} des verbes à particule se place \emph{entre} la particule et le radical (\all{mitzunehmen}, \all{fernzusehen}) ; ne jamais écrire \all{*zu fernsehen}.
\end{corrige}

\begin{corrige}[Exercice 6 — Complète avec um \ldots\ zu, ohne \ldots\ zu ou anstatt \ldots\ zu]
\begin{enumerate}
\item \all{Er lernt Vokabeln, \textbf{um} die Prüfung \textbf{zu} bestehen.} (but : « pour réussir ».)
\item \all{Sie benutzt die App, \textbf{anstatt} ein Buch \textbf{zu} kaufen.} (alternative : « au lieu d'acheter ».)
\item \all{Er verlässt das Zimmer, \textbf{ohne} seine Hausaufgaben \textbf{zu} machen.} (absence de l'action : « sans faire ».)
\item \all{Wir gehen in die Bibliothek, \textbf{um} im Internet \textbf{zu} arbeiten.} (but.)
\item \all{\textbf{Anstatt} den Lehrer \textbf{zu} fragen, sucht er die Antwort bei der KI.} (groupe en tête $\rightarrow$ inversion : \all{sucht er}.)
\end{enumerate}
\textbf{Points de vigilance :} pour choisir, traduire mentalement : but $\rightarrow$ \all{um \ldots\ zu} ; absence $\rightarrow$ \all{ohne \ldots\ zu} ; substitution $\rightarrow$ \all{anstatt \ldots\ zu}. Le sujet des deux actions doit être le même.
\end{corrige}

\begin{corrige}[Exercice 7 — Texte à trous]
\all{In Zukunft \textbf{werden} viele Schüler in Kamerun mit einem \textbf{Tablet} lernen. Sie benutzen eine \textbf{App}, \textbf{um} schneller Deutsch \textbf{zu} \textbf{lernen}. Manche Schüler gehen in die Schule, \textbf{ohne} ihre Bücher mitzunehmen. Der \textbf{Computer} bleibt aber wichtig, denn nicht jede Familie hat ein Tablet.}

\textbf{Justifications :} \all{werden} : Futur I annoncé par \all{in Zukunft} ; \all{mit einem Tablet} : datif instrumental après \all{mit} ; \all{um \ldots\ zu lernen} : but ; \all{ohne ihre Bücher mitzunehmen} : particule séparable \all{mitnehmen} $\rightarrow$ \all{mitzunehmen} ; \all{der Computer} : sujet au nominatif.
\end{corrige}

\begin{corrige}[Exercice 8 — Appariements : mot et traduction]
1--c (\all{die künstliche Intelligenz} = l'intelligence artificielle) ; 2--a (\all{der Bildschirm} = l'écran) ; 3--e (\all{das Lernprogramm} = le programme d'apprentissage) ; 4--b (\all{lernen} = apprendre) ; 5--d (\all{die Hausaufgabe} = le devoir à la maison) ; 6--f (\all{der Unterricht} = le cours).

\textbf{Points de vigilance :} \all{die Hausaufgabe} s'emploie le plus souvent au pluriel \all{die Hausaufgaben} ; ne pas confondre \all{lernen} (apprendre) et \all{lehren} (enseigner).
\end{corrige}

\courssub{Je me prépare au Bac}

\begin{corrige}[Exercice 9 — Compréhension écrite : texte type Bac]
\textbf{Barème indicatif (sur 10) :} questions globales 2 pts ; questions de détail 3 pts ; vrai/faux justifié 1 pt ; questions de langue 3 pts ; qualité de la langue 1 pt.

\textbf{Questions globales :}
\begin{enumerate}
\item \all{Der Text handelt vom Einsatz der künstlichen Intelligenz und der digitalen Medien (Lernprogramme, Apps, Tablets) in österreichischen Schulen. Er zeigt die Vorteile, aber auch die Probleme dieser Entwicklung.}
\item \all{Die Lehrerin aus Wien meint, dass die KI den Unterricht verändern wird, aber den Lehrer nie ersetzen wird.}
\end{enumerate}
\textbf{Questions de détail :}
\begin{enumerate}
\setcounter{enumi}{2}
\item \all{Die Schüler benutzen die Lernprogramme, um schneller zu lernen.}
\item \all{Anstatt lange im Wörterbuch zu suchen, fragen sie den Computer.}
\item \all{Erstens hat nicht jede Familie einen Computer zu Hause. Zweitens schauen manche Kinder zu lange auf den Bildschirm.}
\end{enumerate}
\textbf{Vrai ou faux :}
\begin{enumerate}
\setcounter{enumi}{5}
\item \textbf{Falsch.} Citation : \all{„Die KI wird den Unterricht verändern, aber sie wird den Lehrer nie ersetzen.“}
\end{enumerate}
\textbf{Questions de langue :}
\begin{enumerate}
\setcounter{enumi}{6}
\item Deux phrases au Futur I (au choix) : \all{Die KI wird den Unterricht verändern.} ; \all{In Zukunft werden mehr Schulen Tablets kaufen.} ; \all{Die Schüler werden ihre Hausaufgaben digital machen.}
\item \all{um schneller zu lernen} = « pour apprendre plus vite » ; \all{ohne ein Buch zu öffnen} = « sans ouvrir un livre ».
\end{enumerate}
\textbf{Points de vigilance :} répondre par des phrases complètes en reprenant les termes de la question ; pour le vrai/faux, la justification par une citation exacte est exigée ; ne pas recopier tout le texte.
\end{corrige}

\begin{corrige}[Exercice 10 — Version : traduis en français]
\textbf{Barème indicatif (sur 5) :} respect du sens 3 pts ; correction et qualité du français 2 pts.

\textbf{Traduction proposée :} « Dans certaines écoles de Suisse, les élèves travailleront bientôt avec des tablettes. Ils laisseront leurs livres à la maison et apprendront tout de manière numérique. Un enseignant déclare : "Les enfants utilisent les tablettes pour mieux écrire et calculer." Mais certains parents sont sceptiques : ils pensent que les enfants jouent trop, au lieu d'apprendre. L'école veut donc informer les parents, sans interdire les nouveaux médias. »

\textbf{Points de vigilance :} rendre le Futur I par le futur français ; \all{rechnen} = « calculer » (et non « raconter » : faux ami) ; \all{anstatt zu lernen} = « au lieu d'apprendre » ; \all{ohne die neuen Medien zu verbieten} = « sans interdire les nouveaux médias ».
\end{corrige}

\begin{corrige}[Exercice 11 — Production écrite type Bac]
\textbf{Barème indicatif (sur 10) :} respect du sujet et des consignes 3 pts ; cohérence et organisation 2 pts ; correction grammaticale (Futur I, groupe infinitif, ordre des mots) 3 pts ; richesse du vocabulaire 2 pts.

\textbf{Proposition de rédaction modèle :}

\all{Das Handy in der Schule hat Vor- und Nachteile. Ein Vorteil ist, dass die Schüler schnell Informationen im Internet finden können. Außerdem können sie Lern-Apps benutzen, um jeden Tag ihre Vokabeln zu wiederholen. Ein Nachteil ist, dass viele Schüler im Unterricht abgelenkt sind und Nachrichten schreiben, anstatt dem Lehrer zuzuhören. Auch Cybermobbing ist ein großes Problem. In Zukunft werden die Schulen klare Regeln für die Handys aufstellen, und die Lehrer werden die Handys am Anfang der Stunde einsammeln. Meiner Meinung nach sollte das Handy nur in der Pause erlaubt sein, damit man sich besser auf den Unterricht konzentrieren kann.}

\textbf{Points de vigilance :} vérifier la présence des deux phrases au Futur I (\all{werden \ldots\ aufstellen}, \all{werden \ldots\ einsammeln}) et du groupe infinitif (\all{um \ldots\ zu wiederholen}, \all{anstatt \ldots\ zuzuhören}) ; respecter la place du verbe en position 2 dans les phrases principales et en fin dans les subordonnées avec \all{dass} et \all{damit} ; terminer impérativement par \all{Meiner Meinung nach} + verbe en position 2.
\end{corrige}

\newpage

\coursec{Fiche bilan}

\begin{popfigure}
\begin{tikzpicture}[
  mindmap,
  concept color=popblue,
  text=white,
  font=\sffamily\bfseries,
  level 1 concept/.append style={level distance=4.5cm, sibling angle=60},
  level 2 concept/.append style={level distance=3cm},
  every node/.style={concept, execute at begin node=\hskip0pt},
  popblue/.style={concept color=popblue},
  poporange/.style={concept color=poporange},
  popgreen/.style={concept color=popgreen},
  poppurple/.style={concept color=poppurple},
  poppink/.style={concept color=poppink},
  popgold/.style={concept color=popgold}
]
\node[concept, minimum size=2.8cm, align=center] (center){Lernen\\mit der KI}
  child[popblue] { node[concept, minimum size=2.2cm, align=center]{Textsupport\\Lernen mit der KI} }
  child[poporange] { node[concept, minimum size=2.2cm, align=center]{Wortschatz\\digitale Medien} }
  child[popgreen] { node[concept, minimum size=2.2cm, align=center]{Grammatik 1\\Futur I} }
  child[poppurple] { node[concept, minimum size=2.2cm, align=center]{Grammatik 2\\Infinitivgruppen} }
  child[poppink] { node[concept, minimum size=2.2cm, align=center]{Kompetenz\\Verstehen \& Produzieren} }
  child[popgold] { node[concept, minimum size=2.2cm, align=center]{Methode\\Commentaire guidé} };
\end{tikzpicture}
\legende{Carte mentale de la leçon AL19 : outils numériques et IA}
\end{popfigure}

\begin{bilan}

\end{bilan}

\courssub{Lexique clé : Digitale Medien und KI}

\begin{center}
\begin{tabularx}{\linewidth}{|>{\bfseries}l|X|X|}
\hline
\rowcolor{popblueL}
\textbf{Deutsch (Artikel, Plural)} & \textbf{Français} & \textbf{Famille de mots / Exemple} \\
\hline
\all{die künstliche Intelligenz (KI), -en} & l'intelligence artificielle (IA) & \all{KI-gestützt} (basé sur l'IA), \all{der KI-Algorithmus} \\
\all{das Lernprogramm, die -e} & le logiciel éducatif & \all{programmieren}, \all{der Programmierer} \\
\all{der Computer, -} & l'ordinateur & \all{computergestützt}, \all{am Computer arbeiten} \\
\all{die App, -s} & l'application (mobile) & \all{herunterladen}, \all{installieren} \\
\all{das Tablet, -s} & la tablette & \all{auf dem Tablet lesen}, \all{Tablet-Klasse} \\
\all{lernen} & apprendre & \all{der Lernende}, \all{das Lernziel}, \all{lebenslanges Lernen} \\
\all{das digitale Medium, -ien} & le média numérique & \all{medial}, \all{Medienkompetenz} \\
\all{der Chatbot, -s} & le chatbot / agent conversationnel & \all{chatten}, \all{automatisiert antworten} \\
\all{personalisiert} & personnalisé & \all{personalisieren}, \all{die Personalisierung} \\
\all{der Datenschutz} & la protection des données & \all{datenschutzkonform}, \all{die DSGVO (RGPD)} \\
\hline
\end{tabularx}
\end{center}

\courssub{Grammaire 1 : Das Futur I (Le futur)}

\begin{center}
\begin{tabularx}{\linewidth}{|>{\bfseries}l|X|X|}
\hline
\rowcolor{poporangeL}
\textbf{Aspekt} & \textbf{Règle / Forme} & \textbf{Beispiel / Exemple} \\
\hline
\textbf{Formation} & \all{werden} (conjugué au présent) + \textbf{infinitif} à la fin de la phrase. & \all{Ich werde ein Tablet kaufen.} (J'achèterai une tablette.) \\
\hline
\textbf{Conjugaison de \textit{werden}} & \all{ich werde, du wirst, er/sie/es wird, wir werden, ihr werdet, sie/Sie werden} & \all{Wir werden KI nutzen.} (Nous utiliserons l'IA.) \\
\hline
\textbf{Emploi principal} & Action future, prédiction, intention, promesse. Souvent avec marqueur temporel (\all{morgen, nächstes Jahr, bald}). & \all{Die KI wird das Lernen verändern.} (L'IA changera l'apprentissage.) \\
\hline
\textbf{Emploi modal (supposition)} & Supposition sur le présent ou le futur (probabilité). Souvent avec \all{wohl, sicher, wahrscheinlich}. & \all{Er wird wohl schon zu Hause sein.} (Il est probablement déjà chez lui.) \\
\hline
\textbf{Ordre des mots (Hauptsatz)} & Position 2 : \all{werden} conjugué \quad | \quad Position finale : \textbf{Infinitif du verbe principal}. & \all{Morgen werde ich die App testen.} \\
\hline
\textbf{Ordre des mots (Nebensatz)} & \all{werden} (conjugué) tout à la fin, \textbf{avant} l'infinitif principal. & \all{..., dass ich die App testen werde.} \\
\hline
\textbf{Négation} & \all{nicht} avant l'infinitif (ou avant le complément). \all{kein} pour l'objet. & \all{Ich werde kein Tablet kaufen.} / \all{Ich werde nicht lernen.} \\
\hline
\end{tabularx}
\end{center}

\courssub{Grammaire 2 : Die Infinitivgruppen (Groupes infinitifs avec \textit{zu})}

\begin{center}
\begin{tabularx}{\linewidth}{|>{\bfseries}l|X|X|}
\hline
\rowcolor{popgreenL}
\textbf{Structure} & \textbf{Règle / Cas d'emploi} & \textbf{Beispiel / Exemple} \\
\hline
\textbf{Formation de base} & \all{zu} + \textbf{infinitif} (à la fin de la proposition). Séparé par une virgule si groupe étendu. & \all{Ich habe vor, \textbf{zu lernen}.} (J'ai l'intention d'apprendre.) \\
\hline
\textbf{Verbes à particule séparable} & \all{zu} s'insère \textbf{entre} la particule et le radical. Écriture en un mot. & \all{anfangen $\to$ anzufangen} \quad \all{Ich beginne, \textbf{anzufangen}.} \\
\hline
\textbf{Verbes en \textit{-ieren} / préfixes inseparables} & Pas de \all{ge-} au participe, pas de séparation. \all{zu} devant l'infinitif. & \all{studieren $\to$ zu studieren} \quad \all{versprechen $\to$ zu versprechen} \\
\hline
\textbf{Infinitif passé (\textit{Perfekt})} & \all{zu} + \textbf{avoir/être à l'infinitif} + \textbf{participe passé}. & \all{Er gibt vor, \textbf{gelernt zu haben}.} (Il prétend avoir appris.) \\
\hline
\textbf{Infinitif passif} & \all{zu} + \textbf{werden} (infinitif) + \textbf{participe passé}. & \all{Das Tablet muss \textbf{geladen werden}.} (La tablette doit être chargée.) \\
\hline
\textbf{Conjonctions prépositionnelles} & \all{um ... zu} (but), \all{ohne ... zu} (moyen négatif), \all{statt ... zu} (alternative), \all{an ... zu} (occupation). & \all{Er nutzt KI, \textbf{um besser zu lernen}.} \\
\hline
\textbf{Antécédent du sujet} & Sujet de la principale = sujet de l'infinitif (contrôle de sujet). & \all{Ich versuche, \textbf{zu verstehen}.} (J'essaie de comprendre \emph{moi-même}.) \\
\hline
\textbf{Objet comme sujet (Passif/Inpersonnel)} & Si l'objet de la principale devient sujet de l'infinitif $\to$ passif ou construction impersonnelle. & \all{Das Buch ist leicht \textbf{zu lesen}.} (Le livre est facile à lire.) \\
\hline
\end{tabularx}
\end{center}

\courssub{Méthode express : Commentaire de texte \& Production}

\begin{enumerate}
\item \textbf{Lire globalement} : Identifier le type de texte (article, blog, rapport), le thème (KI im Unterricht), la thèse principale.
\item \textbf{Lexique actif} : Surligner 10--15 mots-clés (voir liste ci-dessus), noter genre/pluriel, chercher dérivés (\all{lernen $\to$ Lernende, Lernziel}).
\item \textbf{Structure} : Découper en parties logiques (Einführung, Argumente Pro/Contra, Fazit). Noter connecteurs (\all{zwar ... aber, außerdem, deshalb}).
\item \textbf{Grammaire en action} : Repérer \textbf{Futur I} (prédictions de l'auteur) et \textbf{Infinitivgruppen} (buts, intentions, moyens). Transformer : Hauptsatz $\leftrightarrow$ Nebensatz mit \all{dass} / Infinitivgruppe.
\item \textbf{Répondre aux questions} : Phrases complètes en allemand, réemploi du vocabulaire du texte, \textbf{Futur I} pour ses propres prédictions.
\item \textbf{Résumé (Zusammenfassung)} : 5--7 phrases, \textbf{Präteritum} ou \textbf{Perfekt} (style écrit), indirect discourse (\all{Konjunktiv I} si niveau requis, sinon indicatif), pas de « je ».
\item \textbf{Production personnelle (Stellungnahme)} : \all{Meiner Meinung nach ... / Ich finde, dass ... / In Zukunft werde ich ...} + arguments + exemple concret (Cameroun/Allemagne).
\end{enumerate}



\begin{aretenir}[Checklist avant le Bac]
$\square$ Je connais le genre, le pluriel et la traduction des 10 mots-clés du lexique (KI, Lernprogramm, App, Tablet, Datenschutz...).
$\square$ Je sais conjuguer \all{werden} au présent sans hésiter et former le \textbf{Futur I} (ordre des mots Hauptsatz/Nebensatz).
$\square$ Je distingue \textbf{Futur I temporel} (action future) et \textbf{Futur I modal} (supposition avec \all{wohl}).
$\square$ Je construis correctement une \textbf{Infinitivgruppe mit \textit{zu}} (place de \all{zu}, virgule, particule séparable \all{anzufangen}).
$\square$ Je maîtrise les 4 conjonctions prépositionnelles : \all{um ... zu}, \all{ohne ... zu}, \all{statt ... zu}, \all{an ... zu}.
$\square$ Je sais transformer une phrase \all{dass} + indicatif en \textbf{Infinitivgruppe} (et inversement) si le sujet est identique.
$\square$ Je sais former l'\textbf{infinitif passé} (\all{zu haben/sein + Partizip II}) et l'\textbf{infinitif passif} (\all{zu werden + Partizip II}).
$\square$ Je suis capable de rédiger un \textbf{résumé structuré} (5--7 phrases, temps du récit, discours indirect).
$\square$ Je peux exprimer mon avis personnel (\all{Meiner Meinung nach / Ich denke, dass ...}) avec du \textbf{Futur I} pour mes projections.
$\square$ Je relève les \textbf{faux-amis} et pièges : \all{das Handy} $\neq$ handy, \all{das Event} $\neq$ événement (terme marketing), \all{kontrollieren} = vérifier (pas contrôler au sens policier).
$\square$ Je gère la \textbf{majuscule systématique} des noms allemands (\all{die KI, das Tablet, der Datenschutz}) et les guillemets „ “.
$\square$ Je connais au moins 3 arguments \textbf{Pour} et 3 arguments \textbf{Contre} l'IA à l'école (personnalisation vs protection des données, motivation vs dépendance technique).
\end{aretenir}

\findecours

\end{document}
